% Rotational Dynamics — Further Mathematics Upper Sixth — Cours signé OnBuch+
% Official MINESEC syllabus. Compiler avec : tectonic FMA16-rotational-dynamics.tex
\newcommand{\DOCMATIERE}{Further Mathematics}
\newcommand{\DOCNIVEAU}{Upper Sixth}
\newcommand{\DOCMODULE}{Upper Sixth — Further mathematics}
\newcommand{\DOCLECON}{16}
\newcommand{\DOCTITRE}{Rotational Dynamics}
% OnBuch+ — préambule « cours signés » ENRICHI (compilé avec tectonic / XeLaTeX)
% Système visuel « Pop » d'OnBuch+ : fond crème (jamais blanc), bordures encre
% épaisses, ombres « dures » décalées, titres Archivo Black en capitales.
% Version MATHÉMATIQUES : graphes (pgfplots/tikz), boîtes pédagogiques
% (définition, propriété, méthode, attention, le savais-tu, exercice résolu,
% exercice, TP, bilan), page de garde et signature OnBuch+ ancrée en bas.
%
% Macros attendues dans main.tex AVANT \input{preamble} :
%   \DOCMATIERE  \DOCNIVEAU  \DOCMODULE  \DOCLECON (numéro)  \DOCTITRE
\documentclass[11pt]{article}

\usepackage{fontspec}
\usepackage[english]{babel}
\usepackage[a4paper,margin=2cm,top=3.4cm,bottom=2.8cm,headheight=1.4cm,footskip=1.3cm]{geometry}
\usepackage{amsmath,amssymb,amsfonts}
\usepackage{mathtools} % \xmapsto, \coloneqq... souvent utilisés par les modèles
\usepackage{cancel} % simplifications barrées (bilans de forces)
% \abs{z} (module, valeur absolue) et \norm{u} (norme), utilisés par les modèles
\providecommand{\abs}{}\let\abs\relax\DeclarePairedDelimiter{\abs}{\lvert}{\rvert}
\providecommand{\norm}{}\let\norm\relax\DeclarePairedDelimiter{\norm}{\lVert}{\rVert}
\usepackage[table]{xcolor} % [table] : fournit aussi \rowcolors (alternance de lignes)
\usepackage{pagecolor}
\usepackage{tikz}
\usetikzlibrary{arrows.meta,positioning,calc,shapes.geometric,shapes.misc,shapes.arrows,decorations.pathmorphing,decorations.pathreplacing,decorations.markings,patterns,shadows,mindmap,trees,fit,backgrounds,matrix,intersections,angles,quotes}
\usepackage{pgfplots}
\usepackage[version=4]{mhchem} % équations nucléaires \ce{^{235}_{92}U}
\usepackage[european,siunitx]{circuitikz} % schémas électriques
\pgfplotsset{compat=1.18}
\usepgfplotslibrary{fillbetween,groupplots} % aires entre courbes (calcul intégral)
\usepackage{enumitem}
\usepackage{fancyhdr}
% [most] charge toutes les bibliothèques tcolorbox (listings, minted, raster,
% documentation...) et gonfle inutilement la mémoire TeX sur un document long
% (~300 boîtes) : on ne charge que ce qui est réellement utilisé.
\usepackage[skins,breakable]{tcolorbox}
\usepackage{graphicx}
\usepackage{colortbl}
\usepackage{booktabs}
\usepackage{tabularx}
\usepackage{diagbox} % case d'en-tête diagonale des tableaux à double entrée
% Colonnes extensibles centrée / gauche / droite (C, L, R), souvent utilisées par les modèles
\newcolumntype{C}{>{\centering\arraybackslash}X}
\newcolumntype{L}{>{\raggedright\arraybackslash}X}
\newcolumntype{R}{>{\raggedleft\arraybackslash}X}
\usepackage{array}
\usepackage{multicol}
\usepackage{siunitx}
% parse-numbers=false : accepte aussi \SI{2,5 \times 10^{-3}}{...}
\sisetup{locale=UK,output-decimal-marker={.},group-minimum-digits=4,parse-numbers=false}
\DeclareSIUnit\ohm{\ensuremath{\symup{\Omega}}} % U+2126 absent de la police
\DeclareSIUnit\division{div} % réglages d'oscilloscope (ms/div, V/div)
\DeclareSIUnit\tr{tr} % tours (vitesse de rotation, ex. \SI{1500}{\tr\per\minute})
\usepackage{qrcode}
% \degree : commande fréquemment utilisée par les modèles pour un angle ou une
% température en dehors de \SI{}{} (non fournie par les paquets chargés).
\providecommand{\degree}{^{\circ}}
% \qed (fin de démonstration), utilisé par les modèles sans amsthm
\providecommand{\qed}{\hfill$\square$}
% \barre : double barre des valeurs interdites dans un tableau de variations (tkz-tab)
\providecommand{\barre}{\ensuremath{\|}}
% environnement proof (amsthm non chargé) utilisé par les modèles
\ifdefined\proof\else\newenvironment{proof}[1][Proof]{\par\noindent\textit{#1.}\ }{\qed\par}\fi
\usepackage{lastpage}
\usepackage{hyperref}
\hypersetup{hidelinks}

% ============================================================
% Palette « Pop » OnBuch+ — mêmes hex que la classe P (plus_ui.dart)
% ============================================================
\definecolor{poporange}{HTML}{F59321}
\definecolor{popdark}{HTML}{E07A0C}
\definecolor{popcream}{HTML}{FFF6E8}
\definecolor{popink}{HTML}{1C1714}
\definecolor{poppurple}{HTML}{7A5AE0}
\definecolor{popgreen}{HTML}{2E9E6A}
\definecolor{poppink}{HTML}{E0457B}
\definecolor{popblue}{HTML}{2F7DE1}
\definecolor{popgold}{HTML}{C98A12}
\definecolor{popmuted}{HTML}{8A7F76}
\definecolor{popline}{HTML}{EADFD2}
\definecolor{popgray}{HTML}{8A7F76}
\colorlet{popgrey}{popgray}
% Alias tolérants (le modèle nomme parfois les couleurs différemment)
\colorlet{popwhite}{white}
\colorlet{popblack}{popink}
\colorlet{popred}{poppink}
\colorlet{popyellow}{popgold}
% Teintes claires pour fonds de tableaux / zones
\colorlet{popblueL}{popblue!10!white}
\colorlet{poporangeL}{poporange!14!white}
\colorlet{popgreenL}{popgreen!12!white}
\colorlet{poppurpleL}{poppurple!12!white}
\colorlet{poppinkL}{poppink!12!white}
\colorlet{popgoldL}{popgold!14!white}

\pagecolor{popcream}

% ============================================================
% Typographie
% ============================================================
\defaultfontfeatures{Ligatures=TeX}
\setmainfont{PlusJakartaSans-Medium.ttf}[Path=../../../fonts/,BoldFont=PlusJakartaSans-Bold.ttf]
% Maths en sans-serif (Fira Math) avec chiffres et lettres droites de Plus
% Jakarta, pour que les formules se fondent dans le texte.
\usepackage[math-style=ISO,bold-style=ISO]{unicode-math}
\setmathfont{FiraMath-Regular.otf}[Path=../../../fonts/]
\setmathfont{PlusJakartaSans-Medium.ttf}[Path=../../../fonts/,range={up/num,up/{Latin,latin},\mathup}]
% (icomma retiré : décimales avec point en anglais)
% Caractères mathématiques Unicode absents des polices de texte (Plus Jakarta,
% Archivo Black) : rendus par leur équivalent mathématique, en texte comme en
% formule — sinon ils s'affichent en carré vide (ex. « Arithmétique dans ℤ »).
\usepackage{newunicodechar}
\newunicodechar{ℝ}{\ensuremath{\mathbb{R}}}
\newunicodechar{ℤ}{\ensuremath{\mathbb{Z}}}
\newunicodechar{ℂ}{\ensuremath{\mathbb{C}}}
\newunicodechar{ℕ}{\ensuremath{\mathbb{N}}}
\newunicodechar{ℚ}{\ensuremath{\mathbb{Q}}}
\newunicodechar{𝔻}{\ensuremath{\mathbb{D}}}
\newunicodechar{α}{\ensuremath{\alpha}}
\newunicodechar{β}{\ensuremath{\beta}}
\newunicodechar{γ}{\ensuremath{\gamma}}
\newunicodechar{δ}{\ensuremath{\delta}}
\newunicodechar{ε}{\ensuremath{\varepsilon}}
\newunicodechar{θ}{\ensuremath{\theta}}
\newunicodechar{λ}{\ensuremath{\lambda}}
\newunicodechar{μ}{\ensuremath{\mu}}
\newunicodechar{σ}{\ensuremath{\sigma}}
\newunicodechar{τ}{\ensuremath{\tau}}
\newunicodechar{φ}{\ensuremath{\varphi}}
\newunicodechar{ψ}{\ensuremath{\psi}}
\newunicodechar{ω}{\ensuremath{\omega}}
\newunicodechar{Σ}{\ensuremath{\Sigma}}
% majuscules produites par \MakeUppercase dans les titres (α-aminés -> Α)
\newunicodechar{Α}{\ensuremath{\alpha}}
\newunicodechar{Β}{\ensuremath{\beta}}
\newunicodechar{Γ}{\ensuremath{\Gamma}}
\newunicodechar{Δ}{\ensuremath{\Delta}}
\newunicodechar{Ε}{\ensuremath{\varepsilon}}
\newunicodechar{Θ}{\ensuremath{\Theta}}
\newunicodechar{Λ}{\ensuremath{\Lambda}}
\newunicodechar{Μ}{\ensuremath{\mu}}
\newunicodechar{Τ}{\ensuremath{\tau}}
\newunicodechar{Φ}{\ensuremath{\Phi}}
\newunicodechar{Ψ}{\ensuremath{\Psi}}
\newunicodechar{Ω}{\ensuremath{\Omega}}
\newunicodechar{↦}{\ensuremath{\mapsto}}
\newunicodechar{∈}{\ensuremath{\in}}
\newunicodechar{∉}{\ensuremath{\notin}}
\newunicodechar{∧}{\ensuremath{\wedge}}
\newunicodechar{⇔}{\ensuremath{\Leftrightarrow}}
\newunicodechar{⟺}{\ensuremath{\Longleftrightarrow}}
\newunicodechar{⇒}{\ensuremath{\Rightarrow}}
\newunicodechar{⟹}{\ensuremath{\Longrightarrow}}
\newunicodechar{∘}{\ensuremath{\circ}}
\newunicodechar{∪}{\ensuremath{\cup}}
\newunicodechar{∩}{\ensuremath{\cap}}
\newunicodechar{≡}{\ensuremath{\equiv}}
\newunicodechar{⊂}{\ensuremath{\subset}}
\newunicodechar{⊥}{\ensuremath{\perp}}
\newunicodechar{∃}{\ensuremath{\exists}}
\newunicodechar{∀}{\ensuremath{\forall}}
\newunicodechar{∛}{\ensuremath{\sqrt[3]{\,}}}
\newunicodechar{∜}{\ensuremath{\sqrt[4]{\,}}}
\newunicodechar{⃗}{\ensuremath{{}^{\to}}}
\newunicodechar{ⁿ}{\ifmmode^{n}\else\textsuperscript{\ensuremath{n}}\fi}
\newunicodechar{ˣ}{\ifmmode^{x}\else\textsuperscript{\ensuremath{x}}\fi}
\newunicodechar{ᵇ}{\ifmmode^{b}\else\textsuperscript{\ensuremath{b}}\fi}
\newunicodechar{ʳ}{\ifmmode^{r}\else\textsuperscript{\ensuremath{r}}\fi}
\newunicodechar{ᵘ}{\ifmmode^{u}\else\textsuperscript{\ensuremath{u}}\fi}
\newunicodechar{ʸ}{\ifmmode^{y}\else\textsuperscript{\ensuremath{y}}\fi}
\newunicodechar{ᵃ}{\ifmmode^{a}\else\textsuperscript{\ensuremath{a}}\fi}
\newunicodechar{ᵉ}{\ifmmode^{e}\else\textsuperscript{\ensuremath{e}}\fi}
\newunicodechar{ᵏ}{\ifmmode^{k}\else\textsuperscript{\ensuremath{k}}\fi}
\newunicodechar{ᵖ}{\ifmmode^{p}\else\textsuperscript{\ensuremath{p}}\fi}
\newunicodechar{ᴺ}{\ifmmode^{N}\else\textsuperscript{\ensuremath{N}}\fi}
\newunicodechar{ᶜ}{\ifmmode^{c}\else\textsuperscript{\ensuremath{c}}\fi}
\newunicodechar{ʰ}{\ifmmode^{h}\else\textsuperscript{\ensuremath{h}}\fi}
\newunicodechar{ᵠ}{\ifmmode^{q}\else\textsuperscript{\ensuremath{q}}\fi}
\newunicodechar{ₙ}{\ifmmode_{n}\else\textsubscript{\ensuremath{n}}\fi}
\newunicodechar{ₐ}{\ifmmode_{a}\else\textsubscript{\ensuremath{a}}\fi}
\newunicodechar{ᵢ}{\ifmmode_{i}\else\textsubscript{\ensuremath{i}}\fi}
\newunicodechar{ᵣ}{\ifmmode_{r}\else\textsubscript{\ensuremath{r}}\fi}
\newunicodechar{ₖ}{\ifmmode_{k}\else\textsubscript{\ensuremath{k}}\fi}
\newunicodechar{ᵦ}{\ifmmode_{\beta}\else\textsubscript{\ensuremath{\beta}}\fi}
\newunicodechar{ₓ}{\ifmmode_{x}\else\textsubscript{\ensuremath{x}}\fi}
\newfontfamily\popdisplay{ArchivoBlack-Regular.ttf}[Path=../../../fonts/]
\newfontfamily\poplogo{SpaceGrotesk-Bold.ttf}[Path=../../../fonts/]
\newfontfamily\popsemi{PlusJakartaSans-SemiBold.ttf}[Path=../../../fonts/]
\newfontfamily\popxbold{PlusJakartaSans-ExtraBold.ttf}[Path=../../../fonts/]
\newfontfamily\popxboldls{PlusJakartaSans-ExtraBold.ttf}[Path=../../../fonts/,LetterSpace=3.0]

\setlist[enumerate]{leftmargin=1.7em,itemsep=5pt,topsep=4pt}
\setlist[itemize]{leftmargin=1.7em,itemsep=4pt,topsep=4pt,label={\color{poporange}\textbullet}}
\setlength{\parindent}{0pt}
\setlength{\parskip}{5pt}
\renewcommand{\arraystretch}{1.25}

% pgfplots : style Pop par défaut
\pgfplotsset{
  every axis/.append style={axis line style={popink,line width=1pt},tick style={popink},
    grid style={popline,line width=0.5pt},label style={font=\small},tick label style={font=\footnotesize}},
  popcurve/.style={poporange,line width=1.8pt,no markers,smooth},
}
% Même style disponible dans un \draw TikZ simple (hors pgfplots)
\tikzset{popcurve/.style={poporange,line width=1.8pt}}
\tikzset{popgrid/.style={popline,very thin}}
\colorlet{popcurve}{poporange} % le nom du style est parfois utilisé comme couleur

% ============================================================
% Pill / squiggle
% ============================================================
\newcommand{\poppill}[3]{%
  \tikz[baseline=-0.55ex]\node[draw=popink,line width=1.2pt,rounded corners=2.6pt,%
    fill=#2,inner xsep=6.5pt,inner ysep=3pt]{%
    \popxboldls\fontsize{7}{7}\selectfont\color{#3}\MakeUppercase{#1}};%
}
\newcommand{\popsquiggle}[1]{%
  \tikz[baseline=-0.4ex]{
    \draw[#1,line width=2.6pt,line cap=round]
      plot [smooth] coordinates {(0,0.09) (0.34,0.01) (0.68,0.21) (1.02,0.01) (1.36,0.10)};
  }%
}
% Mot-clé surligné (marqueur orange sous le texte)
\newcommand{\cle}[1]{{\popxbold\color{popdark}#1}}

% ============================================================
% En-tête / pied
% ============================================================
\pagestyle{fancy}
\fancyhf{}
\renewcommand{\headrulewidth}{0pt}
\renewcommand{\footrulewidth}{0pt}
\lhead{\raisebox{-0.15ex}{\poplogo\fontsize{13}{13}\selectfont\color{popink}OnBuch\color{poporange}+}}
\rhead{\poppill{\DOCMATIERE\ \textbullet\ \DOCNIVEAU\ \textbullet\ Chapter \DOCLECON}{white}{popink}}
\lfoot{\popxbold\fontsize{7.6}{7.6}\selectfont\color{popmuted}\MakeUppercase{Signed course OnBuch+}\ \textbullet\ {\popsemi\DOCTITRE}}
\rfoot{\tikz[baseline=-0.6ex]\node[rounded corners=8.5pt,fill=popink,minimum height=17pt,minimum width=17pt,inner xsep=5pt,inner sep=0]{\popxbold\fontsize{8}{8}\selectfont\color{popcream}\thepage\,/\,\pageref*{LastPage}};}

% ============================================================
% Titres
% ============================================================
\newcommand{\chaptitre}[2]{%
  \begin{tcolorbox}[enhanced,colback=poporange,colframe=popink,boxrule=2.6pt,arc=11pt,
      drop shadow=popink,left=15pt,right=15pt,top=12pt,bottom=11pt,before skip=3pt,after skip=18pt]
    \poppill{#2}{popink}{popcream}\par\vspace{9pt}
    {\raggedright\hyphenpenalty=10000\exhyphenpenalty=10000\popdisplay\fontsize{22}{24}\selectfont\color{popcream}\MakeUppercase{#1}\par}
  \end{tcolorbox}
}
% Section numérotée : pastille orange avec le numéro + titre Archivo
\newcounter{popsec}
\newcommand{\coursec}[1]{%
  \stepcounter{popsec}\setcounter{popsub}{0}%
  \par\vspace{16pt}\noindent
  \begin{minipage}{\linewidth}
  \tikz[baseline=-0.7ex]\node[draw=popink,line width=1.6pt,fill=poporange,rounded corners=4pt,minimum size=22pt,inner sep=0]{\popdisplay\fontsize{12}{12}\selectfont\color{popcream}\thepopsec};%
  \hspace{8pt}{\popdisplay\fontsize{15}{17}\selectfont\color{popink}\MakeUppercase{#1}}\par
  \vspace{4pt}\noindent{\color{poporange}\rule{36pt}{3.6pt}}
  \end{minipage}\par\nopagebreak\vspace{8pt}
}
\newcounter{popsub}
\newcommand{\courssub}[1]{%
  \stepcounter{popsub}%
  \par\vspace{9pt}\noindent{\popxbold\fontsize{11.5}{13}\selectfont\color{popink}{\color{poporange}\thepopsec.\thepopsub}\hspace{6pt}#1}\par\nopagebreak\vspace{4pt}
}

% ============================================================
% Boîtes pédagogiques — même recette : fond blanc, bordure épaisse colorée,
% ombre dure encre, badge plein. Titre optionnel [#1].
% ============================================================
\tcbset{popbase/.style={breakable,enhanced,colback=white,boxrule=2.3pt,arc=9pt,
  shadow={3pt}{-3pt}{0pt}{popink},left=14pt,right=14pt,top=11pt,bottom=11pt,before skip=11pt,after skip=15pt,
  fontupper=\popsemi\fontsize{10.6}{14.5}\selectfont\color{popink},
  fontlower=\popsemi\fontsize{10.6}{14.5}\selectfont\color{popink}}}
% Coche dessinée (le glyphe ✓ n'existe pas dans les polices du document)
\newcommand{\popcheck}[1]{\tikz[baseline=-0.5ex]\draw[#1,line width=1.6pt,line cap=round,line join=round](-0.13,0.0)--(-0.04,-0.09)--(0.13,0.1);}
\providecommand{\coche}{\popcheck{popgreen}}
\providecommand{\resultat}[1]{\par\smallskip\noindent\fcolorbox{popgreen}{popgreenL}{\parbox{\dimexpr\linewidth-2\fboxsep-2\fboxrule\relax}{\textbf{Result.} #1}}\par\smallskip}
\newcommand{\popboxhead}[3]{\poppill{#1}{#2}{white}\ifx\relax#3\relax\else\hspace{6pt}{\popxbold\fontsize{10.6}{12}\selectfont\color{popink}#3}\fi\par\vspace{7pt}}

\newtcolorbox{aretenir}[1][]{popbase,colframe=popgreen,before upper={\popboxhead{Key points}{popgreen}{#1}}}
\newtcolorbox{exemplebox}[1][]{popbase,colframe=poporange,before upper={\popboxhead{Exemple}{poporange}{#1}}}
\newtcolorbox{definition}[1][]{popbase,colframe=popblue,before upper={\popboxhead{Definition}{popblue}{#1}}}
\newtcolorbox{propriete}[1][]{popbase,colframe=poppurple,before upper={\popboxhead{Property}{poppurple}{#1}}}
\newtcolorbox{methode}[1][]{popbase,colframe=popgold,before upper={\popboxhead{Method}{popgold}{#1}}}
\newtcolorbox{attention}[1][]{popbase,colframe=poppink,before upper={\popboxhead{Warning}{poppink}{#1}}}
\newtcolorbox{experience}[1][]{popbase,colframe=popblue,colback=popblueL,before upper={\popboxhead{Experiment}{popblue}{#1}}}
% « Le savais-tu ? » : carte violette pleine (contexte, culture, Cameroun)
\newtcolorbox{savaistu}[1][]{popbase,colback=poppurple,colframe=popink,
  fontupper=\popsemi\fontsize{10.4}{14.2}\selectfont\color{white},
  before upper={\poppill{Did you know?}{white}{poppurple}\ifx\relax#1\relax\else\hspace{6pt}{\popxbold\color{white}#1}\fi\par\vspace{7pt}}}
% Exercice résolu : énoncé en haut, \tcblower puis la solution sur fond crème
\newcounter{popexo}\newcounter{popexr}
\newtcolorbox{exoresolu}[1][]{popbase,colframe=poporange,colbacklower=popcream,
  segmentation style={popink,line width=1.2pt,dashed},
  before upper={\stepcounter{popexr}\popboxhead{Worked example \thepopexr}{poporange}{#1}},
  before lower={\poppill{Solution}{popink}{popcream}\par\vspace{6pt}}}
% Exercice (non résolu en place) — la correction est dans la partie Corrigés
\newtcolorbox{exercice}[1][]{popbase,colframe=popink,
  before upper={\stepcounter{popexo}\popboxhead{Exercise \thepopexo}{popink}{#1}}}
% Corrigé
\newtcolorbox{corrige}[1][]{popbase,colframe=popgreen,colback=popgreenL,
  before upper={\popboxhead{Solution}{popgreen}{#1}}}
% Objectifs / prérequis / bilan
\newtcolorbox{objectifs}{popbase,colframe=popink,colback=poporangeL,
  before upper={\popboxhead{Chapter objectives}{popink}{}}}
\newtcolorbox{prerequis}{popbase,colframe=popink,colback=popblueL,
  before upper={\popboxhead{Prerequisites}{popblue}{}}}
\newtcolorbox{bilan}{popbase,colframe=popink,colback=white,boxrule=2.8pt,
  before upper={\popboxhead{Fiche bilan}{poporange}{L'essentiel à connaître pour l'examen}}}
% Figure encadrée avec légende
\newtcolorbox{popfigure}[1][]{enhanced,breakable=false,colback=white,colframe=popline,boxrule=1.4pt,arc=8pt,
  left=10pt,right=10pt,top=10pt,bottom=8pt,before skip=10pt,after skip=12pt,halign=center,
  lower separated=false,halign lower=center,
  fontlower=\popsemi\fontsize{9}{11.5}\selectfont\color{popmuted},
  #1}
\newcommand{\legende}[1]{\par\vspace{6pt}{\popsemi\fontsize{9}{11.5}\selectfont\color{popmuted}#1}}

% ============================================================
% Signature OnBuch+
% ============================================================
\newcommand{\poplogoinline}[1]{{\poplogo\fontsize{#1}{#1}\selectfont\color{popink}OnBuch\color{poporange}+}}
\newcommand{\signatureonbuch}{%
  \begin{tcolorbox}[enhanced,colback=popink,colframe=popink,boxrule=2.2pt,arc=10pt,
      drop shadow=poporange,left=14pt,right=14pt,top=10pt,bottom=10pt,before skip=0pt,after skip=0pt]
    \tikz[baseline=-0.7ex]\node[circle,fill=popgreen,draw=popcream,line width=1.3pt,minimum size=22pt,inner sep=0]{\popcheck{white}};%
    \hspace{9pt}\begin{minipage}[c]{0.62\linewidth}
      {\popxbold\fontsize{12}{13}\selectfont\color{popcream}Signed\ \textbullet\ {\poplogo OnBuch\color{poporange}+}}\par\vspace{2pt}
      {\raggedright\popsemi\fontsize{8}{10.5}\selectfont\color{popline}\MakeUppercase{OnBuch+ teaching team}\\\MakeUppercase{Follows the official MINESEC syllabus}\par}
    \end{minipage}\hfill
    {\poplogo\fontsize{22}{22}\selectfont\color{popcream}OnBuch\color{poporange}+}
  \end{tcolorbox}%
}

% ============================================================
% Page de garde
% #1 titre · #2 matière · #3 niveau · #4 module · #5 n° leçon · #6 durée ·
% #7 description / objectifs (texte ou itemize)
% ============================================================
\newcommand{\pagegarde}[7]{%
  \thispagestyle{empty}
  \newgeometry{a4paper,margin=1.7cm,top=1.6cm,bottom=1.4cm}%
  \begin{tikzpicture}[remember picture,overlay]
    \fill[poporange!18] ([xshift=-3.2cm,yshift=-3.6cm]current page.north east) circle (4.6cm);
    \fill[poppurple!14] ([xshift=2.4cm,yshift=7.6cm]current page.south west) circle (3.8cm);
  \end{tikzpicture}%
  {\poplogo\fontsize{30}{30}\selectfont\color{popink}OnBuch\color{poporange}+}\hfill
  \raisebox{0.6ex}{\poppill{Signed course \textbullet\ Premium}{popink}{popcream}}\par
  \vspace{1pt}\popsquiggle{poppurple}\par
  \vspace{8pt}
  \begin{tcolorbox}[enhanced,colback=poporange,colframe=popink,boxrule=2.8pt,arc=13pt,
      drop shadow=popink,left=18pt,right=18pt,top=12pt,bottom=13pt,after skip=10pt]
    \poppill{#2 \textbullet\ #3}{popink}{popcream}\hspace{5pt}\poppill{#4}{white}{popink}\par\vspace{8pt}
    {\popdisplay\fontsize{14}{14}\selectfont\color{popink}CHAPTER #5}\par\vspace{4pt}
    {\raggedright\hyphenpenalty=10000\exhyphenpenalty=10000\popdisplay\fontsize{25}{27}\selectfont\color{popcream}\MakeUppercase{#1}\par}
    \vspace{8pt}
    {\popxbold\fontsize{9.5}{9.5}\selectfont\color{popink}\MakeUppercase{Suggested time: #6}}
  \end{tcolorbox}
  \begin{tcolorbox}[enhanced,colback=white,colframe=popink,boxrule=2.2pt,arc=10pt,
      drop shadow=popink,left=16pt,right=16pt,top=10pt,bottom=10pt,after skip=8pt]
    \poppill{In this chapter}{poppurple}{white}\par\vspace{5pt}
    \popsemi\fontsize{9.3}{12.6}\selectfont\color{popink}#7
  \end{tcolorbox}
  \noindent\begin{minipage}[t]{0.487\linewidth}
    \begin{tcolorbox}[enhanced,colback=popgreen,colframe=popink,boxrule=2.2pt,arc=10pt,
        drop shadow=popink,left=11pt,right=11pt,top=10pt,bottom=10pt,height=3.1cm,valign=center]
      \tcbox[colback=white,colframe=popink,boxrule=1.3pt,arc=5pt,boxsep=0pt,nobeforeafter,
        left=3pt,right=3pt,top=3pt,bottom=3pt,valign=center]{\qrcode[height=1.9cm]{https://play.google.com/store/apps/details?id=com.luvvix.onbuch&hl=en}}%
      \hspace{8pt}\begin{minipage}[c]{0.5\linewidth}
        {\popdisplay\fontsize{11}{12.5}\selectfont\color{white}\MakeUppercase{Download OnBuch}}\par\vspace{4pt}
        {\raggedright\popsemi\fontsize{8.4}{11}\selectfont\color{white}Free \textbullet\ Google Play\\AI tutor, past papers, results\par}
      \end{minipage}
    \end{tcolorbox}
  \end{minipage}\hfill
  \begin{minipage}[t]{0.487\linewidth}
    \begin{tcolorbox}[enhanced,colback=poppurple,colframe=popink,boxrule=2.2pt,arc=10pt,
        drop shadow=popink,left=11pt,right=11pt,top=10pt,bottom=10pt,height=3.1cm,valign=center]
      \tikz[baseline=-0.7ex]\node[circle,fill=white,draw=popink,line width=1.3pt,minimum size=1.6cm,inner sep=0]{\popdisplay\fontsize{24}{24}\selectfont\color{poppurple}+};%
      \hspace{8pt}\begin{minipage}[c]{0.6\linewidth}
        {\popdisplay\fontsize{11}{12.5}\selectfont\color{white}\MakeUppercase{Go OnBuch+}}\par\vspace{4pt}
        {\raggedright\popsemi\fontsize{8.4}{11}\selectfont\color{white}All signed courses, unlimited AI tutor, PDF sheets — OnBuch+ tab in the app\par}
      \end{minipage}
    \end{tcolorbox}
  \end{minipage}
  \vspace{8pt}
  \popcontenu
  \vfill
  \signatureonbuch
  \restoregeometry
  \newpage
}

% Rangée « Ce cours contient » (page de garde)
\newcommand{\poptile}[3]{% #1 couleur #2 gros texte #3 légende
  \begin{tcolorbox}[enhanced,colback=#1,colframe=popink,boxrule=2pt,arc=9pt,shadow={2.5pt}{-2.5pt}{0pt}{popink},
      left=6pt,right=6pt,top=9pt,bottom=9pt,halign=center,valign=center,height=1.75cm,width=\linewidth,nobeforeafter]
    {\popdisplay\fontsize{12.5}{13}\selectfont\color{white}\MakeUppercase{#2}}\par\vspace{3pt}
    {\centering\popsemi\fontsize{7.8}{9.5}\selectfont\color{white}#3\par}
  \end{tcolorbox}}
\newcommand{\popcontenu}{%
  \par\noindent{\popxboldls\fontsize{7.5}{7.5}\selectfont\color{popmuted}THIS SIGNED COURSE CONTAINS}\par\vspace{7pt}
  \noindent\begin{minipage}[t]{0.235\linewidth}\poptile{popblue}{Course}{complete, illustrated, step by step}\end{minipage}\hfill
  \begin{minipage}[t]{0.235\linewidth}\poptile{poporange}{Methods}{and worked examples}\end{minipage}\hfill
  \begin{minipage}[t]{0.235\linewidth}\poptile{poppink}{Exercises}{exam style + solutions}\end{minipage}\hfill
  \begin{minipage}[t]{0.235\linewidth}\poptile{popgreen}{Summary sheet}{the essentials to remember}\end{minipage}\par}

% Sommaire
\newtcolorbox{sommaire}{enhanced,colback=white,colframe=popink,boxrule=2.2pt,arc=10pt,
  drop shadow=popink,left=18pt,right=18pt,top=13pt,bottom=13pt,before skip=6pt,after skip=20pt,
  before upper={\poppill{Contents}{poppurple}{white}\par\vspace{9pt}\popsemi\fontsize{10.6}{14.5}\selectfont\color{popink}}}

% Signature de fin de cours — poussée tout en bas de la dernière page
\newcommand{\findecours}{%
  \par\vfill\nopagebreak
  \begin{center}\popsquiggle{poporange}\end{center}\vspace{4pt}
  \signatureonbuch
}
\AtBeginDocument{\def\llbracket{\mathopen{[\mkern-3.5mu[}}\def\rrbracket{\mathclose{]\mkern-3.5mu]}}} % intervalles d'entiers (glyphe ⟦ absent de la police)
% Glyphes absents de Fira Math : redéfinis à partir de glyphes disponibles
\AtBeginDocument{%
  \def\setminus{\mathbin{\backslash}}\let\smallsetminus\setminus
  \def\vdots{\mathord{\vbox{\baselineskip3.2pt\lineskiplimit0pt\kern4pt\hbox{.}\hbox{.}\hbox{.}}}}%
  \def\ddots{\mathinner{\mkern1mu\raise6.5pt\hbox{.}\mkern2mu\raise3.6pt\hbox{.}\mkern2mu\raise0.7pt\hbox{.}\mkern1mu}}%
  \def\triangle{\mathord{\scalebox{0.9}{$\symup{\Delta}$}}}%
  \def\nabla{\mathord{\raisebox{\depth}{\scalebox{1}[-1]{$\symup{\Delta}$}}}}%
  \def\checkmark{\popcheck{popdark}}%
  \def\star{\mathbin{\ast}}\def\bigstar{\mathord{\ast}}%
  \def\diamond{\mathbin{\scalebox{0.6}{$\Diamond$}}}%
  \def\bigcup{\mathop{\mathchoice{\raisebox{-0.3em}{\scalebox{1.7}{$\cup$}}}{\scalebox{1.25}{$\cup$}}{\cup}{\cup}}\displaylimits}%
  \def\bigcap{\mathop{\mathchoice{\raisebox{-0.3em}{\scalebox{1.7}{$\cap$}}}{\scalebox{1.25}{$\cap$}}{\cap}{\cap}}\displaylimits}%
  \def\lesseqqgtr{\mathrel{\lessgtr}}\def\bigoplus{\mathop{\oplus}\displaylimits}%
}
\providecommand{\ssi}{if and only if} % abréviation fréquente
\AtBeginDocument{\providecommand{\wideparen}{\overparen}} % arc de cercle

% Fonctions usuelles que les modèles utilisent sans que amsmath les fournisse
\providecommand{\arsinh}{\operatorname{arsinh}}\providecommand{\arcosh}{\operatorname{arcosh}}
\providecommand{\artanh}{\operatorname{artanh}}\providecommand{\arsech}{\operatorname{arsech}}
\providecommand{\arcsch}{\operatorname{arcsch}}\providecommand{\arcoth}{\operatorname{arcoth}}
\providecommand{\arcsinh}{\operatorname{arsinh}}\providecommand{\arccosh}{\operatorname{arcosh}}
\providecommand{\arctanh}{\operatorname{artanh}}\providecommand{\sech}{\operatorname{sech}}
\providecommand{\csch}{\operatorname{csch}}\providecommand{\arcsec}{\operatorname{arcsec}}
\providecommand{\arccsc}{\operatorname{arccsc}}\providecommand{\arccot}{\operatorname{arccot}}
\providecommand{\sgn}{\operatorname{sgn}}\providecommand{\lcm}{\operatorname{lcm}}
\providecommand{\Var}{\operatorname{Var}}\providecommand{\Cov}{\operatorname{Cov}}

%% FIGFIX-BEGIN v4 — correctifs globaux des figures (docs/onbuch-generation/tools/figfix)
\usepackage{adjustbox}\usepackage{etoolbox}
% toute figure TikZ/pgfplots plus large que la ligne est réduite à la largeur disponible
\BeforeBeginEnvironment{tikzpicture}{\begin{adjustbox}{max width=\linewidth}}
\AfterEndEnvironment{tikzpicture}{\end{adjustbox}}
% légendes pgfplots lisibles par-dessus les courbes
\pgfplotsset{every axis/.append style={legend style={fill=white,fill opacity=0.92,text opacity=1,draw=black!15,inner sep=3pt}}}
% étiquettes d'arêtes (syntaxe edge["..."]) avec fond blanc
\tikzset{every edge quotes/.append style={fill=white,inner sep=1.2pt}}
%% FIGFIX-END
\begin{document}

\pagegarde{Rotational Dynamics}{Further Mathematics}{Upper Sixth}{Upper Sixth — Further mathematics}{16}{10 periods}{This chapter extends Newtonian mechanics to rigid bodies rotating about a fixed axis: moment of inertia, radius of gyration, torque-driven motion, angular momentum, rotational energy, the compound pendulum and impulsive problems. It provides the full theoretical and problem-solving toolkit required by the MINESEC Further Mathematics syllabus and the GCE Advanced Level examination.
\vspace{4pt}
\begin{multicols}{2}\small
\begin{itemize}
\item By the end of this chapter I can define the moment of inertia of a particle and of a system of particles about an axis and compute it for discrete systems.
\item By the end of this chapter I can state and use the radius of gyration and the parallel and perpendicular axes theorems.
\item By the end of this chapter I can derive by integration the moments of inertia of standard bodies (rod, ring, disc, sphere, cylinder, rectangular lamina).
\item By the end of this chapter I can write and solve the equation of motion of a rigid body under a constant torque, including connected-particle systems with a pulley.
\item By the end of this chapter I can compute the moment of momentum (angular momentum) of a rotating body and apply its conservation.
\item By the end of this chapter I can calculate the kinetic energy of a rotating rigid body and apply conservation of energy to rotational problems.
\end{itemize}\end{multicols}}

\begin{sommaire}
\begin{multicols}{2}
\begin{enumerate}
\item Moment of Inertia and Radius of Gyration
\item Standard Moments of Inertia by Integration
\item Equation of Motion under a Constant Torque
\item Moment of Momentum and Its Conservation; Impulse
\item Energy of a Rotating Body and the Force on the Axis
\item The Compound Pendulum and Synthesis
\item Integration activity
\item Methods and worked examples
\item Exercises
\item Solutions to the exercises
\item Summary sheet
\end{enumerate}
\end{multicols}
\end{sommaire}

\begin{objectifs}
\begin{itemize}
\item By the end of this chapter I can define the moment of inertia of a particle and of a system of particles about an axis and compute it for discrete systems.
\item By the end of this chapter I can state and use the radius of gyration and the parallel and perpendicular axes theorems.
\item By the end of this chapter I can derive by integration the moments of inertia of standard bodies (rod, ring, disc, sphere, cylinder, rectangular lamina).
\item By the end of this chapter I can write and solve the equation of motion of a rigid body under a constant torque, including connected-particle systems with a pulley.
\item By the end of this chapter I can compute the moment of momentum (angular momentum) of a rotating body and apply its conservation.
\item By the end of this chapter I can calculate the kinetic energy of a rotating rigid body and apply conservation of energy to rotational problems.
\item By the end of this chapter I can analyse the compound pendulum, find its period and the length of the equivalent simple pendulum.
\item By the end of this chapter I can determine the force exerted on the axis of rotation using the motion of the centre of mass.
\item By the end of this chapter I can solve impulsive problems involving rotating bodies and connected particles.
\end{itemize}
\end{objectifs}

\begin{prerequis}
\begin{itemize}
\item Newton's laws of motion and the dynamics of connected particles (Lower Sixth Mechanics)
\item Circular motion: angular velocity ω, relation v = rω, centripetal acceleration (Lower Sixth)
\item The simple pendulum and simple harmonic motion, T = 2π√(l/g)
\item Definite integration and standard integrals, including use of symmetry (Pure Mathematics)
\item Work, kinetic energy ½mv² and the principle of conservation of mechanical energy
\end{itemize}
\end{prerequis}

\newpage

\coursec{Moment of Inertia and Radius of Gyration}

At the Mokolo market in Yaound\'e, a carpenter spins a grindstone with his foot while a nearby potter shapes clay on a rotating wheel: both know instinctively that a heavy rim keeps the wheel turning long after the push stops. Why does a bicycle wheel resist being tilted once it spins fast, and why does a fufu pounder's pestle swing like a pendulum when hung from its end? These everyday Cameroonian observations hide a deep truth: rotating bodies obey laws as precise as Newton's, but with mass replaced by a new quantity --- the \cle{moment of inertia}. This chapter builds the complete mechanics of rotation needed for the GCE Advanced Level, and it begins with the central question of this section: \emph{how do we measure a body's resistance to being spun?}

\courssub{The Rotational Analogue of Mass}

\textbf{From Newton's second law to rotation.}\par
Consider a single particle $P$ of mass $m$, attached to a fixed vertical axis $\Delta$ by a light rigid rod of length $r$, free to rotate in a horizontal plane. Suppose a force $\vec{F}$ is applied to the particle, tangential to its circular path. Newton's second law gives the tangential acceleration:
\[
F = ma.
\]
The angular acceleration $\ddot{\theta}$ of the particle about $\Delta$ is related to the tangential acceleration by $a = r\ddot{\theta}$. Multiplying Newton's law by $r$:
\[
Fr = m r^2\,\ddot{\theta}.
\]
The left-hand side, $Fr$, is the \emph{moment of the force} about the axis --- the \cle{torque}, usually written $C$ or $\Gamma$. The equation therefore reads
\[
C = \left(mr^2\right)\ddot{\theta}.
\]
Compare this with $F = m\ddot{x}$ for linear motion. The roles correspond exactly:

\begin{center}
\begin{tabularx}{\linewidth}{|l|X|X|}
\hline
\rowcolor{popblueL} & \textbf{Linear motion} & \textbf{Rotation about a fixed axis} \\
\hline
Cause & Force $F$ & Torque $C = Fr$ \\
\hline
Effect & Acceleration $\ddot{x}$ & Angular acceleration $\ddot{\theta}$ \\
\hline
Inertia & Mass $m$ & $mr^2$ \\
\hline
\end{tabularx}
\end{center}

The quantity $mr^2$ plays for rotation the role that mass plays for translation. Crucially, it depends not only on \emph{how much} mass there is, but on \emph{where} that mass sits relative to the axis. Doubling the distance $r$ quadruples the resistance to angular acceleration. This is exactly why the potter's wheel has its mass concentrated in a heavy rim: a large $r$ for most of the mass gives a large $mr^2$, so the wheel's angular velocity barely changes between kicks of the foot.

\begin{definition}[Moment of inertia of a particle and of a system]
The \cle{moment of inertia} of a particle of mass $m$ about an axis $\Delta$, at perpendicular distance $r$ from $\Delta$, is
\[
I = mr^2.
\]
For a system of particles of masses $m_1, m_2, \dots, m_n$ at perpendicular distances $r_1, r_2, \dots, r_n$ from the \emph{same} axis $\Delta$, the moment of inertia is
\[
I = \sum_{i=1}^{n} m_i r_i^2.
\]
The SI unit is the kilogram metre squared, $\mathrm{kg\,m^2}$. The distance $r_i$ is always the \emph{perpendicular} distance from the particle to the axis.
\end{definition}

\begin{attention}[Distance to the axis, not to a point]
The most frequent error at this stage is to measure $r_i$ from a \emph{point} (for example the origin) instead of measuring the \emph{perpendicular distance to the axis}. If the axis is the vertical line $x = a$ in a horizontal plane, a particle at horizontal coordinate $x_i$ has $r_i = |x_i - a|$, whatever its other coordinates may be.
\end{attention}

\courssub{Computing $I$ for Discrete Systems}

\begin{methode}[Checklist for computing $I$ of a discrete system]
\begin{enumerate}
\item \textbf{Identify the axis} of rotation precisely (a line, not a point).
\item \textbf{Measure the perpendicular distance} $r_i$ of each particle from that axis; ignore masses on the axis itself (they contribute $0$) and ignore light rods or frames (their mass is neglected).
\item \textbf{Compute each term} $m_i r_i^2$, taking care to square the distance.
\item \textbf{Sum} the terms: $I = \sum m_i r_i^2$. State the unit, $\mathrm{kg\,m^2}$.
\end{enumerate}
\end{methode}

\begin{popfigure}
\begin{tikzpicture}[scale=1.05, >=stealth]
% axis
\draw[line width=1.4pt, popblue] (0,-0.6) -- (0,3.4) node[above, popblue]{\small axis $\Delta$};
% frame
\draw[line width=1pt, popdark] (-2.6,1.2) -- (3.2,1.2);
\draw[line width=1pt, popdark] (0,1.2) -- (1.8,2.6);
\draw[line width=1pt, popdark] (0,1.2) -- (-1.6,2.4);
% particles
\fill[poporange] (-2.6,1.2) circle (0.14) node[below left=1pt, popdark]{\small $m_1$};
\fill[poporange] (3.2,1.2) circle (0.14) node[below right=1pt, popdark]{\small $m_2$};
\fill[poporange] (1.8,2.6) circle (0.14) node[right=1pt, popdark]{\small $m_3$};
% distances
\draw[<->, poppurple] (-2.6,0.55) -- (0,0.55) node[midway, below, poppurple]{\small $r_1$};
\draw[<->, poppurple] (0,0.85) -- (3.2,0.85) node[midway, below, poppurple]{\small $r_2$};
\draw[dashed, popmuted] (1.8,2.6) -- (0,2.6);
\draw[<->, poppurple] (-0.35,2.6) -- (1.45,2.6) node[midway, above, poppurple]{\small $r_3$};
\end{tikzpicture}
\legende{Three particles on a light frame rotating about a vertical axis $\Delta$. Each $r_i$ is the perpendicular distance from the particle to the axis, and $I = m_1 r_1^2 + m_2 r_2^2 + m_3 r_3^2$.}
\end{popfigure}

\begin{exoresolu}[A dumb-bell about two different axes]
A dumb-bell consists of two particles, each of mass $2\ \mathrm{kg}$, joined by a light rod of length $1.2\ \mathrm{m}$. Find the moment of inertia (a) about an axis perpendicular to the rod through its midpoint $G$; (b) about a parallel axis through one of the masses.
\tcblower
\textbf{Solution.}
(a) Each mass is at distance $r = 0.6\ \mathrm{m}$ from the axis through $G$:
\[
I_G = 2\times(0.6)^2 + 2\times(0.6)^2 = 0.72 + 0.72 = 1.44\ \mathrm{kg\,m^2}.
\]
(b) About the axis through one mass: that mass contributes $2\times 0^2 = 0$; the other is at distance $1.2\ \mathrm{m}$:
\[
I = 2\times(1.2)^2 = 2.88\ \mathrm{kg\,m^2}.
\]
Note that $I = 2 I_G$: the moment of inertia depends essentially on the choice of axis.
\end{exoresolu}

\begin{exoresolu}[Particles at the vertices of a square frame]
Four particles of masses $m$, $2m$, $3m$ and $4m$ are fixed at the vertices $A$, $B$, $C$, $D$ of a light square frame of side $a$. Find the moment of inertia about the side $AD$.
\tcblower
\textbf{Solution.} The axis $AD$ passes through $A$ and $D$, so the masses at $A$ and $D$ contribute nothing. The particles at $B$ (mass $2m$) and $C$ (mass $3m$) are each at perpendicular distance $a$ from $AD$:
\[
I_{AD} = 2m\,a^2 + 3m\,a^2 = 5ma^2.
\]
The same calculation about the diagonal $AC$ would give $I_{AC} = 2m\left(\tfrac{a}{\sqrt{2}}\right)^2 + 4m\left(\tfrac{a}{\sqrt{2}}\right)^2 = 3ma^2$, since $B$ and $D$ are at distance $a/\sqrt{2}$ from the diagonal.
\end{exoresolu}

\courssub{Radius of Gyration}

The total mass of a body is $M = \sum m_i$, but the moment of inertia $\sum m_i r_i^2$ involves the \emph{distribution} of that mass. It is convenient to compress all the information into a single length.

\begin{definition}[Radius of gyration]
The \cle{radius of gyration} $k$ of a body about an axis is the distance defined by
\[
I = Mk^2, \qquad\text{i.e.}\qquad k = \sqrt{\frac{I}{M}},
\]
where $M$ is the total mass. Physically, $k$ is the distance from the axis at which the \emph{whole mass} of the body could be concentrated into a single particle without changing the moment of inertia.
\end{definition}

For the dumb-bell above, $M = 4\ \mathrm{kg}$ and $I_G = 1.44\ \mathrm{kg\,m^2}$, so
\[
k = \sqrt{\frac{1.44}{4}} = 0.6\ \mathrm{m},
\]
which is exactly the distance of each mass from $G$ --- as it must be, since all the mass already lies at that distance. For a body whose mass is spread at different distances, $k$ is a root-mean-square distance: $k^2 = \dfrac{\sum m_i r_i^2}{\sum m_i}$. The radius of gyration is the quantity engineers quote because it separates the two factors $M$ and $k^2$: two flywheels of equal mass but different $k$ behave very differently.

\courssub{The Theorem of Parallel Axes}

Computing $I$ about every conceivable axis would be tedious. Two theorems allow us to \emph{transfer} a known result from one axis to another. The first and most important concerns parallel axes.

\begin{propriete}[Theorem of parallel axes]
Let $G$ be the centre of mass of a body of mass $M$, and let $I_G$ be its moment of inertia about an axis through $G$. The moment of inertia about a \emph{parallel} axis $\Delta$ at distance $h$ from the first is
\[
I = I_G + Mh^2.
\]
(Proof examinable --- see below.)
\end{propriete}

\begin{popfigure}
\begin{tikzpicture}[scale=1.0, >=stealth]
% body
\draw[line width=1pt, popdark, fill=popblueL] plot[smooth cycle, tension=0.9] coordinates {(-2.2,0.3) (-1.4,1.7) (0.2,2.1) (1.8,1.5) (2.4,0.2) (1.2,-0.9) (-0.8,-1.0)};
% G and axes
\fill[popdark] (0,0.5) circle (0.07) node[above right=0pt, popdark]{\small $G$};
\draw[line width=1.2pt, popblue] (0,-1.5) -- (0,2.7) node[above, popblue]{\small axis through $G$};
\draw[line width=1.2pt, poporange] (1.6,-1.5) -- (1.6,2.7) node[above, poporange]{\small axis $\Delta$};
\draw[<->, poppurple] (0,-1.3) -- (1.6,-1.3) node[midway, below, poppurple]{\small $h$};
% element P
\fill[popgreen] (-1.1,1.2) circle (0.07) node[left=0pt, popdark]{\small $P(m)$};
\draw[dashed, popmuted] (-1.1,1.2) -- (0,1.2);
\draw[dashed, popmuted] (-1.1,1.2) -- (1.6,1.2);
\node[popmuted, align=center] at (-0.55,1.42) {\small $x$};
\node[popmuted, align=center] at (0.85,1.42) {\small $h-x$};
\end{tikzpicture}
\legende{Parallel axes: the axis through the centre of mass $G$ and a parallel axis $\Delta$ at distance $h$. A general element $P$ of mass $m$ has coordinate $x$ measured from the axis through $G$.}
\end{popfigure}

\textbf{Proof.}\par
Choose the axis through $G$ as the line $x = 0$ in a plane perpendicular to both axes, and let $\Delta$ be the parallel line $x = h$. A general particle $P$ of mass $m_i$ has perpendicular distance $x_i$ from the axis through $G$; its perpendicular distance from $\Delta$ is $(x_i - h)$. (Distances perpendicular to the axes in the other direction are unchanged, so we may work with signed coordinates $x_i$ in the plane containing both axes.) Then
\begin{align*}
I &= \sum m_i (x_i - h)^2 = \sum m_i x_i^2 - 2h\sum m_i x_i + h^2 \sum m_i \\
&= I_G - 2h\sum m_i x_i + M h^2.
\end{align*}
But $\sum m_i x_i = M\bar{x}$, where $\bar{x}$ is the coordinate of the centre of mass measured from the axis through $G$; since that axis passes \emph{through} $G$, $\bar{x} = 0$ and the middle term vanishes. Hence
\[
I = I_G + Mh^2. \qquad \blacksquare
\]

\begin{attention}[One axis must pass through $G$]
The theorem is valid \emph{only} when one of the two axes passes through the centre of mass. It is \textbf{false} that $I_A = I_B + Mh^2$ for two arbitrary parallel axes $A$ and $B$. If you know $I_A$ about an axis not through $G$ and want $I_B$, you must go via $G$: first $I_G = I_A - Mh_1^2$, then $I_B = I_G + Mh_2^2$. Note also the corollary: among all parallel axes, the moment of inertia is \emph{least} about the axis through $G$.
\end{attention}

\begin{exemplebox}[Checking the dumb-bell]
For the dumb-bell, $I_G = 1.44\ \mathrm{kg\,m^2}$, $M = 4\ \mathrm{kg}$, and the axis through one mass is at $h = 0.6\ \mathrm{m}$ from $G$:
\[
I = I_G + Mh^2 = 1.44 + 4(0.6)^2 = 1.44 + 1.44 = 2.88\ \mathrm{kg\,m^2},
\]
agreeing with the direct computation.
\end{exemplebox}

\courssub{The Theorem of Perpendicular Axes}

The second transfer theorem applies only to \emph{laminae} --- plane bodies of negligible thickness.

\begin{propriete}[Theorem of perpendicular axes]
Let a lamina lie in the $xy$-plane, and let $Ox$, $Oy$, $Oz$ be three mutually perpendicular axes, with $Oz$ perpendicular to the lamina. Then
\[
I_z = I_x + I_y,
\]
where $I_x$, $I_y$, $I_z$ are the moments of inertia of the lamina about the three axes. (Proof examinable --- see below.)
\end{propriete}

\begin{popfigure}
\begin{tikzpicture}[scale=1.0, >=stealth]
% lamina (xy-plane seen in perspective)
\draw[line width=1pt, popdark, fill=popgreenL] plot[smooth cycle, tension=0.9] coordinates {(-2.0,-0.4) (-0.9,0.9) (0.8,1.1) (2.1,0.3) (1.4,-1.0) (-0.6,-1.2)};
% axes
\draw[->, line width=1pt, popblue] (0,0) -- (2.9,0) node[right, popblue]{\small $x$};
\draw[->, line width=1pt, popblue] (0,0) -- (-1.5,-1.5) node[below left, popblue]{\small $y$};
\draw[->, line width=1pt, poporange] (0,0) -- (0,2.3) node[above, poporange]{\small $z$};
% element P
\fill[popdark] (1.1,0.55) circle (0.06);
\node[popdark] at (1.35,0.75) {\small $P(m)$};
\draw[dashed, popmuted] (1.1,0.55) -- (1.1,0);
\draw[dashed, popmuted] (1.1,0.55) -- (0,0.55);
\node[popmuted] at (1.25,0.28) {\small $y$};
\node[popmuted] at (0.55,0.72) {\small $x$};
\fill[popdark] (0,0) circle (0.06) node[below right=0pt]{\small $O$};
\end{tikzpicture}
\legende{A lamina in the $xy$-plane with a general element $P$ of mass $m$ at $(x,y)$. Its distance from $Oz$ is $\sqrt{x^2+y^2}$, from $Ox$ is $|y|$, from $Oy$ is $|x|$.}
\end{popfigure}

\textbf{Proof.}\par
Let a general element of the lamina have mass $m_i$ and coordinates $(x_i, y_i)$. Its perpendicular distances from the three axes are:
\[
\text{to } Ox:\ |y_i|, \qquad \text{to } Oy:\ |x_i|, \qquad \text{to } Oz:\ \sqrt{x_i^2 + y_i^2}.
\]
Therefore
\[
I_z = \sum m_i (x_i^2 + y_i^2) = \sum m_i x_i^2 + \sum m_i y_i^2 = I_y + I_x. \qquad \blacksquare
\]

\begin{attention}[Laminae only]
The perpendicular axes theorem requires the body to be \emph{plane}: for a three-dimensional solid, the distance of an element from $Oz$ involves $z_i$ as well, and the relation fails. Do not apply it to spheres, cones or cylinders.
\end{attention}

\courssub{Exam Technique: Transferring Standard Results}

The real power of the two theorems appears when they are combined with the standard results proved by integration in the next section (Lesson 110). Anticipating those results, the method is always the same.

\begin{methode}[Transferring a standard moment of inertia to a new axis]
\begin{enumerate}
\item \textbf{Recall or derive} $I$ about the standard axis (usually through $G$).
\item \textbf{Parallel axes?} If the new axis is parallel to the standard one, add $Mh^2$ where $h$ is the distance between the axes --- checking first that the standard axis passes through $G$.
\item \textbf{Perpendicular axes?} For a lamina, if you know two of $I_x, I_y, I_z$, deduce the third from $I_z = I_x + I_y$. Symmetry often gives $I_x = I_y$, so $I_z = 2I_x$.
\item \textbf{Chain the theorems} if necessary, always passing through the centre of mass.
\end{enumerate}
\end{methode}

\begin{exoresolu}[Rod about an end; square lamina about a side and a corner]
(a) A uniform rod $AB$ of mass $M$ and length $2a$ has $I_G = \tfrac{1}{3}Ma^2$ about the perpendicular axis through its centre. Find $I$ about the parallel axis through $A$.
(b) A uniform square lamina of mass $M$ and side $2a$ has $I = \tfrac{1}{3}Ma^2$ about each of the axes through its centre $G$ parallel to the sides. Find its moment of inertia (i) about the perpendicular axis through $G$; (ii) about one side; (iii) about the perpendicular axis through a corner.
\tcblower
\textbf{Solution.}
(a) $A$ is at distance $a$ from $G$, so by parallel axes:
\[
I_A = \tfrac{1}{3}Ma^2 + Ma^2 = \tfrac{4}{3}Ma^2.
\]
(b)(i) By perpendicular axes, with $I_x = I_y = \tfrac{1}{3}Ma^2$:
\[
I_z = I_x + I_y = \tfrac{2}{3}Ma^2.
\]
(ii) A side is parallel to a central axis, at distance $a$ from $G$:
\[
I_{\text{side}} = \tfrac{1}{3}Ma^2 + Ma^2 = \tfrac{4}{3}Ma^2.
\]
(iii) A corner is at distance $a\sqrt{2}$ from $G$ along the perpendicular direction measured in the plane; the perpendicular axis through the corner is parallel to $Oz$, so by parallel axes:
\[
I_{\text{corner}} = I_z + M(a\sqrt{2})^2 = \tfrac{2}{3}Ma^2 + 2Ma^2 = \tfrac{8}{3}Ma^2.
\]
\end{exoresolu}

\begin{aretenir}[Key points of this section]
\begin{itemize}
\item Moment of inertia of a particle: $I = mr^2$; of a system: $I = \sum m_i r_i^2$, in $\mathrm{kg\,m^2}$, where $r_i$ is the \emph{perpendicular} distance to the \emph{axis}.
\item Radius of gyration: $I = Mk^2$; $k$ is the distance at which the whole mass could be concentrated.
\item Parallel axes: $I = I_G + Mh^2$ --- valid only when one axis passes through the centre of mass $G$; $I$ is minimal about the axis through $G$.
\item Perpendicular axes (laminae only): $I_z = I_x + I_y$.
\item To transfer between two axes neither of which passes through $G$, always go via $G$.
\end{itemize}
\end{aretenir}

\begin{exercice}[Practice]
\begin{enumerate}
\item Three particles of masses $1\ \mathrm{kg}$, $2\ \mathrm{kg}$ and $3\ \mathrm{kg}$ are at the vertices of an equilateral triangle of side $0.6\ \mathrm{m}$, the $1\ \mathrm{kg}$ mass being at the top vertex. Find $I$ about the base of the triangle, and the corresponding radius of gyration.
\item A body has $I = 5\ \mathrm{kg\,m^2}$ about an axis through its centre of mass and $M = 2\ \mathrm{kg}$. Find $I$ about a parallel axis $1.5\ \mathrm{m}$ away.
\item A body of mass $4\ \mathrm{kg}$ has $I = 10\ \mathrm{kg\,m^2}$ about an axis $A$, and $G$ is $1\ \mathrm{m}$ from $A$. Find $I$ about a parallel axis $B$ on the other side of $G$, $0.5\ \mathrm{m}$ from $G$.
\item A uniform circular lamina has $I = \tfrac{1}{2}Ma^2$ about the perpendicular axis through its centre. Using symmetry and the perpendicular axes theorem, find $I$ about a diameter, and then about a tangent line in the plane of the lamina.
\end{enumerate}
\end{exercice}

\begin{corrige}[Exercise 1]
\textbf{1.} The base passes through the $2\ \mathrm{kg}$ and $3\ \mathrm{kg}$ masses (zero contribution); the $1\ \mathrm{kg}$ mass is at the height of the triangle, $h = 0.6\times\frac{\sqrt{3}}{2} = 0.3\sqrt{3}\ \mathrm{m}$. So $I = 1\times(0.3\sqrt{3})^2 = 0.27\ \mathrm{kg\,m^2}$, and $k = \sqrt{0.27/6} \approx 0.212\ \mathrm{m}$.
\textbf{2.} $I = 5 + 2(1.5)^2 = 9.5\ \mathrm{kg\,m^2}$.
\textbf{3.} Via $G$: $I_G = 10 - 4(1)^2 = 6\ \mathrm{kg\,m^2}$; then $I_B = 6 + 4(0.5)^2 = 7\ \mathrm{kg\,m^2}$.
\textbf{4.} By symmetry $I_x = I_y$ for two perpendicular diameters, and $I_z = I_x + I_y = \tfrac{1}{2}Ma^2$ gives $I_{\text{diameter}} = \tfrac{1}{4}Ma^2$. A tangent in the plane is parallel to a diameter at distance $a$: $I_{\text{tangent}} = \tfrac{1}{4}Ma^2 + Ma^2 = \tfrac{5}{4}Ma^2$.
\end{corrige}

\begin{savaistu}[Flywheels in Cameroonian workshops]
The grindstones and potters' wheels of Mokolo market are flywheels: their rims are deliberately heavy so that $k$ --- and hence $I = Mk^2$ --- is large. A large moment of inertia smooths out the irregular pushes of the foot, storing rotational kinetic energy $\tfrac{1}{2}I\omega^2$ (Section~5) and releasing it steadily. The same principle steadies the engines of the \emph{bendskins} (motorcycle taxis) of Douala and the generators that keep Yaound\'e lit during load-shedding.
\end{savaistu}

\coursec{Standard Moments of Inertia by Integration}

In the previous section we defined the moment of inertia of a system of particles as $I=\sum m_i r_i^2$ and introduced the radius of gyration $k$ through $I=Mk^2$. That discrete formula is perfect for a few point masses, but a real rigid body --- a rod, a disc, a flywheel, a sphere --- contains an effectively infinite number of mass elements. To compute its moment of inertia we must replace the sum by an integral. This section develops that technique, derives the standard results quoted throughout the GCE Advanced Level examination, and shows how to combine them for composite bodies.

\courssub{From Sums to Integrals}

Imagine a lamina divided into $n$ tiny pieces, the $i$-th piece having mass $\delta m_i$ at perpendicular distance $r_i$ from the axis. Then
\[ I=\sum_{i=1}^{n} r_i^2\,\delta m_i. \]
As the pieces are made arbitrarily small, this sum becomes an integral taken over the whole body:
\begin{definition}[Moment of inertia of a continuous body]
\[ I=\int r^2\,dm, \]
where $r$ is the perpendicular distance of the mass element $dm$ from the axis of rotation, and the integral extends over the whole body.
\end{definition}

The whole art lies in choosing a \cle{mass element} $dm$ whose points are all at (approximately) the same distance $r$ from the axis, and expressing $dm$ through the density:
\begin{itemize}
\item \textbf{Linear density} $\lambda$ (mass per unit length, $\mathrm{kg\,m^{-1}}$): $dm=\lambda\,dx$ for rods.
\item \textbf{Surface density} $\sigma$ (mass per unit area, $\mathrm{kg\,m^{-2}}$): $dm=\sigma\,dA$ for laminas, rings, discs.
\item \textbf{Volume density} $\rho$ (mass per unit volume, $\mathrm{kg\,m^{-3}}$): $dm=\rho\,dV$ for solid cylinders and spheres.
\end{itemize}
For a \emph{uniform} body the density is constant and equals total mass divided by total length, area or volume.

\begin{methode}[General strategy for finding $I$ by integration]
\begin{enumerate}
\item Draw the body and mark the axis clearly.
\item Choose a thin element (strip, ring, disc) all of whose points lie at the same distance $r$ from the axis.
\item Write the mass of the element: $dm=\lambda\,dx$, $\sigma\,dA$ or $\rho\,dV$.
\item Write $dI=r^2\,dm$ and integrate between the geometric limits of the body.
\item Replace the density by $M$ over total length/area/volume to express $I$ in terms of $M$.
\item (Optional) Read off the radius of gyration from $I=Mk^2$.
\end{enumerate}
\end{methode}

\courssub{The Uniform Rod}

Consider a uniform rod $AB$ of length $2a$ and mass $M$, so its linear density is $\lambda=\dfrac{M}{2a}$.

\textbf{Axis perpendicular to the rod through its centre.} Place the origin at the centre $G$ and the axis perpendicular to the rod. An element of length $dx$ at distance $x$ from $G$ has mass $dm=\lambda\,dx$ and every point of it is at distance $x$ from the axis.

\begin{popfigure}
\begin{center}
\begin{tikzpicture}[scale=1.0]
\draw[line width=3pt,popblue] (-3,0) -- (3,0);
\draw[poporange,line width=4pt] (1.2,0) -- (1.7,0);
\draw[->,popdark,dashed] (0,-1.3) -- (0,1.3) node[above]{\small axis};
\fill[popdark] (0,0) circle(1.5pt) node[below left=1pt]{\small $G$};
\draw[<->,popmuted] (0,-0.45) -- (1.45,-0.45) node[midway,below]{\small $x$};
\draw[poporange] (1.45,0.18) node[above]{\small $dx$};
\draw[<->,popmuted] (-3,0.5) -- (3,0.5) node[midway,above]{\small $2a$};
\node[below] at (-3,-0.1) {\small $-a$};
\node[below] at (3,-0.1) {\small $a$};
\end{tikzpicture}
\legende{Rod with element of length $dx$ at distance $x$ from the central axis.}
\end{center}
\end{popfigure}

Hence
\[ I_G=\int_{-a}^{a} x^2\,\lambda\,dx=\lambda\left[\frac{x^3}{3}\right]_{-a}^{a}=\frac{2\lambda a^3}{3}=\frac{2a^3}{3}\cdot\frac{M}{2a}=\boxed{\tfrac{1}{3}Ma^2}. \]

\textbf{Axis perpendicular to the rod through one end.} Measuring $x$ from the end $A$, the limits become $0$ and $2a$:
\[ I_A=\int_{0}^{2a} x^2\,\lambda\,dx=\lambda\frac{(2a)^3}{3}=\frac{8\lambda a^3}{3}=\frac{4}{3}Ma^2. \]
Note the consistency check with the parallel axes theorem: $I_A=I_G+M a^2=\tfrac13 Ma^2+Ma^2=\tfrac43 Ma^2$. 

\begin{attention}[The axis position changes everything]
For a rod of length $2a$, $I=\tfrac13 Ma^2$ about the \emph{centre} but $\tfrac43 Ma^2$ about an \emph{end} --- four times larger. Many candidates lose marks by quoting $\tfrac13 Ma^2$ when the rod swings about its end. Always identify the axis first, and check whether the question uses length $2a$ or length $l$ (with $l=2a$, the results are $\tfrac{1}{12}Ml^2$ and $\tfrac13 Ml^2$).
\end{attention}

\courssub{The Uniform Ring and the Thin Hollow Cylinder}

For a uniform ring (hoop) of mass $M$ and radius $a$ about the axis through its centre perpendicular to its plane, \emph{every} mass element is at the same distance $a$ from the axis. No real integration is needed:
\[ I=\int a^2\,dm=a^2\int dm=Ma^2. \]
The same argument applies to a thin hollow cylinder (a ``hoop'' extended along the axis) about its central axis: $I=Ma^2$. This is why a hoop rolls downhill more slowly than a disc of the same mass and radius --- more of its mass sits far from the axis.

\courssub{The Uniform Circular Disc}

Let a uniform disc have mass $M$ and radius $a$, with surface density $\sigma=\dfrac{M}{\pi a^2}$, and take the axis through its centre perpendicular to its plane. Points of the disc lie at all distances from $0$ to $a$, so we slice the disc into \cle{concentric rings}: a ring of radius $r$ and width $dr$ has area $2\pi r\,dr$, and every point of it is at distance $r$ from the axis.

\begin{popfigure}
\begin{center}
\begin{tikzpicture}[scale=1.1]
\draw[fill=popblueL,draw=popblue] (0,0) circle(2);
\draw[fill=white,draw=poporange,line width=1pt] (0,0) circle(1.25);
\draw[draw=poporange,line width=1pt] (0,0) circle(1.05);
\draw[->,popdark,dashed] (0,0) -- (0,2.5) node[above]{\small axis};
\draw[->,popmuted] (0,0) -- (1.15,0.42) node[midway,above left=-2pt]{\small $r$};
\draw[<->,popmuted] (1.05,-0.55) -- (1.25,-0.62) node[midway,below right=-3pt]{\small $dr$};
\draw[->,popmuted] (0,0) -- (1.6,-1.2) node[midway,below left=-2pt]{\small $a$};
\fill[popdark] (0,0) circle(1.5pt);
\end{tikzpicture}
\legende{Disc divided into concentric rings of radius $r$ and width $dr$.}
\end{center}
\end{popfigure}

The ring's mass is $dm=\sigma\cdot 2\pi r\,dr$, so
\begin{align*}
I&=\int_0^a r^2\,\sigma\,2\pi r\,dr=2\pi\sigma\int_0^a r^3\,dr=2\pi\sigma\,\frac{a^4}{4}=\frac{\pi\sigma a^4}{2}.
\end{align*}
Substituting $\sigma=\dfrac{M}{\pi a^2}$:
\[ \boxed{I=\tfrac12 Ma^2}. \]
A solid cylinder of mass $M$ and radius $a$ about its own axis is just a stack of such discs, so the same formula holds: $I=\tfrac12 Ma^2$ (independent of the height).

\textbf{Solid cylinder about a diameter through its centre.} Slice the cylinder of height $h$ into discs of thickness $dx$ perpendicular to the axis, at distance $x$ from the centre. Each disc has mass $dm=\dfrac{M}{h}dx$ and, by the perpendicular axes theorem applied to a disc, moment of inertia $\tfrac14 dm\,a^2$ about a diameter \emph{of that disc}; the parallel axes theorem shifts this to the diameter through $G$:
\[ dI=\tfrac14 a^2\,dm+x^2\,dm. \]
Integrating from $-h/2$ to $h/2$:
\[ I=\frac{M}{h}\int_{-h/2}^{h/2}\left(\frac{a^2}{4}+x^2\right)dx=M\left(\frac{a^2}{4}+\frac{h^2}{12}\right). \]

\courssub{The Uniform Solid Sphere about a Diameter}

Let the sphere have mass $M$, radius $a$, and volume density $\rho=\dfrac{M}{\frac43\pi a^3}$. Take the diameter as the $x$-axis and slice the sphere perpendicular to it into thin discs. The disc at position $x$ has radius $y=\sqrt{a^2-x^2}$ (Pythagoras on a cross-section) and thickness $dx$.

\begin{popfigure}
\begin{center}
\begin{tikzpicture}[scale=1.1]
\draw[fill=popblueL,draw=popblue] (0,0) circle(2);
\draw[popblue,dashed] (-2,0) arc(180:360:2 and 0.55);
\draw[popblue] (2,0) arc(0:180:2 and 0.55);
\draw[fill=poporange!40,draw=poporange] (1.1,0) ellipse(0.14 and 1.67);
\draw[->,popdark] (-2.6,0) -- (2.8,0) node[right]{\small axis $x$};
\fill[popdark] (0,0) circle(1.5pt) node[below left]{\small $O$};
\draw[->,popmuted] (0,-0.35) -- (1.1,-0.35) node[midway,below]{\small $x$};
\draw[->,popmuted] (1.1,0) -- (1.1,1.67) node[midway,right]{\small $y$};
\draw[popmuted] (0,0) -- (1.1,1.67) node[midway,above left=-2pt]{\small $a$};
\node[poporange] at (1.45,1.95) {\small disc};
\end{tikzpicture}
\legende{Sphere sliced into discs perpendicular to a diameter; $y=\sqrt{a^2-x^2}$.}
\end{center}
\end{popfigure}

The disc's mass is $dm=\rho\pi y^2\,dx=\rho\pi(a^2-x^2)\,dx$, and its moment of inertia about the axis (its own central axis) is $\tfrac12 y^2\,dm$:
\begin{align*}
I&=\int_{-a}^{a}\frac12 y^2\,dm=\frac{\pi\rho}{2}\int_{-a}^{a}(a^2-x^2)^2\,dx\\
&=\frac{\pi\rho}{2}\int_{-a}^{a}\left(a^4-2a^2x^2+x^4\right)dx
=\frac{\pi\rho}{2}\left[a^4x-\frac{2a^2x^3}{3}+\frac{x^5}{5}\right]_{-a}^{a}\\
&=\frac{\pi\rho}{2}\cdot 2\left(a^5-\frac{2a^5}{3}+\frac{a^5}{5}\right)
=\pi\rho a^5\cdot\frac{15-10+3}{15}=\frac{8\pi\rho a^5}{15}.
\end{align*}
Replacing $\rho=\dfrac{3M}{4\pi a^3}$:
\[ \boxed{I=\frac{2}{5}Ma^2}. \]
This derivation is a classic ``show that'' question at the GCE; learn the slicing idea, not just the result.

\courssub{The Uniform Rectangular Lamina}

Let a uniform rectangular lamina have mass $M$, sides $2a$ (along $x$) and $2b$ (along $y$), and surface density $\sigma=\dfrac{M}{4ab}$.

\textbf{About the axis in its plane through $G$ parallel to the side $2b$.} Slice into strips parallel to that axis, of width $dx$ at distance $x$: each strip is itself a rod of mass $\sigma\cdot 2b\,dx$ with all points at distance $x$:
\[ I=\int_{-a}^{a}x^2\,\sigma\,2b\,dx=2b\sigma\,\frac{2a^3}{3}=\frac{4ab\sigma a^2}{3}=\tfrac13 Ma^2. \]
Similarly, about the axis parallel to the side $2a$: $I=\tfrac13 Mb^2$.

\textbf{About the perpendicular axis through $G$.} By the perpendicular axes theorem for a lamina, $I_z=I_x+I_y$:
\[ I=\tfrac13 M(a^2+b^2). \]

\courssub{Summary of Standard Results}

\begin{propriete}[Standard moments of inertia --- to memorise]
All bodies are uniform, of mass $M$; the axis passes through the centre of mass unless stated. Proofs by integration (as above) are examinable when the question says ``prove'', ``show'' or ``from first principles''.
\begin{center}
\small
\begin{tabularx}{\linewidth}{|l|X|c|c|}
\hline
\rowcolor{popblueL} \textbf{Body} & \textbf{Axis} & $\boldsymbol{I}$ & $\boldsymbol{k^2}$ \\
\hline
Rod, length $2a$ & perpendicular, through centre & $\tfrac13 Ma^2$ & $\tfrac13 a^2$ \\
\hline
Rod, length $2a$ & perpendicular, through one end & $\tfrac43 Ma^2$ & $\tfrac43 a^2$ \\
\hline
Ring / hoop, radius $a$ & central, perpendicular to plane & $Ma^2$ & $a^2$ \\
\hline
Disc / solid cylinder, radius $a$ & central axis & $\tfrac12 Ma^2$ & $\tfrac12 a^2$ \\
\hline
Solid sphere, radius $a$ & any diameter & $\tfrac25 Ma^2$ & $\tfrac25 a^2$ \\
\hline
Rectangular lamina, sides $2a,2b$ & in plane, through centre, $\parallel$ side $2b$ & $\tfrac13 Ma^2$ & $\tfrac13 a^2$ \\
\hline
Rectangular lamina, sides $2a,2b$ & perpendicular, through centre & $\tfrac13 M(a^2+b^2)$ & $\tfrac13(a^2+b^2)$ \\
\hline
Solid cylinder, radius $a$, height $h$ & diameter through centre & $M\left(\tfrac{a^2}{4}+\tfrac{h^2}{12}\right)$ & $\tfrac{a^2}{4}+\tfrac{h^2}{12}$ \\
\hline
\end{tabularx}
\end{center}
\end{propriete}

\begin{attention}[Exam technique]
In the GCE you may \emph{quote} any standard result from this table without re-deriving it --- unless the question explicitly says ``prove'', ``show by integration'' or ``from first principles''. When a proof is required, present the element, its mass, and the integral exactly as in the models above.
\end{attention}

\courssub{Composite Bodies: Adding and Subtracting}

\begin{methode}[Moments of inertia of composite bodies]
\begin{enumerate}
\item Split the body into standard parts (rods, discs, spheres, point masses).
\item Moments of inertia about the \emph{same} axis \textbf{add}: $I=I_1+I_2+\cdots$
\item For a \textbf{hole}, subtract the moment of inertia of the missing piece (computed about the same axis).
\item Use the parallel axes theorem $I=I_G+Md^2$ whenever a part's own centre is not on the axis.
\end{enumerate}
\end{methode}

\begin{exoresolu}[Disc with a circular hole]
A uniform circular disc of radius $2a$ has a concentric circular hole of radius $a$. The mass of the remaining body is $M$. Find its moment of inertia about the central axis perpendicular to its plane.
\tcblower
\textbf{Solution.} Treat the body as (full disc of radius $2a$) minus (disc of radius $a$). The surface density is uniform; the full disc has area $4\pi a^2$ and the hole $\pi a^2$, so the remaining area is $3\pi a^2$. Hence the full disc had mass $\tfrac43 M$ and the removed piece mass $\tfrac13 M$.
\begin{align*}
I_{\text{full}}&=\tfrac12\left(\tfrac43 M\right)(2a)^2=\tfrac83 Ma^2,\\
I_{\text{hole}}&=\tfrac12\left(\tfrac13 M\right)a^2=\tfrac16 Ma^2,\\
I&=\tfrac83 Ma^2-\tfrac16 Ma^2=\tfrac{15}{6}Ma^2=\boxed{\tfrac52 Ma^2}.
\end{align*}
Check: $k^2=\tfrac52 a^2$, between $a^2$ (ring) and $\tfrac12(2a)^2=2a^2$ (full disc) --- sensible since mass is concentrated at larger radii.
\end{exoresolu}

\begin{exoresolu}[Rod carrying two point masses]
A uniform rod $AB$ of length $2a$ and mass $3m$ carries a particle of mass $2m$ at $A$ and a particle of mass $m$ at its midpoint $G$. Find the moment of inertia of the system about a perpendicular axis through $A$, and the radius of gyration.
\tcblower
\textbf{Solution.} Add the contributions about the axis through $A$:
\begin{align*}
I_{\text{rod}}&=\tfrac43(3m)a^2=4ma^2,\\
I_{2m}&=2m\cdot 0^2=0 \quad(\text{on the axis}),\\
I_{m}&=m\,a^2.
\end{align*}
Total: $I=4ma^2+ma^2=5ma^2$. The total mass is $3m+2m+m=6m$, so
\[ k=\sqrt{\frac{I}{M}}=\sqrt{\frac{5ma^2}{6m}}=a\sqrt{\frac56}. \]
\end{exoresolu}

\begin{experience}[Rolling race --- an experimental insight]
Take a hoop, a solid disc, and a solid sphere, all of the same mass and radius. Release them from rest at the top of an inclined plane. The sphere reaches the bottom first, then the disc, then the hoop. This order reflects their moments of inertia: $I_{\text{sphere}}=\tfrac25 Ma^2$, $I_{\text{disc}}=\tfrac12 Ma^2$, $I_{\text{hoop}}=Ma^2$. The smaller the moment of inertia, the more gravitational potential energy is converted into translational kinetic energy rather than rotational kinetic energy, so the greater the linear acceleration. This experiment beautifully illustrates why the distribution of mass matters, not just the total mass.
\end{experience}

\begin{savaistu}[Historical note]
The concept of moment of inertia was introduced by Leonhard Euler in 1765 in his work \emph{Theoria motus corporum solidorum seu rigidorum}. Euler also derived the principal axes theorem and the equations of motion for a rigid body. The term ``moment of inertia'' was coined later; Euler called it ``momentum inertiae''. The parallel axes theorem is due to Christiaan Huygens (1673) in his study of compound pendulums, and the perpendicular axes theorem for laminas is often attributed to Jakob Steiner (1826).
\end{savaistu}

\begin{exercice}[Flywheel with a hole]
A uniform solid cylinder of mass $6m$ and radius $2a$ has a coaxial cylindrical hole of radius $a$ drilled through it. Show that the moment of inertia of the remaining body about its axis is $\dfrac{45}{4}ma^2$, and state the radius of gyration.
\end{exercice}

\begin{corrige}[Flywheel with a hole]
The cylinder \emph{before} drilling has mass $6m$ and radius $2a$. Its cross-sectional area is $\pi(2a)^2=4\pi a^2$. The hole has radius $a$, area $\pi a^2$. The remaining cross-sectional area is $3\pi a^2$. Since the density is uniform, masses are proportional to areas.
\[
\text{Mass of removed cylinder} = \frac{\pi a^2}{4\pi a^2}\times 6m = \frac{3}{2}m = 1.5m.
\]
\[
\text{Mass of remaining body} = 6m - 1.5m = 4.5m = \frac{9}{2}m.
\]
Now compute moments of inertia about the common axis:
\[
I_{\text{original}} = \frac12 (6m)(2a)^2 = 12ma^2.
\]
\[
I_{\text{hole}} = \frac12 \left(\frac{3}{2}m\right)a^2 = \frac{3}{4}ma^2.
\]
\[
I_{\text{remaining}} = 12ma^2 - \frac{3}{4}ma^2 = \frac{48-3}{4}ma^2 = \boxed{\frac{45}{4}ma^2}.
\]
The radius of gyration $k$ satisfies $I = M_{\text{rem}} k^2$:
\[
\frac{45}{4}ma^2 = \frac{9}{2}m\,k^2 \quad\Rightarrow\quad k^2 = \frac{45}{4}\cdot\frac{2}{9}a^2 = \frac{5}{2}a^2 \quad\Rightarrow\quad \boxed{k = a\sqrt{\frac{5}{2}}}.
\]
\end{corrige}

\begin{aretenir}[Key points]
\begin{itemize}
\item $I=\displaystyle\int r^2\,dm$; choose elements whose points are all at the same distance from the axis (strips, rings, discs).
\item Standard results: rod $\tfrac13 Ma^2$ (centre) and $\tfrac43 Ma^2$ (end); ring $Ma^2$; disc and solid cylinder $\tfrac12 Ma^2$; solid sphere $\tfrac25 Ma^2$.
\item Moments of inertia add for composite bodies; subtract for holes; shift with $I_G+Md^2$ when needed.
\item Quote standard results freely unless the question demands a proof by integration.
\end{itemize}
\end{aretenir}

\coursec{Equation of Motion under a Constant Torque}

In the previous sections we learned how to compute the moment of inertia $I$ of a rigid body by integration, and how to express the result compactly through the radius of gyration, $I = Mk^{2}$. The moment of inertia, however, is not an end in itself: it is the rotational analogue of mass, and its true role appears when we ask the dynamical question --- \emph{how does a rigid body respond when forces try to make it rotate?} Just as Newton's second law $F = ma$ answers this question for translation, there is a rotational counterpart, $C = I\alpha$, which is the central result of this section and one of the most frequently examined results of the whole chapter. We shall build it carefully, starting from the notion of the moment of a force, and then apply it to braking problems, to pulleys with inertia, and to bodies swinging under gravity.

\courssub{The Moment of a Force about a Fixed Axis}

\textbf{Intuition.} Everyone who has opened a door knows that the effect of a force on a rotating body depends not only on the size of the force but also on \emph{where} and \emph{in what direction} it is applied. Push near the hinges and the door barely moves; push at the handle, perpendicular to the door, and it swings open easily. Push along the door towards the hinges and nothing rotates at all. The quantity that measures this ``turning effect'' is the \cle{moment} of the force, also called the \cle{torque}.

\begin{definition}[Moment (torque) of a force about a fixed axis]
Let a force of magnitude $F$ act on a body free to rotate about a fixed axis $\Delta$, and suppose the force lies in a plane perpendicular to the axis. The \cle{moment} of the force about the axis is
\[
C = F\,d,
\]
where $d$ is the \cle{lever arm}: the perpendicular distance from the axis to the line of action of the force. The SI unit of moment is the newton-metre ($\mathrm{N\,m}$). The moment is given a \textbf{sign}: we choose a positive sense of rotation (usually anticlockwise) and count $C > 0$ if the force tends to turn the body in that sense, $C<0$ otherwise. A moment tending to slow a rotation is called a \emph{resisting couple} or \emph{braking torque}.
\end{definition}

\begin{attention}[Torque is not work]
Although the newton-metre is dimensionally the same as the joule, a torque is \emph{not} an energy. Never write a torque in joules, and never add a torque to an energy. Keep the unit $\mathrm{N\,m}$ for moments and reserve $\mathrm{J}$ for work and energy.
\end{attention}

If several forces act, the total moment about the axis is the algebraic sum $C = \sum F_i d_i$, each term carrying the sign of the sense in which it tends to turn the body. Note also that a force whose line of action passes \emph{through} the axis has $d = 0$ and hence zero moment: the reaction at a hinge or axle never appears in the moment equation. This is precisely why the moment equation is so useful: it lets us study the rotation \emph{without ever knowing the reaction at the axle}.

\begin{popfigure}
\begin{tikzpicture}[scale=1.05]
  % wheel
  \draw[thick, popblue] (0,0) circle (2);
  \fill[poppink] (0,0) circle (0.07);
  \node[below left, popink] at (0,0) {$O$};
  % point of application
  \fill[poporange] (2,0) circle (0.07);
  \node[right, popink] at (2.05,-0.15) {$A$};
  % lever arm
  \draw[thick, popgreen, ->] (0,0) -- (2,0) node[midway, below, popgreen] {$d$};
  % force tangential
  \draw[very thick, poporange, ->] (2,0) -- (2,1.8) node[right, poporange] {$\vec{F}$};
  % line of action dashed
  \draw[dashed, popmuted] (2,-1.2) -- (2,2.2);
  % right angle marker
  \draw[popmuted] (1.8,0) -- (1.8,0.2) -- (2,0.2);
  % rotation arrow
  \draw[thick, poppurple, ->] (150:2.6) arc (150:60:2.6);
  \node[poppurple] at (90:2.85) {sense of rotation};
  % axis label
  \node[popink] at (-2.5,-2.3) {fixed axis through $O$};
\end{tikzpicture}
\legende{A force $\vec{F}$ applied tangentially at $A$ to a wheel free to turn about a fixed axis through $O$. The lever arm is $d = OA$, and the moment is $C = Fd$, positive in the anticlockwise sense shown.}
\end{popfigure}

\courssub{Derivation of the Equation of Motion $C = I\alpha$}

Consider a rigid body rotating about a fixed axis through $O$ with angular acceleration $\alpha$. Think of the body as a collection of particles $P_i$ of mass $m_i$, each at a perpendicular distance $r_i$ from the axis. Because the body is rigid, every particle has the \emph{same} angular acceleration $\alpha$, and particle $P_i$ moves on a circle of radius $r_i$ with tangential acceleration $a_i = r_i\alpha$.

Apply Newton's second law to particle $P_i$ in the tangential direction. The tangential force on $P_i$ has two origins: external forces applied to the body (giving a tangential component $F_i$) and internal forces exerted by the other particles. Thus
\[
m_i r_i \alpha = F_i + (\text{tangential internal forces on } P_i).
\]
Multiply by $r_i$ and sum over all particles:
\[
\alpha \sum_i m_i r_i^{2} = \sum_i F_i r_i + \sum_i r_i \times (\text{internal forces}).
\]
Now two remarkable simplifications occur:

\begin{itemize}
\item By Newton's \emph{third} law, internal forces occur in equal and opposite pairs acting along the same line, so their total moment about the axis is \textbf{zero}: the last sum vanishes.
\item The sum $\sum_i m_i r_i^{2}$ is exactly the moment of inertia $I$ of the body about the axis, and $\sum_i F_i r_i$ is the total moment $C$ of the external forces about the axis.
\end{itemize}

\begin{propriete}[Equation of rotational motion about a fixed axis]
For a rigid body rotating about a \textbf{fixed axis}, if $C$ is the total moment of the external forces about the axis, $I$ the moment of inertia about the same axis, and $\alpha$ the angular acceleration, then
\[
\boxed{\,C = I\alpha\,}
\]
This is the rotational analogue of Newton's second law $F = ma$: torque plays the role of force, moment of inertia the role of mass, and angular acceleration the role of acceleration. The derivation above (from $F = ma$ applied to each particle, using Newton's third law to eliminate internal forces) is examinable and should be reproduced in full.
\end{propriete}

\begin{aretenir}[The translation--rotation dictionary]
\begin{center}
\begin{tabularx}{\linewidth}{|l|X|X|}
\hline
\rowcolor{popblueL} & \textbf{Translation} & \textbf{Rotation} \\
\hline
Displacement & $s$ (m) & $\theta$ (rad) \\
\hline
Velocity & $v = \dot{s}$ & $\omega = \dot{\theta}$ \\
\hline
Acceleration & $a = \dot{v}$ & $\alpha = \dot{\omega}$ \\
\hline
Inertia & mass $m$ & moment of inertia $I$ \\
\hline
Cause of motion & force $F$ & moment $C = Fd$ \\
\hline
Second law & $F = ma$ & $C = I\alpha$ \\
\hline
\end{tabularx}
\end{center}
\end{aretenir}

\courssub{Constant Torque: the Kinematics of Rotation}

When the total torque $C$ is \emph{constant} and the body is rigid (so $I$ is constant), the equation $C = I\alpha$ gives a \textbf{constant angular acceleration} $\alpha = C/I$. Rotational motion with constant $\alpha$ is governed by formulae exactly parallel to those of uniformly accelerated linear motion, obtained by integrating $\dot{\omega} = \alpha$:

\begin{propriete}[Constant angular acceleration formulae]
If $\alpha$ is constant, $\omega_0$ the initial angular velocity and $\theta$ the angle turned through in time $t$, then
\[
\omega = \omega_0 + \alpha t, \qquad
\theta = \omega_0 t + \tfrac{1}{2}\alpha t^{2}, \qquad
\omega^{2} = \omega_0^{2} + 2\alpha\theta .
\]
Angles must be in \textbf{radians}, angular velocities in $\mathrm{rad\,s^{-1}}$, angular accelerations in $\mathrm{rad\,s^{-2}}$.
\end{propriete}

\begin{attention}[Keep the units consistent]
Before substituting into $C = I\alpha$, convert everything to SI units: torque in $\mathrm{N\,m}$, moment of inertia in $\mathrm{kg\,m^{2}}$, $\alpha$ in $\mathrm{rad\,s^{-2}}$. The most frequent error is to leave a rotation speed in revolutions per minute: convert with $\omega = \dfrac{2\pi N}{60}\ \mathrm{rad\,s^{-1}}$ where $N$ is in $\mathrm{rev\,min^{-1}}$. Similarly, a number of revolutions $n$ corresponds to an angle $\theta = 2\pi n$ radians.
\end{attention}

\begin{exoresolu}[Flywheel brought to rest by a braking couple]
A flywheel of moment of inertia $I = 0.8\ \mathrm{kg\,m^{2}}$ rotates at $300\ \mathrm{rev\,min^{-1}}$. A constant resisting couple of magnitude $2\ \mathrm{N\,m}$ is applied. Find (a) the angular acceleration, (b) the time taken to stop, (c) the number of revolutions made before stopping.
\tcblower
\textbf{Solution.} Convert the initial speed:
\[
\omega_0 = \frac{2\pi \times 300}{60} = 10\pi\ \mathrm{rad\,s^{-1}} \approx 31.4\ \mathrm{rad\,s^{-1}}.
\]
(a) The couple resists the motion, so with the sense of rotation positive, $C = -2\ \mathrm{N\,m}$. Then
\[
\alpha = \frac{C}{I} = \frac{-2}{0.8} = -2.5\ \mathrm{rad\,s^{-2}}.
\]
(b) Stopping means $\omega = 0$. From $\omega = \omega_0 + \alpha t$:
\[
0 = 10\pi - 2.5\,t \quad\Longrightarrow\quad t = \frac{10\pi}{2.5} = 4\pi \approx 12.6\ \mathrm{s}.
\]
(c) Use $\omega^{2} = \omega_0^{2} + 2\alpha\theta$ with $\omega = 0$:
\[
\theta = \frac{-\omega_0^{2}}{2\alpha} = \frac{-(10\pi)^{2}}{2(-2.5)} = \frac{100\pi^{2}}{5} = 20\pi^{2}\ \mathrm{rad}.
\]
The number of revolutions is
\[
n = \frac{\theta}{2\pi} = \frac{20\pi^{2}}{2\pi} = 10\pi \approx 31.4\ \text{revolutions}.
\]
Note how the sign convention (resisting couple negative) automatically produces a negative $\alpha$ and a positive stopping angle.
\end{exoresolu}

\begin{popfigure}
\begin{tikzpicture}
\begin{axis}[
  width=10.5cm, height=6.5cm,
  xlabel={$t$ (s)}, ylabel={$\omega$ (rad\,s$^{-1}$)},
  xmin=0, xmax=15, ymin=0, ymax=36,
  xtick={0,4,8,12,12.566}, xticklabels={$0$,$4$,$8$,$12$,$4\pi$},
  ytick={0,10,20,31.4}, yticklabels={$0$,$10$,$20$,$10\pi$},
  grid=both, grid style={popline, very thin},
  axis line style={popink},
]
\addplot[very thick, popblue, domain=0:12.566] {31.416 - 2.5*x};
\addplot[dashed, popmuted] coordinates {(12.566,0) (12.566,31.416)};
\addplot[dashed, popmuted] coordinates {(0,31.416) (12.566,31.416)};
\addplot[only marks, mark=*, poporange, mark size=2pt] coordinates {(12.566,0)};
\node[poporange, anchor=south west] at (axis cs:12.7,1) {stops at $t = 4\pi$ s};
\node[popblue, anchor=south west] at (axis cs:1,26) {$\omega = 10\pi - 2.5t$};
\end{axis}
\end{tikzpicture}
\legende{Angular velocity decreasing linearly under a constant braking couple: $\omega(t)$ is a straight line of slope $\alpha = -2.5\ \mathrm{rad\,s^{-2}}$, reaching zero at the stopping time $t = 4\pi\ \mathrm{s}$.}
\end{popfigure}

\courssub{The Pulley with Non-negligible Moment of Inertia}

In Mechanics you solved Atwood machines with a ``light, frictionless pulley'', for which the tension is the same on both sides. When the pulley has a \emph{non-negligible moment of inertia} $I$ and radius $r$, this is no longer true: a difference of tensions is precisely what provides the torque that spins the pulley up. This is one of the most classic GCE traps.

\begin{attention}[With a massive pulley, $T_1 \neq T_2$]
If the pulley has inertia and the string does not slip, the tensions on the two sides are \textbf{different}. Writing $T_1 = T_2$ forces the net torque on the pulley to be zero, contradicting the fact that the pulley angularly accelerates. Always introduce two distinct unknowns $T_1$ and $T_2$.
\end{attention}

\begin{methode}[Solving pulley-with-inertia problems]
\begin{enumerate}
\item \textbf{Draw three free-body diagrams}: mass $m_1$, mass $m_2$, and the pulley. Mark the two tensions $T_1$, $T_2$ as distinct.
\item \textbf{Write $F = ma$ for each mass} along its direction of motion:
\[
m_1 g - T_1 = m_1 a, \qquad T_2 - m_2 g = m_2 a \quad (m_1 > m_2).
\]
\item \textbf{Write $C = I\alpha$ for the pulley}: the axle reaction passes through the axis so it has no moment; only the tensions contribute:
\[
(T_1 - T_2)\,r = I\alpha .
\]
\item \textbf{Add the no-slip condition} linking string and pulley:
\[
a = r\alpha .
\]
\item \textbf{Solve the four equations} for $a$, $\alpha$, $T_1$, $T_2$. A quick route: substitute $T_1$, $T_2$ from step 2 into step 3 and use step 4 to obtain
\[
a = \frac{(m_1 - m_2)g}{m_1 + m_2 + I/r^{2}} .
\]
\end{enumerate}
\end{methode}

\begin{popfigure}
\begin{tikzpicture}[scale=1.0]
  % support
  \draw[thick, popink] (-1.5,3.4) -- (1.5,3.4);
  \draw[popmuted] (-1.3,3.4) -- (-1.0,3.8) (-0.7,3.4) -- (-0.4,3.8) (-0.1,3.4) -- (0.2,3.8) (0.5,3.4) -- (0.8,3.8);
  \draw[thick, popink] (0,3.4) -- (0,2.6);
  % pulley
  \draw[thick, popblue, fill=popblueL] (0,2) circle (0.6);
  \fill[poppink] (0,2) circle (0.05);
  \node[popink] at (0.35,2.35) {$I,\ r$};
  % strings
  \draw[thick, popink] (-0.6,2) -- (-0.6,0.4);
  \draw[thick, popink] (0.6,2) -- (0.6,-0.6);
  % masses
  \draw[thick, popgreen, fill=popgreenL] (-0.95,0.4) rectangle (-0.25,-0.3);
  \node[popink] at (-0.6,0.05) {$m_1$};
  \draw[thick, poporange, fill=white] (0.25,-0.6) rectangle (0.95,-1.3);
  \node[popink] at (0.6,-0.95) {$m_2$};
  % tensions
  \draw[thick, poppurple, ->] (-0.95,0.2) -- (-0.95,1.3) node[left, poppurple] {$T_1$};
  \draw[thick, poppurple, ->] (0.95,-0.8) -- (0.95,0.4) node[right, poppurple] {$T_2$};
  % weights
  \draw[thick, popdark, ->] (-0.6,-0.3) -- (-0.6,-1.1) node[below, popdark] {$m_1 g$};
  \draw[thick, popdark, ->] (0.6,-1.3) -- (0.6,-1.9) node[below, popdark] {$m_2 g$};
  % rotation arrow on pulley
  \draw[thick, poppurple, ->] (140:0.95) arc (140:40:0.95);
  % acceleration arrows
  \draw[thick, popgreen, ->] (-1.35,0.05) -- (-1.35,-0.75) node[left, popgreen] {$a$};
  \draw[thick, poporange, ->] (1.35,-0.95) -- (1.35,-0.15) node[right, poporange] {$a$};
\end{tikzpicture}
\legende{Atwood machine with a massive pulley of moment of inertia $I$ and radius $r$. The tensions $T_1$ and $T_2$ differ: their difference $(T_1 - T_2)r$ supplies the torque $I\alpha$, and the no-slip condition gives $a = r\alpha$.}
\end{popfigure}

\begin{exoresolu}[Atwood machine with a massive pulley]
Two masses $m_1 = 4\ \mathrm{kg}$ and $m_2 = 2\ \mathrm{kg}$ hang from a string passing over a pulley of radius $r = 0.2\ \mathrm{m}$ and moment of inertia $I = 0.08\ \mathrm{kg\,m^{2}}$. The string does not slip and the axle is frictionless. Take $g = 9.8\ \mathrm{m\,s^{-2}}$. Find the acceleration of the masses and the two tensions.
\tcblower
\textbf{Solution.} With $m_1$ descending, the three equations and the no-slip condition are
\[
4g - T_1 = 4a, \qquad T_2 - 2g = 2a, \qquad (T_1 - T_2)(0.2) = 0.08\,\alpha, \qquad a = 0.2\,\alpha .
\]
From the pulley equation, $T_1 - T_2 = \dfrac{0.08}{0.2}\alpha = 0.4\,\alpha$, and since $\alpha = \dfrac{a}{0.2} = 5a$, we get $T_1 - T_2 = 2a$.
Adding the two mass equations: $2g - (T_1 - T_2) = 6a$, so
\[
2g - 2a = 6a \quad\Longrightarrow\quad a = \frac{2g}{8} = \frac{19.6}{8} = 2.45\ \mathrm{m\,s^{-2}}.
\]
Then
\[
T_1 = 4(g - a) = 4(9.8 - 2.45) = 29.4\ \mathrm{N}, \qquad
T_2 = 2(g + a) = 2(9.8 + 2.45) = 24.5\ \mathrm{N}.
\]
Check: $T_1 - T_2 = 4.9\ \mathrm{N} = 2a$ \checkmark, and $\alpha = a/r = 12.25\ \mathrm{rad\,s^{-2}}$. Note that $T_1 \neq T_2$: the difference of $4.9\ \mathrm{N}$ is what accelerates the pulley. Observe also that the denominator of the general formula is $m_1 + m_2 + I/r^{2} = 4 + 2 + 2 = 8\ \mathrm{kg}$: the pulley behaves like an extra \emph{effective mass} $I/r^{2} = 2\ \mathrm{kg}$.
\end{exoresolu}

\courssub{Rigid Body Rotating under Gravity about a Fixed Horizontal Axis}

Consider now a rigid body (a rod, a lamina, any pendulum-like object) free to rotate about a fixed \emph{horizontal} axis, displaced from its equilibrium position and released. Which forces have a moment about the axis? The reaction at the axis passes through it, so its moment is zero. The weight $Mg$ acts at the centre of gravity $G$; if $G$ is at distance $h$ from the axis and the line $OG$ makes an angle $\theta$ with the downward vertical, the lever arm of the weight is the \emph{horizontal distance} of $G$ from the axis, namely $h\sin\theta$. Hence
\[
C = -Mg\,h\sin\theta,
\]
the minus sign showing that gravity always pulls $G$ back towards the vertical (a \emph{restoring} moment). The equation of motion is therefore
\[
I\ddot{\theta} = -Mgh\sin\theta .
\]
For small oscillations this will reduce to simple harmonic motion --- that is the subject of the compound pendulum, treated in Section 6. For now, note two things: the equation is exact (no small-angle assumption is needed to \emph{write} it), and when the body is released from rest in a given position $\theta_0$, the initial angular acceleration is $\alpha_0 = -Mgh\sin\theta_0 / I$.

\begin{exemplebox}[A uniform rod released from the horizontal]
A uniform rod $AB$ of length $2l$ and mass $M$ is free to rotate about a horizontal axis through $A$. It is held horizontal and released. With $I_A = \frac{4}{3}Ml^{2}$ (standard result from Section 2) and $G$ at distance $l$ from $A$, the weight has lever arm $l$ in the horizontal position, so
\[
\frac{4}{3}Ml^{2}\,\alpha_0 = Mgl \quad\Longrightarrow\quad \alpha_0 = \frac{3g}{4l}.
\]
The tangential acceleration of the free end $B$ at release is $2l\,\alpha_0 = \tfrac{3}{2}g$ --- greater than $g$! This surprising result is a favourite examiner's question: the end of a falling rod initially accelerates \emph{faster} than a freely falling particle.
\end{exemplebox}

\begin{savaistu}[Why flywheels matter]
The equation $C = I\alpha$ explains why engineers fit heavy flywheels to engines and presses: a large $I$ means a small $\alpha$ for any disturbing torque, so the rotation speed stays almost constant between power strokes. The same principle steadies the potter's wheel and the traditional grinding stones used in many Cameroonian workshops: once spinning, a massive stone resists any sudden change of speed.
\end{savaistu}

\begin{exercice}[Braking a grindstone]
A grindstone of moment of inertia $0.5\ \mathrm{kg\,m^{2}}$ runs at $600\ \mathrm{rev\,min^{-1}}$. A tool is pressed against it, producing a constant resisting couple of $5\ \mathrm{N\,m}$. Find (a) the angular retardation, (b) the time to stop, (c) the number of revolutions completed while stopping.
\end{exercice}

\begin{corrige}[Exercise 1]
$\omega_0 = 2\pi(600)/60 = 20\pi\ \mathrm{rad\,s^{-1}}$.
(a) $\alpha = C/I = -5/0.5 = -10\ \mathrm{rad\,s^{-2}}$.
(b) $t = \omega_0/|\alpha| = 20\pi/10 = 2\pi \approx 6.28\ \mathrm{s}$.
(c) $\theta = \omega_0^{2}/(2|\alpha|) = 400\pi^{2}/20 = 20\pi^{2}\ \mathrm{rad}$, so $n = \theta/(2\pi) = 10\pi \approx 31.4$ revolutions.
\end{corrige}

\begin{exercice}[Massive pulley]
In the worked Atwood example above, the pulley is replaced by one of radius $0.1\ \mathrm{m}$ and moment of inertia $0.02\ \mathrm{kg\,m^{2}}$, the masses being unchanged. Find the new acceleration and the two tensions.
\end{exercice}

\begin{corrige}[Exercise 2]
$T_1 - T_2 = \dfrac{I}{r^{2}}a = \dfrac{0.02}{0.01}a = 2a$ as before, so $a = \dfrac{2g}{6 + 2} = 2.45\ \mathrm{m\,s^{-2}}$ (the ratio $I/r^{2}$ is unchanged), $T_1 = 29.4\ \mathrm{N}$, $T_2 = 24.5\ \mathrm{N}$. The lesson: what matters dynamically is the \emph{effective mass} $I/r^{2}$ of the pulley.
\end{corrige}

\begin{exercice}[Rod released from rest]
A uniform rod of length $1.2\ \mathrm{m}$ and mass $3\ \mathrm{kg}$ is free to rotate about a horizontal axis through one end. It is released from rest in the horizontal position. Taking $g = 9.8\ \mathrm{m\,s^{-2}}$, find (a) the moment of inertia about the axis, (b) the initial angular acceleration, (c) the initial tangential acceleration of the free end.
\end{exercice}

\begin{corrige}[Exercise 3]
Here $2l = 1.2\ \mathrm{m}$, so $l = 0.6\ \mathrm{m}$.
(a) $I_A = \tfrac{4}{3}Ml^{2} = \tfrac{4}{3}(3)(0.36) = 1.44\ \mathrm{kg\,m^{2}}$.
(b) $\alpha_0 = \dfrac{Mgl}{I_A} = \dfrac{3 \times 9.8 \times 0.6}{1.44} = 12.25\ \mathrm{rad\,s^{-2}}$ (equivalently $\alpha_0 = \dfrac{3g}{4l} = \dfrac{29.4}{2.4} = 12.25$).
(c) $a_B = 2l\,\alpha_0 = 1.2 \times 12.25 = 14.7\ \mathrm{m\,s^{-2}} = \tfrac{3}{2}g$, greater than $g$ as expected.
\end{corrige}

\begin{aretenir}[Key points of this section]
\begin{itemize}
\item Moment of a force about a fixed axis: $C = Fd$, with a sign fixed by the chosen positive sense of rotation; forces through the axis have zero moment.
\item Equation of motion of a rigid body about a fixed axis: $C = I\alpha$, derived from $F = ma$ particle by particle, internal forces cancelling by Newton's third law.
\item Constant torque $\Rightarrow$ constant $\alpha$: use $\omega = \omega_0 + \alpha t$, $\theta = \omega_0 t + \tfrac12 \alpha t^{2}$, $\omega^{2} = \omega_0^{2} + 2\alpha\theta$, with angles in radians.
\item Braking problems: resisting couple negative; stopping time from $\omega = 0$; revolutions from $n = \theta/(2\pi)$.
\item Massive pulley: $T_1 \neq T_2$; solve $F = ma$ for each mass, $C = I\alpha$ for the pulley, and $a = r\alpha$ for the no-slip link; the pulley contributes an effective mass $I/r^{2}$.
\item Gravity on a body about a horizontal axis: $C = -Mgh\sin\theta$, where $h\sin\theta$ is the horizontal distance of $G$ from the axis.
\end{itemize}
\end{aretenir}

This section has given us the dynamical law of rotation. In the next section we shall multiply this law by time to obtain the \emph{moment of momentum} (angular momentum) and its conservation --- the rotational analogue of momentum and impulse --- which governs collisions and sudden changes in rotating systems.

\coursec{Moment of Momentum and Its Conservation; Impulse}

In the previous section we established the fundamental equation of rotational motion, $G = I\ddot{\theta}$, which plays for rotation the role that Newton's second law $\vec{F} = m\vec{a}$ plays for translation. In linear mechanics, rewriting Newton's second law as $\vec{F} = \dfrac{d}{dt}(m\vec{v})$ leads to the notion of \emph{momentum} and to two of the most powerful tools of dynamics: the \emph{impulse--momentum equation} and the \emph{conservation of linear momentum}. Exactly the same programme can be carried out for rotation. Integrating $G = I\ddot{\theta}$ with respect to time will give us the \emph{moment of momentum} (also called \emph{angular momentum}), the \emph{impulsive couple} equation, and the celebrated \emph{law of conservation of moment of momentum}. These tools are indispensable whenever forces act for a very short time — collisions, jerks, strings suddenly becoming taut — situations in which the GCE examiner delights.

\courssub{Moment of Momentum of a Particle about an Axis}

\textbf{An observation to build intuition.}\par
Imagine a small stone of mass $m$ whirled at the end of a string on a smooth horizontal table. The stone moves fast, and the faster it moves, or the longer the string, the harder it is to stop the whirling motion. The ``quantity of rotational motion'' clearly depends on three things: the mass $m$, the speed $v$, and the distance from the centre. But which distance matters? If you push a door, the effect of your push depends on the \emph{perpendicular} distance from the line of the push to the hinge — a push aimed straight at the hinge produces no rotation at all, however large the force. The same is true for momentum: only the perpendicular distance from the axis to the line of the velocity counts.

\begin{popfigure}
\begin{tikzpicture}[scale=1.1, >=stealth]
  \draw[popline, dashed] (0,-1.2) -- (0,3.2) node[above, popmuted] {axis through $O$};
  \fill[popdark] (0,0) circle (2pt) node[below left, popdark] {$O$};
  \fill[poporange] (2.6,1.6) circle (3.5pt) node[right, popdark] {$m$};
  \draw[->, very thick, popblue] (2.6,1.6) -- (4.4,2.6) node[right, popblue] {$\vec{v}$};
  \draw[thick, popgreen] (0,0) -- (2.6,1.6);
  \draw[dashed, popmuted] (0,0) -- (3.52,2.12);
  \draw[dashed, popmuted] (3.52,2.12) -- (4.4,2.6);
  \node[popgreen] at (1.15,1.05) {$\vec{r}$};
  \draw[<->, thick, poppurple] (0.62,0.39) -- (3.0,1.85);
  \node[poppurple] at (1.6,1.95) {$r_{\perp}$};
  \node[popmuted, align=center, text width=4.6cm] at (5.6,0.2) {$r_{\perp}$ = perpendicular\\ distance from the axis\\ to the line of $\vec{v}$};
\end{tikzpicture}
\legende{A particle of mass $m$ moving with velocity $\vec{v}$; its moment of momentum about the axis through $O$ is $m v r_{\perp}$.}
\end{popfigure}

\begin{definition}[Moment of momentum of a particle about a fixed axis]
Let a particle of mass $m$ move with velocity $\vec{v}$ in a plane perpendicular to a fixed axis through $O$. The \cle{moment of momentum} (or \cle{angular momentum}) of the particle about the axis is
\[
L = m v r_{\perp},
\]
where $r_{\perp}$ is the \emph{perpendicular} distance from the axis to the line of action of $\vec{v}$. The sign of $L$ is fixed by a chosen positive sense of rotation. The SI unit of moment of momentum is the $\mathrm{kg\,m^{2}\,s^{-1}}$.
\end{definition}

\begin{attention}[Perpendicular distance, not position distance]
A very common error is to use the distance $|\vec{r}|$ from $O$ to the particle instead of the perpendicular distance from the axis to the \emph{line} of $\vec{v}$. If the particle moves \emph{towards} $O$ (radially), its moment of momentum about the axis is \textbf{zero}, whatever its speed. Always drop a perpendicular from $O$ onto the line of the velocity.
\end{attention}

\courssub{Moment of Momentum of a Rigid Body: $L = I\omega$}

Consider now a rigid body rotating with angular velocity $\omega$ about a fixed axis. Every particle $P_i$ of the body, of mass $m_i$ at perpendicular distance $r_i$ from the axis, moves in a circle of radius $r_i$ with speed $v_i = r_i\omega$. Since the velocity is tangential, the perpendicular distance from the axis to the line of $\vec{v_i}$ is exactly $r_i$. Hence the moment of momentum of the particle is
\[
L_i = m_i v_i r_i = m_i r_i^{2}\,\omega .
\]
The moments of momentum of all the particles have the same sense, so they add:
\[
L = \sum_i m_i r_i^{2}\,\omega = \omega \sum_i m_i r_i^{2}.
\]
But $\sum_i m_i r_i^{2}$ is precisely the moment of inertia $I$ of the body about the axis. We have therefore proved the following fundamental result.

\begin{propriete}[Moment of momentum of a rigid body about a fixed axis]
A rigid body of moment of inertia $I$ rotating with angular velocity $\omega$ about a fixed axis has moment of momentum about that axis
\[
\boxed{L = I\omega}.
\]
Unit: $\mathrm{kg\,m^{2}\,s^{-1}}$. This is the rotational analogue of the linear momentum $p = mv$. The derivation above (summing $m_i r_i^2 \omega$) is examinable and should be reproducible.
\end{propriete}

\begin{aretenir}[The translation--rotation dictionary]
\begin{center}
\begin{tabularx}{\linewidth}{|l|X|X|}
\hline
\rowcolor{popblueL} \textbf{Quantity} & \textbf{Translation} & \textbf{Rotation} \\
\hline
Displacement & $s$ & $\theta$ \\
\hline
Velocity & $v$ & $\omega$ \\
\hline
Inertia & mass $m$ & moment of inertia $I$ \\
\hline
Momentum & $p = mv$ & $L = I\omega$ \\
\hline
Force / torque & $F = \dfrac{dp}{dt}$ & $G = \dfrac{dL}{dt}$ \\
\hline
Impulse & $J = \Delta(mv)$ & impulsive couple $= \Delta(I\omega)$ \\
\hline
Kinetic energy & $\tfrac{1}{2}mv^{2}$ & $\tfrac{1}{2}I\omega^{2}$ \\
\hline
\end{tabularx}
\end{center}
\end{aretenir}

Indeed, the torque equation itself can be rewritten in momentum form: since $G = I\ddot{\theta} = \dfrac{d}{dt}(I\omega)$ when $I$ is constant,
\[
G = \frac{dL}{dt},
\]
which is the exact rotational analogue of $\vec{F} = \dfrac{d\vec{p}}{dt}$: \emph{the resultant external torque about a fixed axis equals the rate of change of the moment of momentum about that axis}.

\courssub{Conservation of Moment of Momentum}

From $G = \dfrac{dL}{dt}$ an immediate and enormously useful consequence follows.

\begin{propriete}[Conservation of moment of momentum]
If the resultant external torque about a fixed axis is zero during some interval, then the moment of momentum of the system about that axis remains constant:
\[
G = 0 \quad\Longrightarrow\quad I\omega = \text{constant}.
\]
For a system whose moment of inertia changes from $I_1$ to $I_2$, this reads
\[
I_1\omega_1 = I_2\omega_2 .
\]
(Proof: $G = \dfrac{dL}{dt} = 0$ implies $L$ constant. The statement is examinable.)
\end{propriete}

\begin{attention}[Choose the axis once, and state it]
Conservation of moment of momentum holds about \emph{one chosen fixed axis}, and only if the external torques \emph{about that same axis} vanish. In every examination answer, write explicitly: ``Taking moments of momentum about the fixed axis through $O$ \dots''. Marks are routinely lost by candidates who never state their axis.
\end{attention}

\begin{savaistu}[The skater and the diver]
Conservation of moment of momentum is the secret of the figure skater's spin: with arms stretched wide, $I$ is large and $\omega$ modest; by pulling the arms in, $I$ decreases sharply, so $\omega$ must increase to keep $I\omega$ constant — the spin becomes spectacularly fast. High divers use the same trick, tucking into a ball to multiply their rotation rate before straightening for entry into the water. No external torque about the vertical axis is needed: the rotation speeds up ``for free''.
\end{savaistu}

\begin{exoresolu}[A point mass dropped onto a rotating turntable]
A horizontal disc of mass $M = 2\ \mathrm{kg}$ and radius $a = 0.4\ \mathrm{m}$ rotates freely about a fixed vertical axis through its centre at $\omega_1 = 6\ \mathrm{rad\,s^{-1}}$. A small lump of clay of mass $m = 0.5\ \mathrm{kg}$ is dropped vertically onto the disc at a distance $r = 0.3\ \mathrm{m}$ from the axis and sticks to it. Find the new angular velocity.
\tcblower
\textbf{Solution.} The clay falls vertically, so it has no moment of momentum about the vertical axis before sticking. The forces during the impact are internal to the system (disc $+$ clay) or act along the axis; hence the external torque about the axis is zero and moment of momentum is conserved.
\[
I_{\text{disc}} = \tfrac{1}{2}Ma^{2} = \tfrac{1}{2}(2)(0.4)^{2} = 0.16\ \mathrm{kg\,m^{2}}, \qquad
I_{\text{clay}} = mr^{2} = (0.5)(0.3)^{2} = 0.045\ \mathrm{kg\,m^{2}}.
\]
Conservation about the axis: $I_1\omega_1 = I_2\omega_2$:
\[
0.16 \times 6 = (0.16 + 0.045)\,\omega_2
\quad\Longrightarrow\quad
\omega_2 = \frac{0.96}{0.205} \approx 4.68\ \mathrm{rad\,s^{-1}}.
\]
Note that the kinetic energy falls from $\tfrac{1}{2}(0.16)(36) = 2.88\ \mathrm{J}$ to $\tfrac{1}{2}(0.205)(4.68)^{2} \approx 2.25\ \mathrm{J}$: energy is lost in the inelastic ``sticking'', exactly as in a linear collision.
\end{exoresolu}

\begin{popfigure}
\begin{tikzpicture}[scale=0.95, >=stealth]
  % Before
  \draw[fill=popblueL, draw=popblue, thick] (0,0) ellipse (1.5 and 0.55);
  \fill[popdark] (0,0) circle (1.5pt);
  \draw[->, very thick, poporange] (1.9,0.15) arc (8:-60:1.9 and 0.75);
  \node[poporange] at (2.35,-0.75) {$\omega_1$};
  \node[popdark, align=center] at (0,-1.3) {Before: $I_1 = \tfrac12 Ma^2$};
  % arrow
  \draw[->, ultra thick, popmuted] (3.1,-0.3) -- (4.3,-0.3);
  \node[popmuted] at (3.7,0.15) {clay};
  % After
  \begin{scope}[xshift=7.4cm]
    \draw[fill=popblueL, draw=popblue, thick] (0,0) ellipse (1.5 and 0.55);
    \fill[popdark] (0,0) circle (1.5pt);
    \fill[popgreen] (0.9,0.28) circle (4pt) node[above right=-2pt, popgreen] {$m$};
    \draw[->, very thick, poporange] (1.9,0.15) arc (8:-60:1.9 and 0.75);
    \node[poporange] at (2.35,-0.75) {$\omega_2$};
    \node[popdark, align=center] at (0,-1.3) {After: $I_2 = I_1 + mr^2$};
  \end{scope}
  \node[popdark] at (3.7,-1.9) {$I_1\omega_1 = I_2\omega_2$};
\end{tikzpicture}
\legende{A lump of clay dropped onto a freely rotating disc: the moment of inertia increases, so the angular velocity decreases, with $I\omega$ constant.}
\end{popfigure}

\courssub{Impulsive Couples and the Angular Impulse--Momentum Equation}

In linear mechanics, when a very large force acts for a very short time, we bypass the details of the force and work with the \emph{impulse} $J = \int F\,dt = \Delta(mv)$. The same device works for rotation. Integrating $G = \dfrac{dL}{dt}$ over the short duration of the blow:
\[
\int G\,dt = \Delta(I\omega) = I\omega_2 - I\omega_1 .
\]

\begin{definition}[Impulsive couple]
The \cle{impulsive couple} (or \cle{angular impulse}) of a torque about a fixed axis is $\displaystyle\int G\,dt$. If an impulsive force of impulse $J$ acts at perpendicular distance $d$ from the axis, its impulsive couple about the axis is $Jd$. The impulsive couple equals the change of moment of momentum:
\[
\boxed{\text{impulsive couple} = \Delta(I\omega)}.
\]
Unit: $\mathrm{kg\,m^{2}\,s^{-1}}$ (equivalently $\mathrm{N\,m\,s}$).
\end{definition}

\begin{methode}[Solving impulsive rotation problems]
\begin{enumerate}
  \item \textbf{Choose and state the fixed axis} about which moments of momentum and impulsive couples will be taken.
  \item \textbf{Draw the impulsive forces}: impulsive tensions in strings, impulsive reactions at hinges or axles, impulsive contact forces. Ignore all finite forces (weights, ordinary tensions) during the infinitesimal impact.
  \item \textbf{Apply to each rotating body}: (impulsive couple about the axis) $= \Delta(I\omega)$.
  \item \textbf{Apply to each translating particle}: (impulse) $= \Delta(mv)$ along the direction of motion.
  \item \textbf{Add the kinematic constraints}: inextensible string $\Rightarrow$ equal speeds; string on a pulley rim $\Rightarrow$ $v = a\omega$; no slipping.
  \item \textbf{Solve} the simultaneous equations for the unknown impulses and velocities.
\end{enumerate}
\end{methode}

\begin{attention}[Kinetic energy is NOT conserved in impacts]
In any problem involving an impulsive blow, a collision, or a string suddenly becoming taut, kinetic energy is almost always \emph{lost} (converted to heat, sound, deformation). Never write $\tfrac12 I\omega_1^2 = \tfrac12 I\omega_2^2$ across an impact. The correct conserved quantities are \emph{linear momentum} (when no external impulse acts) and \emph{moment of momentum} (when no external impulsive couple acts about the chosen axis). Energy methods may be used \emph{before} and \emph{after} the impact, but not across it. Always check: is there an impact? If yes, momentum methods, not energy.
\end{attention}

\courssub{Connected Particles and Impulsive Tensions}

A classic GCE situation: two particles hang from a string passing over a pulley \emph{with inertia}; the system is held at rest with the string slack, or one particle is given a jerk, and the string suddenly becomes taut. During the jerk, \emph{impulsive tensions} appear in the two portions of the string (they differ, because the pulley has inertia and needs a net impulsive couple to be set spinning), and an \emph{impulsive reaction} appears at the axle. We treat each element separately.

\begin{popfigure}
\begin{tikzpicture}[scale=1.0, >=stealth]
  % pulley
  \draw[fill=popblueL, draw=popblue, thick] (0,0) circle (0.8);
  \fill[popdark] (0,0) circle (2pt) node[below right=1pt, popdark] {$O$};
  \node[popblue] at (0,1.15) {pulley, $I$};
  % strings
  \draw[thick, popdark] (-0.8,0) -- (-0.8,-2.2);
  \draw[thick, popdark] (0.8,0) -- (0.8,-2.6);
  % masses
  \draw[fill=popgreenL, draw=popgreen, thick] (-1.15,-2.7) rectangle (-0.45,-2.2);
  \node[popdark] at (-0.8,-2.45) {$m_1$};
  \draw[fill=popgreenL, draw=popgreen, thick] (0.45,-3.1) rectangle (1.15,-2.6);
  \node[popdark] at (0.8,-2.85) {$m_2$};
  % impulsive tensions
  \draw[->, very thick, poporange] (-0.8,-1.9) -- (-0.8,-1.1) node[midway, left, poporange] {$J_1$};
  \draw[->, very thick, poporange] (0.8,-2.3) -- (0.8,-1.5) node[midway, right, poporange] {$J_2$};
  % velocities
  \draw[->, very thick, poppurple] (-1.5,-2.45) -- (-1.5,-1.65) node[midway, left, poppurple] {$u$};
  \draw[->, very thick, poppurple] (1.5,-2.85) -- (1.5,-3.65) node[midway, right, poppurple] {$u$};
  % reaction
  \draw[->, very thick, poppink] (0,0) -- (0,1.0) node[right, poppink] {$R$};
  \node[popmuted, align=center, text width=4.2cm] at (3.6,-1.6) {String suddenly taut:\\ impulsive tensions $J_1 \ne J_2$,\\ impulsive reaction $R$ at axle};
\end{tikzpicture}
\legende{String over a pulley with moment of inertia $I$: when the string suddenly becomes taut, impulsive tensions $J_1$, $J_2$ act on the two portions and an impulsive reaction $R$ acts at the axle.}
\end{popfigure}

\begin{exoresolu}[Particles connected over a pulley with inertia]
A light inextensible string passes over a pulley of radius $a = 0.2\ \mathrm{m}$ and moment of inertia $I = 0.04\ \mathrm{kg\,m^{2}}$, free to turn about a fixed horizontal axis through its centre $O$. Particles of masses $m_1 = 2\ \mathrm{kg}$ and $m_2 = 1\ \mathrm{kg}$ hang from the ends. The system is held with the string just slack and $m_1$ is projected \emph{downwards} with speed $V = 3\ \mathrm{m\,s^{-1}}$; the string suddenly becomes taut. Find the common speed $u$ just after the jerk, the impulsive tensions, and the impulsive reaction at the axle.
\tcblower
\textbf{Solution.} Let $J_1$, $J_2$ be the impulsive tensions on the $m_1$ and $m_2$ sides (upward impulses on the particles), and $\omega = u/a$ the angular speed of the pulley just after the jerk.

\emph{Particle $m_1$} (moving down, impulse up): taking downward positive,
\[
-\,J_1 = m_1(u - V) = 2(u - 3). \tag{1}
\]
\emph{Particle $m_2$} (moving up, impulse up): taking upward positive,
\[
J_2 = m_2(u - 0) = u. \tag{2}
\]
\emph{Pulley} (impulsive couple $= \Delta(I\omega)$ about $O$): the tensions act at the rim, distance $a$:
\[
(J_1 - J_2)\,a = I\omega = \frac{Iu}{a}
\quad\Longrightarrow\quad
J_1 - J_2 = \frac{Iu}{a^{2}} = \frac{0.04}{0.04}\,u = u. \tag{3}
\]
Substituting (1) and (2) into (3):
\[
2(3 - u) - u = u \quad\Longrightarrow\quad 6 = 4u \quad\Longrightarrow\quad u = 1.5\ \mathrm{m\,s^{-1}}.
\]
Then $J_2 = 1.5\ \mathrm{N\,s}$ and $J_1 = 2(3 - 1.5) = 3\ \mathrm{N\,s}$.

\emph{Impulsive reaction at the axle.} The pulley itself acquires no linear momentum, so the upward impulsive reaction $R$ balances the two downward impulsive tensions on the rim:
\[
R = J_1 + J_2 = 4.5\ \mathrm{N\,s}.
\]
\emph{Check (momentum view).} Taking moments of momentum about $O$ for the whole system kills $R$ (it acts at $O$): $m_1 V a = (m_1 + m_2)ua + I\omega$, i.e.\ $2(3)(0.2) = 3(1.5)(0.2) + 0.04(7.5)$, giving $1.2 = 0.9 + 0.3$. \checkmark\ Note also that kinetic energy is lost in the jerk: before, $\tfrac12(2)(9) = 9\ \mathrm{J}$; after, $\tfrac12(3)(1.5)^2 + \tfrac12(0.04)(7.5)^2 = 3.375 + 1.125 = 4.5\ \mathrm{J}$.
\end{exoresolu}

\begin{methode}[Quick route: conservation of moment of momentum about the axle]
When a system is set in motion by a jerk and the only external impulsive forces act \emph{at the axle} (the impulsive reaction) or parallel to it, take moments of momentum about the axle $O$ for the \emph{whole system}: the unknown reaction disappears, and
\[
\sum m v r_{\perp} + I\omega \quad\text{is conserved across the jerk}.
\]
This gives the common speed in one line; the individual impulsive tensions are then recovered from $\Delta(mv)$ on each particle.
\end{methode}

\courssub{Linear Momentum, Angular Momentum, Energy: Which Survives the Impact?}

A recurring examination trap is to reach for the wrong conservation law. The table below summarises what each law requires and when it applies.

\begin{center}
\begin{tabularx}{\linewidth}{|l|X|X|}
\hline
\rowcolor{popblueL} \textbf{Law} & \textbf{Condition} & \textbf{Use across an impact?} \\
\hline
Linear momentum $mv$ & No external \emph{impulse} on the system & Yes, if external impulses vanish (watch the axle reaction!) \\
\hline
Moment of momentum $I\omega$ & No external \emph{impulsive couple} about the chosen axis & Yes — the standard tool for jerks and collisions \\
\hline
Kinetic energy & No inelastic deformation, no impulsive constraints & \textbf{No} — energy is generally lost; use only before/after \\
\hline
\end{tabularx}
\end{center}

\begin{attention}[The axle reaction destroys linear momentum]
When a body is pivoted at a fixed axis, the impulsive reaction at the axle is an \emph{external impulse} on the system, so linear momentum of the system is generally \emph{not} conserved — but its moment of momentum \emph{about the axle} is, because the reaction has zero moment about $O$. This is precisely why moment of momentum, and not linear momentum, is the right tool for pivoted bodies.
\end{attention}

\begin{exercice}[Sliding bead on a rotating rod]
A smooth uniform rod $AB$ of mass $3m$ and length $2a$ rotates freely in a horizontal plane about a fixed vertical axis through its centre $O$ with angular velocity $\Omega$. A bead of mass $2m$ is initially held at $O$ on the rod, then released and slides outwards. Show that when the bead reaches the end $A$, the angular velocity is $\dfrac{\Omega}{3}$, and explain why the kinetic energy has changed.
\end{exercice}

\begin{corrige}[Exercise 1]
The moment of inertia of a uniform rod of mass $M$ and length $2a$ about a perpendicular axis through its centre is $\tfrac13 Ma^{2}$; with $M = 3m$ this gives $I_{\text{rod}} = \tfrac13(3m)a^{2} = ma^{2}$.

\emph{Conservation of moment of momentum.} The rod is smooth, so the forces between bead and rod are internal to the system; the weights are vertical and have zero moment about the vertical axis; the reaction at the axle acts at $O$. Hence the external torque about the vertical axis is zero and moment of momentum is conserved. Initially the bead is at $O$, contributing nothing:
\[
I_1\omega_1 = \big(ma^{2} + 2m(0)^{2}\big)\Omega = ma^{2}\Omega .
\]
When the bead reaches $A$, at distance $a$ from the axis,
\[
I_2 = ma^{2} + 2m\,a^{2} = 3ma^{2}.
\]
Therefore
\[
ma^{2}\Omega = 3ma^{2}\,\omega_2
\quad\Longrightarrow\quad
\omega_2 = \frac{\Omega}{3},
\]
as required.

\emph{Energy.} The kinetic energy changes from
\[
T_1 = \tfrac12 ma^{2}\Omega^{2}
\quad\text{to}\quad
T_2 = \tfrac12 (3ma^{2})\left(\frac{\Omega}{3}\right)^{2} = \tfrac16 ma^{2}\Omega^{2}.
\]
Two-thirds of the kinetic energy is dissipated through the internal constraint forces as the bead is dragged into circular motion at $A$ (the bead does work against these forces while sliding outwards). Moment of momentum is conserved because no external torque acts; kinetic energy is not, because the internal forces do net work during the radial sliding.
\end{corrige}

\begin{exercice}[A jerked string over a rough-rimmed pulley]
A light inextensible string passes over a pulley of radius $0.1\ \mathrm{m}$ and moment of inertia $0.02\ \mathrm{kg\,m^{2}}$ mounted on a fixed horizontal axis. A particle $A$ of mass $3\ \mathrm{kg}$ hangs at one end; a particle $B$ of mass $1\ \mathrm{kg}$ hangs at the other. With the string slack, $B$ is projected \emph{upwards} with speed $4\ \mathrm{m\,s^{-1}}$ and the string suddenly becomes taut. Find the common speed just after the jerk and the impulsive tensions in the two portions of the string.
\end{exercice}

\begin{corrige}[Exercise 2]
Let $u$ be the common speed just after the jerk ($B$ moving up, $A$ moving down), $\omega = u/a = 10u$ the angular speed of the pulley, and $J_A$, $J_B$ the impulsive tensions on the two sides.

\emph{Particle $B$} (upward positive): $J_B = 1(u - 4) = u - 4$ \dots\ with sign care: the impulse is upward, the velocity changes from $+4$ to $+u$, so $J_B = u - 4$, which will come out negative if $u < 4$ — the magnitude is $4 - u$. Write instead, taking the impulse as $J_B$ upward:
\[
J_B = 1\cdot(u - 4). \tag{1}
\]
\emph{Particle $A$} (downward positive, impulse upward): $-J_A = 3(u - 0)$, i.e.
\[
J_A = -3u \quad\text{(impulse downward-positive)}, \tag{2}
\]
so the upward impulsive tension on $A$ has magnitude $3u$.

\emph{Pulley} about the axle: the tension on $B$'s side drives the pulley one way, that on $A$'s side the other:
\[
\big((4 - u) - 3u\big)(0.1) = I\omega = 0.02 \times 10u
\quad\Longrightarrow\quad
0.1(4 - 4u) = 0.2u,
\]
\[
0.4 - 0.4u = 0.2u \quad\Longrightarrow\quad u = \tfrac{2}{3}\ \mathrm{m\,s^{-1}}.
\]
Hence the impulsive tensions have magnitudes $J_B = 4 - \tfrac23 = \tfrac{10}{3}\ \mathrm{N\,s}$ on $B$'s side and $J_A = 3u = 2\ \mathrm{N\,s}$ on $A$'s side. \emph{Check} by moment of momentum about the axle for the whole system: before, $1 \times 4 \times 0.1 = 0.4$; after, $(3+1)\times\tfrac23\times 0.1 + 0.02\times\tfrac{20}{3} = 0.2667 + 0.1333 = 0.4$. \checkmark
\end{corrige}

\begin{aretenir}[Key points of the section]
\begin{itemize}
  \item Moment of momentum of a particle about an axis: $mvr_{\perp}$, with $r_{\perp}$ the \emph{perpendicular} distance to the line of $\vec{v}$.
  \item Rigid body about a fixed axis: $L = I\omega$, and $G = \dfrac{dL}{dt}$.
  \item Zero external torque about the axis $\Rightarrow$ $I\omega$ constant: $I_1\omega_1 = I_2\omega_2$.
  \item Impulsive couple $= \Delta(I\omega)$; impulse $= \Delta(mv)$ for each particle; constraints $v = a\omega$ link them.
  \item In every impact or jerk, kinetic energy is generally lost — use momentum methods across the impact, energy methods only before and after.
  \item Always state the axis about which moments of momentum are taken.
\end{itemize}
\end{aretenir}

\coursec{Energy of a Rotating Body and the Force on the Axis}

In the previous section we saw that the \cle{moment of momentum} of a rigid body about a fixed axis, $L = I\omega$, plays exactly the role that linear momentum plays for a particle: it is changed by impulses and conserved when no external couple acts. We now complete the parallel by introducing the rotational analogue of kinetic energy. Energy methods are extraordinarily powerful in this chapter: they give us the angular speed of a swinging body \emph{without} solving the equation of motion, and — combined with the motion of the centre of mass — they allow us to compute the force that the body exerts on its axis of rotation, a question of obvious practical importance (think of the axle of a grinding wheel in a Douala workshop, or the pivot of a swing in a school playground).

\courssub{Kinetic Energy of a Body Rotating about a Fixed Axis}

\textbf{Intuition.} When a rigid body spins about a fixed axis, every particle of the body is moving, so every particle carries kinetic energy. Particles far from the axis move faster than particles near it. The total kinetic energy must therefore depend on how the mass is distributed relative to the axis — precisely what the moment of inertia measures.

Consider a rigid body rotating with angular speed $\omega$ about a fixed axis. A particle of the body, of mass $m_i$, at perpendicular distance $r_i$ from the axis, describes a circle of radius $r_i$ with speed $v_i = r_i\omega$. Its kinetic energy is
\[
\frac{1}{2}m_i v_i^2 = \frac{1}{2}m_i r_i^2\omega^2.
\]
Summing over all particles of the body:
\[
E_k = \sum_i \frac{1}{2}m_i r_i^2\omega^2 = \frac{1}{2}\omega^2 \sum_i m_i r_i^2 = \frac{1}{2}I\omega^2,
\]
since $\sum_i m_i r_i^2 = I$ is the moment of inertia about the axis.

\begin{propriete}[Kinetic energy of rotation]
A rigid body of moment of inertia $I$ rotating with angular speed $\omega$ about a \textbf{fixed axis} has kinetic energy
\[
\boxed{E_k = \tfrac{1}{2}I\omega^2}
\]
This is the rotational analogue of $\tfrac{1}{2}mv^2$, with $I$ replacing $m$ and $\omega$ replacing $v$. The derivation above is examinable.
\end{propriete}

\begin{exemplebox}[Kinetic energy of a flywheel]
A flywheel of moment of inertia $I = 0.4\ \mathrm{kg\,m^2}$ rotates at $300$ revolutions per minute. Its angular speed is $\omega = \dfrac{300\times 2\pi}{60} = 10\pi\ \mathrm{rad\,s^{-1}}$, so
\[
E_k = \tfrac{1}{2}(0.4)(10\pi)^2 = 20\pi^2 \approx 197.4\ \mathrm{J}.
\]
\end{exemplebox}

\courssub{Work Done by a Couple and the Work--Energy Principle}

\textbf{Intuition.} A couple of moment $C$ turning a body through an angle does work, just as a force moving a particle through a distance does work. Angle replaces distance.

Let a couple of constant moment $C$ act on a body free to rotate about a fixed axis. The couple consists of two equal and opposite forces; when the body turns through a small angle $\delta\theta$, the points of application move along their circles and the total work done by the two forces is $C\,\delta\theta$ (the forces perpendicular to the motion do no work). Hence, for a rotation through a finite angle $\theta$:

\begin{propriete}[Work of a couple; work--energy theorem for rotation]
The work done by a couple of constant moment $C$ turning a body through an angle $\theta$ (in radians) is
\[
\boxed{W = C\theta}
\]
and the work--energy principle for rotation about a fixed axis states:
\[
\text{work done by all couples} = \Delta E_k = \tfrac{1}{2}I\omega_2^2 - \tfrac{1}{2}I\omega_1^2.
\]
\end{propriete}

\begin{attention}[Angles in radians]
The formula $W = C\theta$ requires $\theta$ in \textbf{radians}. A couple of moment $5\ \mathrm{N\,m}$ turning a shaft through one complete revolution does work $W = 5\times 2\pi = 10\pi \approx 31.4\ \mathrm{J}$, not $5\times 360$ joules!
\end{attention}

\courssub{Conservation of Mechanical Energy for a Body Swinging about a Fixed Horizontal Axis}

Consider a rigid body free to rotate about a fixed \emph{horizontal} axis through $O$, with centre of mass $G$ at distance $a$ from $O$. The forces on the body are its weight $Mg$ (acting at $G$) and the reaction at the axis. The reaction acts \emph{at the axis}, which does not move, so it does \textbf{no work}. Weight is conservative. Therefore mechanical energy is conserved:
\[
\tfrac{1}{2}I\omega^2 + Mgh = \text{constant},
\]
where $h$ is the \textbf{height of $G$} above a fixed reference level.

\begin{attention}[The height is the height of $G$]
In the potential energy $Mgh$, $h$ is the height of the \cle{centre of mass} $G$ — not the height of the top of the body, nor of its lowest point. For a uniform rod of length $2a$ pivoted at one end, $G$ is the midpoint, at distance $a$ from the pivot. Using the wrong point is the most common error in this topic.
\end{attention}

\begin{popfigure}
\begin{center}
\begin{tikzpicture}[scale=1.05, >=stealth]
% Pivot
\fill[popdark] (0,0) circle (0.07);
\node[above left, popdark] at (0,0) {$O$};
% Raised rod (150 degrees)
\draw[line width=2.2pt, popblue] (0,0) -- (-1.386,0.8);
\fill[popblue] (-1.386,0.8) circle (0.09);
\node[popblue, above] at (-1.386,0.85) {$G$};
\node[popblue, left] at (-2.77,1.6) {raised position};
\draw[line width=2.2pt, popblue] (-1.386,0.8) -- (-2.77,1.6);
% Lowest rod
\draw[line width=2.2pt, poporange] (0,0) -- (0,-3.2);
\fill[poporange] (0,-1.6) circle (0.09);
\node[poporange, right] at (0.05,-1.6) {$G'$};
\node[poporange, below] at (0,-3.25) {lowest position};
% Circular path of G
\draw[dashed, popmuted] (-1.386,0.8) arc (150:270:1.6);
% Height drop of G
\draw[dotted, popgreen] (-1.386,0.8) -- (2.2,0.8);
\draw[dotted, popgreen] (0,-1.6) -- (2.2,-1.6);
\draw[<->, popgreen, thick] (2.2,0.8) -- (2.2,-1.6) node[midway, right, popgreen] {drop of $G$};
% angle swept
\draw[->, poppurple] (-2.25,1.3) arc (150:270:2.6);
\node[poppurple] at (-2.3,-1.5) {$\theta$};
\end{tikzpicture}
\end{center}
\legende{A rod pivoted at $O$ swings from a raised position to the lowest position. The loss of potential energy is $Mg\times(\text{drop of }G)$, converted into $\tfrac12 I\omega^2$.}
\end{popfigure}

\begin{methode}[The two-step strategy: energy, then centre-of-mass motion]
To analyse a body swinging about a fixed horizontal axis:
\begin{enumerate}
\item \textbf{Energy} gives the speed: apply $\tfrac12 I\omega^2 + Mgh = \text{constant}$ between the release position and the position of interest to find $\omega(\theta)$.
\item \textbf{Angular acceleration}: use $C = I\alpha$ at the position of interest, \emph{or} differentiate the energy equation with respect to time.
\item \textbf{Reaction at the axis}: write $M\vec{a}_G = \vec{R} + M\vec{g}$ and resolve along and perpendicular to $OG$, using the radial and transverse accelerations of $G$.
\end{enumerate}
\end{methode}

\courssub{Motion of the Centre of Mass; Acceleration Components of $G$}

Since the body is rigid and the axis fixed, $G$ moves on a \textbf{circle} of radius $a = OG$ centred at $O$. If $\theta$ is the angle that $OG$ makes with the downward vertical, then the acceleration of $G$ has the two standard components of circular motion:
\begin{itemize}
\item a \textbf{radial component} (along $GO$, towards $O$): $a\omega^2 = a\dot{\theta}^2$;
\item a \textbf{transverse component} (perpendicular to $OG$, in the direction of increasing $\theta$): $a\alpha = a\ddot{\theta}$.
\end{itemize}

\begin{popfigure}
\begin{center}
\begin{tikzpicture}[scale=1.15, >=stealth]
\fill[popdark] (0,0) circle (0.07);
\node[above, popdark] at (0,0) {$O$};
% rod at angle -35 degrees
\draw[line width=2.2pt, popblue] (0,0) -- (2.785,-1.95);
\fill[popblue] (1.393,-0.975) circle (0.1);
\node[popblue, below left] at (1.35,-1.0) {$G$};
% dashed circle of radius 1.7
\draw[dashed, popmuted] (0,0) circle (1.7);
% downward vertical
\draw[dotted, popmuted] (0,0) -- (0,-2.1);
\draw[->, poppurple] (0,-1.3) arc (-90:-35:1.3);
\node[poppurple] at (0.62,-1.25) {$\theta$};
% radial acceleration (towards O)
\draw[->, very thick, poporange] (1.393,-0.975) -- (0.53,-0.37) node[midway, above right, poporange] {$a\omega^2$};
% transverse acceleration (perpendicular to OG, direction of increasing theta)
\draw[->, very thick, popgreen] (1.393,-0.975) -- (2.17,0.13) node[midway, below right, popgreen] {$a\alpha$};
% weight
\draw[->, very thick, poppink] (1.393,-0.975) -- (1.393,-2.475) node[below, poppink] {$Mg$};
\end{tikzpicture}
\end{center}
\legende{$G$ moves on a circle about $O$. Its acceleration has a radial component $a\omega^2$ (towards $O$) and a transverse component $a\alpha$ (perpendicular to $OG$).}
\end{popfigure}

\courssub{The Force Exerted on the Axis}

The pivot exerts a reaction $\vec{R}$ on the body. The centre-of-mass theorem (Newton's second law applied to the whole body) gives
\[
M\vec{a}_G = \vec{R} + M\vec{g}.
\]
We resolve this vector equation \textbf{along $GO$} (radially, towards $O$) and \textbf{perpendicular to $OG$} (transversally). With $\theta$ measured from the downward vertical, the weight has radial component $-Mg\cos\theta$ (i.e. $Mg\cos\theta$ directed away from $O$ when $G$ is below $O$) and transverse component $-Mg\sin\theta$. Writing $R_r$ and $R_t$ for the components of $\vec{R}$ along $GO$ (towards $O$) and transversally (in the direction of increasing $\theta$):
\[
R_r - Mg\cos\theta = Ma\omega^2, \qquad R_t - Mg\sin\theta = Ma\alpha.
\]
Hence
\[
\boxed{R_r = Mg\cos\theta + Ma\omega^2, \qquad R_t = Mg\sin\theta + Ma\alpha.}
\]

\begin{attention}[The reaction is not vertical]
The reaction at the axis has \textbf{two components}, radial and transverse. Do not assume $\vec{R}$ is vertical, and do not assume it balances the weight: the pivot must also supply the centripetal force $Ma\omega^2$ that keeps $G$ on its circle. At the lowest point of a fast swing, the reaction can be several times the weight.
\end{attention}

\courssub{Complete Worked Problem: Rod Released from the Horizontal}

\begin{exoresolu}[Rod falling from the horizontal position]
A uniform rod $AB$ of length $2a$ and mass $M$ is freely pivoted at its end $A$ about a fixed horizontal axis. It is released from rest in the horizontal position. Find (a) its angular speed when it first reaches the vertical position; (b) the magnitude and direction of the reaction at the pivot in that position. \;[Take $I_A = \tfrac{4}{3}Ma^2$.]
\tcblower
\textbf{(a) Angular speed by energy conservation.} $G$ is the midpoint, with $AG = a$. Between the horizontal and the vertical positions, $G$ falls through a height $a$. Energy conservation:
\[
0 + Mga = \tfrac{1}{2}I_A\omega^2 + 0
\quad\Longrightarrow\quad
Mga = \tfrac{1}{2}\cdot\tfrac{4}{3}Ma^2\,\omega^2 = \tfrac{2}{3}Ma^2\omega^2,
\]
\[
\boxed{\omega^2 = \frac{3g}{2a}}.
\]

\textbf{Angular acceleration at the vertical position.} When the rod hangs vertically, the weight acts along the line $AG$, so its moment about $A$ is zero: $C = I_A\alpha = 0$, hence $\alpha = 0$ at that instant.

\textbf{(b) Reaction at the pivot.} At the vertical position ($\theta = 0$), resolve $M\vec{a}_G = \vec{R} + M\vec{g}$.

\emph{Transverse direction:} $R_t - 0 = Ma\alpha = 0$, so $R_t = 0$: the reaction is purely vertical here.

\emph{Radial direction (upwards, towards $A$):}
\[
R_r - Mg = Ma\omega^2 = Ma\cdot\frac{3g}{2a} = \frac{3}{2}Mg
\quad\Longrightarrow\quad
R_r = Mg + \frac{3}{2}Mg = \frac{5}{2}Mg.
\]
\[
\boxed{R = \tfrac{5}{2}Mg \text{, directed vertically upwards.}}
\]
The pivot therefore supports two and a half times the weight of the rod at the bottom of the swing — the weight $Mg$ plus the centripetal force $\tfrac32 Mg$.
\end{exoresolu}

\begin{popfigure}
\begin{center}
\begin{tikzpicture}
\begin{axis}[
    width=11cm, height=6.2cm,
    xlabel={$\theta$ (degrees from downward vertical)},
    ylabel={$R_r/Mg$},
    xmin=0, xmax=90, ymin=0, ymax=2.8,
    xtick={0,15,30,45,60,75,90},
    ytick={0,0.5,1,1.5,2,2.5},
    grid=both, grid style={popline, very thin},
    axis line style={popdark},
    tick label style={font=\small},
    label style={popdark},
]
\addplot[domain=0:90, samples=120, very thick, popblue] {2.5*cos(x)};
\node[popblue, font=\small] at (axis cs:38,1.75) {$R_r = \tfrac52 Mg\cos\theta$};
\addplot[domain=0:90, samples=2, dashed, poporange] {1};
\node[poporange, font=\small] at (axis cs:20,1.12) {$R = Mg$};
\end{axis}
\end{tikzpicture}
\end{center}
\legende{Radial component of the reaction at the axis for the rod released from the horizontal, as a function of the angle $\theta$ from the downward vertical. Energy gives $a\omega^2 = \tfrac32 g\cos\theta$, so $R_r/Mg = \tfrac52\cos\theta$, rising from $0$ at release (the rod is momentarily at rest, and the weight has no radial component) to $\tfrac52$ at the lowest point.}
\end{popfigure}

\begin{savaistu}[Why engineers care about the force on the axis]
The factor $\tfrac52$ above is not a curiosity: any rotating machine — a ceiling fan, a water pump, the turbines of the Song Loulou hydroelectric dam on the Sanaga river — exerts on its bearings forces that exceed the static weight whenever it moves. Bearings and axles must be designed for these \emph{dynamic} loads, computed exactly as in this section: energy for the speed, centre-of-mass motion for the reaction.
\end{savaistu}

\begin{exercice}[Swinging rod released at $60^{\circ}$]
A uniform rod of mass $M$ and length $2a$ is freely pivoted at one end about a fixed horizontal axis. It is released from rest with the rod making an angle of $60^{\circ}$ with the downward vertical (below the pivot). Find (a) the angular speed when the rod passes through the vertical; (b) the reaction at the pivot at that instant.
\end{exercice}

\begin{corrige}[Exercise 1]
(a) $G$ is at distance $a$ from the pivot. Between the release position ($\theta = 60^{\circ}$) and the vertical ($\theta = 0$), $G$ drops through a height
\[
a\cos 0^{\circ} - a\cos 60^{\circ} = a\left(1 - \tfrac12\right) = \tfrac{a}{2}.
\]
Energy conservation:
\[
Mg\cdot\tfrac{a}{2} = \tfrac12\cdot\tfrac43 Ma^2\omega^2
\;\Longrightarrow\;
\omega^2 = \frac{3g}{4a}.
\]
(b) At the vertical position the moment of the weight about the pivot is zero, so $\alpha = 0$ and the transverse equation gives $R_t = 0$. Radially (upwards, towards the pivot):
\[
R - Mg = Ma\omega^2 = Ma\cdot\frac{3g}{4a} = \tfrac34 Mg
\;\Longrightarrow\;
\boxed{R = \tfrac74 Mg \text{, vertically upwards.}}
\]
\end{corrige}

\begin{aretenir}[Key points of this section]
\begin{itemize}
\item Kinetic energy of rotation about a fixed axis: $E_k = \tfrac12 I\omega^2$; work of a couple: $W = C\theta$ ($\theta$ in radians).
\item For a body swinging about a fixed horizontal axis, $\tfrac12 I\omega^2 + Mgh = \text{constant}$, with $h$ the height of $G$.
\item $G$ moves on a circle: radial acceleration $a\omega^2$ towards $O$, transverse acceleration $a\alpha$.
\item The reaction at the axis comes from $M\vec{a}_G = \vec{R} + M\vec{g}$ resolved radially and transversally: it has two components and generally exceeds the weight.
\item Strategy: energy $\rightarrow \omega$; $C = I\alpha$ $\rightarrow \alpha$; centre-of-mass equation $\rightarrow \vec{R}$.
\end{itemize}
\end{aretenir}

\coursec{The Compound Pendulum and Synthesis}

In the previous section we studied the energy of a rotating body and the force exerted on its axis. We saw that a rigid body pivoted at a fixed point can swing under gravity, and that energy methods tell us its speed at any position. We now close this circle: when a rigid body pivoted about a fixed horizontal axis is displaced slightly from its equilibrium position and released, it \emph{oscillates}. Such a body is called a \cle{compound pendulum}, and it is the most important application of everything in this chapter: the equation of motion $I_O\ddot{\theta}=C$, the parallel axes theorem, energy conservation, and even impulse all appear in its study. It is also a classic GCE practical topic, because it provides one of the most accurate school-laboratory methods for measuring $g$.

\courssub{Definition and Setting Up the Equation of Motion}

\begin{definition}[Compound pendulum]
A \cle{compound pendulum} (or \emph{physical pendulum}) is a rigid body of mass $M$, free to rotate about a fixed \textbf{horizontal} axis through a point $O$, oscillating under gravity alone. Let $G$ be the centre of mass of the body and $h = OG$ the distance from the axis to $G$.
\end{definition}

A simple pendulum (a point mass on a light string) is an idealisation; every real pendulum — a swinging metre rule in the laboratory, a bell, a swinging door, the leg of a walking person — is a compound pendulum, because its mass is distributed.

Consider the body displaced through an angle $\theta$ from the vertical equilibrium position. Two forces act on it: the weight $Mg$ acting vertically downwards at $G$, and the reaction $\vec{R}$ of the axis at $O$. The reaction has zero moment about $O$ (it acts at the axis), so the only torque about $O$ is that of the weight. The line of action of $Mg$ passes at a perpendicular distance $h\sin\theta$ from $O$, and it acts so as to \emph{restore} the body towards equilibrium. Hence the equation of motion $I_O\ddot{\theta} = C_O$ becomes
\[ I_O\ddot{\theta} = -Mgh\sin\theta, \]
where $I_O$ is the moment of inertia of the body about the axis through $O$.

\begin{popfigure}
\begin{center}
\begin{tikzpicture}[scale=1.1]
\fill[poppurple!60!black] (0,0) circle (0.06);
\draw[popline, thick] (-1.4,0) -- (1.4,0);
\foreach \x in {-1.2,-0.9,-0.6,-0.3,0,0.3,0.6,0.9,1.2}{\draw[popline] (\x,0) -- (\x+0.15,0.18);}
\draw[dashed, popmuted] (0,0) -- (0,-3.2);
\draw[very thick, popblue, rounded corners=2pt] (0.42,-0.35) -- (1.05,-2.6) arc(-68:-112:2.82) -- (-0.42,-0.35) arc(200:340:0.44) -- cycle;
\fill[poporange] (0,-1.6) circle (0.09) node[right=2pt, popdark]{\textbf{G}};
\draw[<->, popdark] (0.12,-0.25) -- (0.12,-1.5) node[midway, right]{$h$};
\draw[->, popdark, thick] (0,-1.6) -- (0,-2.9) node[below]{$Mg$};
\draw[->, popgreen, thick] (0,0) -- (0.75,0.55) node[right]{$\vec{R}$};
\draw[->, poporange, thick] (0,-1.1) arc(-90:-68:1.1) node[midway, below right]{$\theta$};
\node[popdark] at (0.28,0.28){$O$};
\end{tikzpicture}
\end{center}
\legende{A compound pendulum: rigid body pivoted at $O$, centre of mass $G$ with $OG=h$, displaced by $\theta$. The weight $Mg$ provides the restoring torque $-Mgh\sin\theta$; the reaction $\vec{R}$ at the axis has no moment about $O$.}
\end{popfigure}

\courssub{Small Oscillations: Simple Harmonic Motion}

The equation $I_O\ddot{\theta} = -Mgh\sin\theta$ is \emph{not} simple harmonic in general, because of the $\sin\theta$. But for small oscillations we may linearise it.

\begin{attention}[The small-angle approximation]
The approximation $\sin\theta \approx \theta$ is valid \textbf{only when $\theta$ is measured in radians} and is small (in practice $|\theta| \lesssim 10^{\circ}$, i.e. $0.17$ rad, gives an error below $0.5\%$). You must \textbf{state this approximation explicitly} in every derivation — examiners award a mark for it. Never write $\sin\theta \approx \theta$ with $\theta$ in degrees.
\end{attention}

For small oscillations, $\sin\theta \approx \theta$, so
\[ I_O\ddot{\theta} = -Mgh\theta \quad\Longleftrightarrow\quad \ddot{\theta} = -\frac{Mgh}{I_O}\,\theta. \]
This is the SHM equation $\ddot{\theta} = -\omega^2\theta$ with $\omega^2 = \dfrac{Mgh}{I_O}$.

\begin{propriete}[Period of a compound pendulum]
A compound pendulum performing small oscillations about a fixed horizontal axis through $O$ executes simple harmonic motion with period
\[ T = 2\pi\sqrt{\frac{I_O}{Mgh}}, \]
where $I_O$ is the moment of inertia about the axis, $M$ the mass, and $h = OG$. (The derivation above — restoring torque, small-angle approximation, identification with SHM — is examinable.)
\end{propriete}

\begin{definition}[Length of the equivalent simple pendulum]
The \cle{length of the equivalent simple pendulum} is the length $l$ of a simple pendulum having the same period as the compound pendulum:
\[ 2\pi\sqrt{\frac{l}{g}} = 2\pi\sqrt{\frac{I_O}{Mgh}} \quad\Longrightarrow\quad \boxed{\,l = \frac{I_O}{Mh} = \frac{k_O^{\,2}}{h}\,} \]
where $k_O$ is the radius of gyration about the axis through $O$ (since $I_O = Mk_O^{\,2}$).
\end{definition}

Using the parallel axes theorem, $k_O^{\,2} = k_G^{\,2} + h^2$, where $k_G$ is the radius of gyration about the parallel axis through $G$. Hence
\[ l = \frac{k_G^{\,2} + h^2}{h} = h + \frac{k_G^{\,2}}{h}. \]

\begin{savaistu}[Centres of suspension and oscillation (enrichment)]
Mark the point $O'$ on the line $OG$ produced, at distance $l$ from $O$. Then $OO' = l = h + k_G^2/h$, so $GO' = k_G^2/h$. A beautiful result (Kater's theorem) states that if the pendulum is inverted and suspended from $O'$, its period is \emph{unchanged}: the centres of \textbf{suspension} and \textbf{oscillation} are interchangeable. This is the principle of Kater's reversible pendulum, historically the most accurate instrument for measuring $g$. You will not be asked to prove it, but recognising it can earn credit in structured questions.
\end{savaistu}

\courssub{Worked Method and Examples}

\begin{methode}[Solving compound pendulum problems]
\begin{enumerate}
\item \textbf{Locate $G$} and compute $h = OG$.
\item \textbf{Compute $I_O$}: use the standard result about $G$ (or about an end) and the \textbf{parallel axes theorem} $I_O = I_G + Mh^2$ if necessary.
\item \textbf{Apply the period formula} $T = 2\pi\sqrt{I_O/(Mgh)}$, or the equivalent length $l = I_O/(Mh)$.
\item For experimental questions, \textbf{rearrange to make $g$ the subject}: $g = \dfrac{4\pi^2 I_O}{MhT^2}$.
\item Check units and state the small-angle assumption where required.
\end{enumerate}
\end{methode}

\begin{exemplebox}[Uniform rod swinging about one end]
A uniform rod $AB$ of length $2a$ and mass $M$ is pivoted at $A$ and makes small oscillations in a vertical plane. Find the period and the length of the equivalent simple pendulum.

\textbf{Solution.} Here $h = AG = a$ and $I_A = \dfrac{4}{3}Ma^2$ (standard result). Hence
\[ T = 2\pi\sqrt{\frac{\tfrac{4}{3}Ma^2}{Mga}} = 2\pi\sqrt{\frac{4a}{3g}}, \qquad l = \frac{I_A}{Ma} = \frac{4a}{3}. \]
The rod of length $2a$ swings like a simple pendulum of length $\tfrac{4a}{3}$ — shorter than the rod, so the rod swings \emph{faster} than a simple pendulum of its own length.
\end{exemplebox}

\begin{exoresolu}[Bar pendulum used to find $g$]
A uniform metre rule of mass $M$ is pivoted about a horizontal axis through a hole drilled $10\ \mathrm{cm}$ from one end. It performs small oscillations with a measured period $T = 1.53\ \mathrm{s}$. Taking $g = 9.81\ \mathrm{m\,s^{-2}}$ first verify the theory, then show how the experiment yields $g$.
\tcblower
\textbf{Solution.} The rule has length $L = 1.00\ \mathrm{m}$, so $G$ is at its midpoint and
\[ h = OG = 0.50 - 0.10 = 0.40\ \mathrm{m}. \]
About its centre, $I_G = \dfrac{1}{12}ML^2$; by the parallel axes theorem,
\[ I_O = M\left(\frac{L^2}{12} + h^2\right) = M\left(\frac{1}{12} + 0.16\right) = 0.2433\,M\ \mathrm{m^2}. \]
The predicted period is
\[ T = 2\pi\sqrt{\frac{I_O}{Mgh}} = 2\pi\sqrt{\frac{0.2433}{9.81 \times 0.40}} = 2\pi\sqrt{0.0620} \approx 1.56\ \mathrm{s}, \]
in good agreement with the measured $1.53\ \mathrm{s}$ (the small difference comes from timing and pivot friction). Conversely, from the measurement,
\[ g = \frac{4\pi^2 I_O}{MhT^2} = \frac{4\pi^2 \times 0.2433}{0.40 \times 1.53^2} \approx 10.3\ \mathrm{m\,s^{-2}}, \]
which is the experimental estimate of $g$ (a more careful timing over $50$ oscillations reduces the error). Note that $M$ cancels: \textbf{the mass of the pendulum is irrelevant}, which is why the bar pendulum is such a clean experiment.
\end{exoresolu}

\courssub{The Bar Pendulum: Period as a Function of $h$}

A bar pendulum is a rectangular bar with a row of holes; it can be suspended from any hole, at distance $h$ from $G$. Since $l = h + \dfrac{k_G^{\,2}}{h}$, the period is
\[ T = 2\pi\sqrt{\frac{h^2 + k_G^{\,2}}{gh}}. \]
As $h \to 0$ (pivot near $G$), $T \to \infty$ — the body barely oscillates because the restoring torque vanishes. As $h \to \infty$, $T \to \infty$ as well. Between the two, $T$ has a \textbf{minimum}. Differentiating $T^2 = \dfrac{4\pi^2}{g}\left(h + \dfrac{k_G^2}{h}\right)$ with respect to $h$:
\[ \frac{d(T^2)}{dh} = \frac{4\pi^2}{g}\left(1 - \frac{k_G^{\,2}}{h^2}\right) = 0 \quad\Longrightarrow\quad h = k_G. \]

\begin{propriete}[Minimum period (GCE-useful enrichment)]
The period of a compound pendulum is minimum when the axis is at distance $h = k_G$ from the centre of mass, and the minimum period is
\[ T_{\min} = 2\pi\sqrt{\frac{2k_G}{g}}. \]
\end{propriete}

\begin{popfigure}
\begin{center}
\begin{tikzpicture}[scale=0.9]
% Bar pendulum with holes
\draw[very thick, popblue, rounded corners=1pt] (-0.35,0) rectangle (0.35,-4);
\foreach \y in {-0.5,-1.0,-1.5,-2.0,-2.5,-3.0,-3.5}{\fill[white] (0,\y) circle (0.09); \draw[popblue] (0,\y) circle (0.09);}
\fill[poppurple] (0,-0.5) circle (0.05); \node[left, popdark] at (-0.4,-0.5){$O$};
\fill[poporange] (0,-2) circle (0.08); \node[right, popdark] at (0.4,-2){$G$};
\draw[<->, popdark] (0.55,-0.5) -- (0.55,-2) node[midway, right]{$h$};
% Equivalent simple pendulum
\begin{scope}[xshift=4.5cm]
\fill[poppurple] (0,0) circle (0.05); \node[left, popdark] at (-0.1,0){$O$};
\draw[popdark, thick] (0,0) -- (0,-2.67);
\fill[poporange] (0,-2.67) circle (0.12);
\draw[<->, popdark] (0.35,0) -- (0.35,-2.67) node[midway, right]{$l = \dfrac{k_O^2}{h}$};
\draw[dashed, popmuted] (0,0) -- (0,-0.6);
\draw[->, poporange] (0,-0.5) arc(-90:-72:0.5) node[midway, below right]{$\theta$};
\end{scope}
\end{tikzpicture}
\end{center}
\legende{Bar pendulum with several suspension holes, and the equivalent simple pendulum of length $l = k_O^2/h$ having the same period.}
\end{popfigure}

\begin{popfigure}
\begin{center}
\begin{tikzpicture}
\begin{axis}[
width=0.72\linewidth, height=6cm,
xlabel={$h$ (m)}, ylabel={$T$ (s)},
xmin=0, xmax=0.6, ymin=1.2, ymax=3,
domain=0.03:0.6, samples=120,
axis lines=left, thick,
xtick={0.289}, xticklabels={$k_G$},
ytick=\empty,
]
\addplot[popblue, very thick] {2*pi*sqrt(max(0,(x^2+0.0833)/(9.81*x)))};
\addplot[dashed, popmuted] coordinates {(0.289,1.2) (0.289,1.552)};
\addplot[dashed, popmuted] coordinates {(0,1.552) (0.289,1.552)};
\addplot[only marks, mark=*, poporange] coordinates {(0.289,1.552)};
\node[popdark, anchor=west] at (axis cs:0.31,1.5){$T_{\min}$};
\end{axis}
\end{tikzpicture}
\end{center}
\legende{Period $T$ against suspension distance $h$ for a metre rule ($k_G = L/\sqrt{12} \approx 0.289\ \mathrm{m}$). The minimum occurs at $h = k_G$. For a given $T > T_{\min}$ there are two values $h_1$ and $h_2$; their sum $h_1+h_2$ equals the equivalent length $l$.}
\end{popfigure}

\begin{experience}[Finding $g$ with a bar pendulum]
Suspend the bar from each hole in turn, on each side of $G$, and time $50$ small oscillations to find $T$ for each $h$. Plot $T$ against $h$ (distance from $G$): the curve is symmetric about $G$. For a chosen $T > T_{\min}$, the horizontal line cuts the curve at two points $h_1$ and $h_2$ on the same side of $G$ (or four points if both sides are plotted). The equivalent length is $l = h_1 + h_2$, and $g = \dfrac{4\pi^2 l}{T^2}$. This graphical method averages out errors in locating $G$ and measuring $h$ — a favourite GCE practical question.
\end{experience}

\courssub{Synthesis Problems: The Whole Chapter at Work}

GCE questions on rotational dynamics frequently combine several lessons in one problem. The following worked exercise is typical of the hardest questions on the paper.

\begin{exoresolu}[Rod falling, hitting a stop, then oscillating]
A uniform rod $AB$ of length $2a = 1.2\ \mathrm{m}$ and mass $M = 2\ \mathrm{kg}$ is smoothly pivoted at $A$. It is held horizontal and released from rest. (a) Find its angular velocity when it first reaches the vertical. (b) As it passes the vertical, a small lump of plasticine of mass $m = 0.5\ \mathrm{kg}$ sticks to the rod at $B$. Find the new angular velocity immediately after the impact. (c) Find the period of small oscillations of the loaded rod about $A$. Take $g = 9.81\ \mathrm{m\,s^{-2}}$.
\tcblower
\textbf{(a) Energy.} Falling from horizontal to vertical, $G$ drops by $a = 0.6\ \mathrm{m}$. With $I_A = \tfrac{4}{3}Ma^2 = \tfrac{4}{3}(2)(0.36) = 0.96\ \mathrm{kg\,m^2}$:
\[ Mga = \tfrac{1}{2}I_A\omega^2 \;\Longrightarrow\; \omega^2 = \frac{2 \times 2 \times 9.81 \times 0.6}{0.96} = 24.525, \quad \omega \approx 4.95\ \mathrm{rad\,s^{-1}}. \]
\textbf{(b) Conservation of moment of momentum about $A$} (the impulse at the pivot has no moment about $A$). Just before impact, $L = I_A\omega = 0.96 \times 4.95 = 4.754\ \mathrm{kg\,m^2\,s^{-1}}$. After impact the inertia is
\[ I_A' = 0.96 + m(2a)^2 = 0.96 + 0.5 \times 1.44 = 1.68\ \mathrm{kg\,m^2}, \]
so $\omega' = \dfrac{4.754}{1.68} \approx 2.83\ \mathrm{rad\,s^{-1}}$. (Kinetic energy is \emph{not} conserved in the impact — check: $\tfrac12 I_A\omega^2 = 11.77$ J before, $\tfrac12 I_A'\omega'^2 = 6.73$ J after.)
\textbf{(c) Period.} The new centre of mass is at distance $h$ from $A$ with
\[ h = \frac{Ma + m(2a)}{M+m} = \frac{2(0.6) + 0.5(1.2)}{2.5} = 0.72\ \mathrm{m}. \]
Then
\[ T = 2\pi\sqrt{\frac{I_A'}{(M+m)gh}} = 2\pi\sqrt{\frac{1.68}{2.5 \times 9.81 \times 0.72}} \approx 1.94\ \mathrm{s}. \]
This single problem used energy (Section 5), moment of momentum and impulse (Section 4), the parallel axes theorem (Sections 1--2) and the compound pendulum (this section).
\end{exoresolu}

\begin{exercice}[Compound pendulum synthesis]
A uniform circular disc of mass $1.5\ \mathrm{kg}$ and radius $0.2\ \mathrm{m}$ is pivoted about a horizontal axis perpendicular to its plane through a point $O$ on its rim.
\begin{enumerate}
\item Show that $I_O = 0.09\ \mathrm{kg\,m^2}$.
\item Find the period of small oscillations in the plane of the disc.
\item Find the length of the equivalent simple pendulum.
\item The disc is released from rest with $G$ level with $O$. Find the angular velocity when $G$ is directly below $O$, and the reaction at the axis at that instant.
\end{enumerate}
\end{exercice}

\begin{corrige}[Exercise 1]
\textbf{1.} $I_G = \tfrac12 Mr^2 = 0.03\ \mathrm{kg\,m^2}$; parallel axes: $I_O = I_G + Mr^2 = 0.03 + 1.5(0.04) = 0.09\ \mathrm{kg\,m^2}$. \quad \textbf{2.} $h = r = 0.2$ m: $T = 2\pi\sqrt{0.09/(1.5 \times 9.81 \times 0.2)} = 2\pi\sqrt{0.03058} \approx 1.10\ \mathrm{s}$. \quad \textbf{3.} $l = I_O/(Mh) = 0.09/0.3 = 0.30\ \mathrm{m}$. \quad \textbf{4.} $G$ drops $0.2$ m: $\tfrac12 I_O\omega^2 = Mgh \Rightarrow \omega^2 = 2(1.5)(9.81)(0.2)/0.09 = 65.4$, $\omega \approx 8.09\ \mathrm{rad\,s^{-1}}$. At the lowest point the reaction $R$ is vertical; the centre of mass moves in a circle of radius $h$, so $R - Mg = Mh\omega^2$, giving $R = 1.5(9.81) + 1.5(0.2)(65.4) = 14.7 + 19.6 = 34.3\ \mathrm{N}$ upwards.
\end{corrige}

\courssub{Revision Map: The Translation--Rotation Analogy}

Nearly every formula of this chapter is the rotational twin of a translational formula you already know. Learning the dictionary below is the fastest way to revise the whole chapter.

\begin{aretenir}[Translation--rotation analogy table]
\begin{center}
\begin{tabularx}{\linewidth}{|l|X|X|}
\hline
\rowcolor{popblueL} \textbf{Quantity / law} & \textbf{Translation} & \textbf{Rotation} \\
\hline
Displacement & $x$ & $\theta$ \\
\hline
Velocity & $v = \dot{x}$ & $\omega = \dot{\theta}$ \\
\hline
Acceleration & $a = \dot{v}$ & $\ddot{\theta}$ \\
\hline
Inertia & mass $m$ & moment of inertia $I = \sum mr^2$ \\
\hline
Force / torque & $F$ & $C = Fr$ \\
\hline
Equation of motion & $F = ma$ & $C = I\ddot{\theta}$ \\
\hline
Momentum & $p = mv$ & $L = I\omega$ \\
\hline
Impulse & $J = \Delta p$ & $J\,r = \Delta L$ \\
\hline
Kinetic energy & $\tfrac12 mv^2$ & $\tfrac12 I\omega^2$ \\
\hline
Work / power & $Fx$, $Fv$ & $C\theta$, $C\omega$ \\
\hline
SHM period & $2\pi\sqrt{m/k}$ & $2\pi\sqrt{I_O/(Mgh)}$ \\
\hline
\end{tabularx}
\end{center}
\medskip
Key chapter results: $I_O = I_G + Mh^2$; $k_O^2 = k_G^2 + h^2$; $L = I\omega$ conserved when no external torque; energy $\tfrac12 I\omega^2 + Mgy_G$ conserved for smooth pivots; compound pendulum $T = 2\pi\sqrt{I_O/(Mgh)}$, equivalent length $l = k_O^2/h$, minimum period at $h = k_G$.
\end{aretenir}

\begin{attention}[Final checklist of classic errors]
\begin{itemize}
\item Forgetting the parallel axes theorem: $I_O \neq I_G$ unless the pivot is at $G$.
\item Mixing $h$ (pivot to $G$) with the full length of the body.
\item Using $\sin\theta \approx \theta$ without stating it, or with $\theta$ in degrees.
\item Conserving kinetic energy through an inelastic impact — only \emph{moment of momentum} about the pivot is conserved there.
\item In $T = 2\pi\sqrt{I_O/(Mgh)}$, using $I_G$ instead of $I_O$ — the single most common error in this topic.
\end{itemize}
\end{attention}

You have now completed the chapter on rotational dynamics: from the definition of moment of inertia, through the equation of motion under a constant torque, moment of momentum, energy, and finally the compound pendulum. Master the analogy table, rehearse the derivations flagged as examinable, and practise synthesis problems like the ones above — that is exactly the profile of the GCE Advanced Level questions on this chapter.

\newpage

\coursec{Integration activity}

\courssub{Aims of the activity}

This closing activity brings together \emph{all} the ideas of the chapter in one authentic engineering study. By the end of it, you should be able to:

\begin{itemize}
\item select and use the standard \cle{moment of inertia} of a uniform disc, $I=\tfrac12 MR^2$, and of a hoop, $I=MR^2$;
\item apply the \cle{work--energy principle} for rotation, $W_C=\Delta E_k$, with $E_k=\tfrac12 I\omega^2$ and the work of a constant couple $W=C\theta$;
\item use the equation of motion of a rigid body under a constant torque, $C=I\alpha$, together with the constant-acceleration kinematic relations $\omega^2=\omega_0^2+2\alpha\theta$ and $\omega=\omega_0+\alpha t$;
\item apply \cle{conservation of moment of momentum} (angular momentum) when a mass is suddenly added to a rotating body, and compute the associated \cle{energy loss};
\item compare two designs (disc versus rim-loaded wheel) and \emph{justify} an engineering choice with numbers.
\end{itemize}

\begin{aretenir}[Formulae you will need]
Uniform disc about its central axis: $I=\dfrac12 MR^2$; hoop (mass at the rim): $I=MR^2$. Rotational kinetic energy: $E_k=\dfrac12 I\omega^2$. Work of a constant couple: $W=C\theta$. Equation of motion: $C=I\alpha$. Angular momentum about a fixed axis: $L=I\omega$, conserved when the external torque about the axis is zero.
\end{aretenir}

\courssub{Situation}

\begin{experience}[The potter's wheel of Bamessing]
In the village of Bamessing (Ndop, North-West Region), pottery is a centuries-old craft. A craftsman, Pa Nformi, uses a traditional kick-wheel: a heavy horizontal wheel mounted on a vertical axle in smooth bearings. The potter gives the wheel one strong kick, then shapes the clay while the wheel spins freely. If the wheel slows down too quickly, the pot becomes lopsided, so the wheel must \emph{keep spinning long enough} after a single kick.

Pa Nformi asks your Further Mathematics class to design the wheel. He gives you the following specifications and measurements:

\begin{itemize}
\item The wheel is modelled as a \textbf{uniform disc} of radius $R=0.30$ m and negligible thickness, rotating about a fixed vertical axis through its centre.
\item One strong kick supplies the wheel with a kinetic energy of exactly $E_0=120$ J.
\item Friction in the bearings and air resistance are modelled by a \textbf{constant resisting couple} $C_r=0.5\ \mathrm{N\,m}$.
\item \textbf{Requirement:} after the kick, the wheel must complete \textbf{at least $N=40$ revolutions} before coming to rest.
\item Later, while the wheel spins freely, a lump of clay of mass $m=2$ kg is dropped vertically onto the wheel \textbf{at the rim} and sticks there.
\end{itemize}

Take $g=9.81\ \mathrm{m\,s^{-2}}$ where needed. The axle is frictionless apart from the resisting couple already given, and its own moment of inertia is neglected.
\end{experience}

\begin{popfigure}
\begin{tikzpicture}[scale=1.2]
% Top view
\draw[thick, fill=popblueL!30] (0,0) circle (1.5);
\draw[thick] (0,0) circle (0.1);
\node at (0,0) {$\times$};
\draw[->, thick, popblue] (1.5,0) arc (0:90:1.5) node[midway, right] {$\omega$};
% Clay lump at rim
\draw[fill=poporange] (1.5,0) circle (0.15);
\node[below] at (1.5,-0.35) {Clay $m$};
% Side view shifted
\begin{scope}[xshift=5.5cm]
\draw[thick] (-1.5,-0.2) rectangle (1.5,0); % base
\draw[thick] (0,0) -- (0,2); % axle
\draw[thick, fill=popblueL!30] (-1.5,2) -- (1.5,2) -- (1.5,2.2) -- (-1.5,2.2) -- cycle; % disc side
\node at (0,2.1) {Disc $M,R$};
\draw[->, thick, popblue] (1.5,2.1) arc (0:90:1.5) node[midway, right] {$\omega$};
\end{scope}
\end{tikzpicture}
\legende{Top view (left) and side view (right) of the potter's wheel. The clay lump is dropped onto the rim.}
\end{popfigure}

\courssub{Tasks}

Work in groups of three or four. Present every numerical answer with its unit and a short sentence of interpretation.

\begin{enumerate}
\item \textbf{Minimum moment of inertia.} Using the work--energy principle between the instant just after the kick and the instant the wheel stops, show that the wheel needs a minimum moment of inertia $I_{\min}$ and compute its value. (Hint: the resisting couple does work $-C_r\theta$ over the total angle $\theta=2\pi N$.)
\item \textbf{Mass of the disc.} Deduce the mass $M$ of a uniform disc of radius $0.30$ m having this moment of inertia. Comment on whether such a wheel could reasonably be made of stone or of hardwood.
\item \textbf{Initial angular velocity.} Compute the angular velocity $\omega_0$ just after the kick, in rad\,s$^{-1}$ and in revolutions per minute (rev/min).
\item \textbf{Deceleration and stopping time.} Using $C=I\alpha$, find the angular deceleration $\alpha$. Check, with the kinematic relation $\omega^2=\omega_0^2+2\alpha\theta$, that the wheel indeed stops after exactly $40$ revolutions, and find the time taken to stop.
\item \textbf{Consistency check.} Recover the same stopping angle by a purely energetic argument and confirm that the two methods agree. Why must they agree?
\item \textbf{The lump of clay.} The wheel is spinning at $\omega_0$ when the $2$ kg lump of clay is dropped vertically onto the rim and sticks. Explain why the moment of momentum about the axle is conserved during the drop, even though the resisting couple acts. Compute the new angular speed $\omega_1$ and the kinetic energy lost in the process. Where does this energy go?
\item \textbf{Design discussion: disc or ring?} Suppose instead that the same mass $M$ is concentrated at the rim (a heavy ring with light spokes). Compute the new moment of inertia, the new number of revolutions after the same $120$ J kick, and the new initial angular speed. Which design is better for Pa Nformi, and what practical disadvantages might the ring design have?
\item \textbf{Final report.} Write a short justified recommendation to Pa Nformi: shape, mass, radius, expected performance.
\end{enumerate}

\courssub{Model answer}

\textbf{1. Minimum moment of inertia.}\par
The wheel starts with kinetic energy $E_0=120$ J and stops after turning through an angle $\theta$. The only work done is that of the resisting couple: $W=-C_r\theta$. The work--energy principle gives
\[
0 - E_0 = -C_r\theta \quad\Longrightarrow\quad \theta = \frac{E_0}{C_r}.
\]
With $C_r=0.5\ \mathrm{N\,m}$, $\theta = 120/0.5 = 240$ rad, which corresponds to
\[
N = \frac{\theta}{2\pi} = \frac{240}{2\pi} = \frac{120}{\pi} \approx 38.2\ \text{revolutions}.
\]
This is \emph{less} than the required $40$ revolutions. Crucially, $\theta$ (and hence $N$) does not depend on the moment of inertia $I$ at all. Therefore \textbf{no choice of $I$ can make the wheel reach 40 revolutions} with the given $E_0$ and $C_r$. The specification is infeasible as written.

To meet the requirement we must either increase the kick energy or reduce the resisting couple. The maximum allowable resisting couple for $40$ revolutions is
\[
C_{r,\max} = \frac{E_0}{2\pi N} = \frac{120}{80\pi} = \frac{1.5}{\pi} \approx 0.477\ \mathrm{N\,m}.
\]
We assume the bearings are improved so that $C_r' = 0.45\ \mathrm{N\,m} < C_{r,\max}$. Then the wheel will stop after
\[
\theta = \frac{120}{0.45} = 266.\overline{6}\ \mathrm{rad} \approx 42.4\ \text{revolutions},
\]
satisfying the requirement with a margin. We continue the design with $C_r' = 0.45\ \mathrm{N\,m}$.

\begin{attention}[A lesson in modelling]
Always check that a specification is \emph{feasible} before computing detailed parameters. Here the work--energy test immediately shows the original brief is impossible: no moment of inertia, however large, can make $120$ J do more than $120$ J of work. A larger $I$ makes the wheel spin \emph{slower} but for the \emph{same total angle}.
\end{attention}

\textbf{2. Mass of the disc.}\par
The moment of inertia does not affect the number of revolutions, but Pa Nformi also wants a wheel that spins \emph{slowly enough to work on comfortably}. We impose a maximum initial angular speed of $150$ rev/min $= 5\pi\ \mathrm{rad\,s^{-1}}$. From $E_0 = \tfrac12 I\omega_0^2$,
\[
I = \frac{2E_0}{\omega_0^2} = \frac{240}{(5\pi)^2} = \frac{240}{25\pi^2} = \frac{9.6}{\pi^2} \approx 0.972\ \mathrm{kg\,m^2}.
\]
For a uniform disc, $I = \tfrac12 MR^2$, so
\[
M = \frac{2I}{R^2} = \frac{2 \times 0.972}{0.09} \approx 21.6\ \mathrm{kg}.
\]
A stone disc of radius $0.30$ m and mass about $22$ kg is entirely realistic for a traditional kick-wheel (density of stone $\approx 2500\ \mathrm{kg\,m^{-3}}$ gives a thickness of $\approx 3$ cm).

\textbf{3. Initial angular velocity.}\par
With $I = 0.972\ \mathrm{kg\,m^2}$ and $E_0 = 120$ J:
\[
\omega_0 = \sqrt{\frac{2E_0}{I}} = \sqrt{\frac{240}{0.972}} \approx 15.7\ \mathrm{rad\,s^{-1}} = 150\ \mathrm{rev/min},
\]
consistent with our design choice.

\textbf{4. Deceleration and stopping time.}\par
The equation of motion $C = I\alpha$ with $C = -C_r' = -0.45\ \mathrm{N\,m}$ gives
\[
\alpha = \frac{-0.45}{0.972} \approx -0.463\ \mathrm{rad\,s^{-2}}.
\]
Kinematic check using $\omega^2 = \omega_0^2 + 2\alpha\theta$ with $\omega=0$:
\[
\theta = \frac{\omega_0^2}{2|\alpha|} = \frac{(15.7)^2}{2 \times 0.463} \approx 266\ \mathrm{rad} \approx 42.4\ \text{revolutions} \ge 40.\quad \checkmark
\]
Stopping time:
\[
t = \frac{\omega_0}{|\alpha|} = \frac{15.7}{0.463} \approx 33.9\ \mathrm{s}.
\]

\textbf{5. Consistency check.}\par
Energetically, $\theta = \dfrac{E_0}{C_r'} = \dfrac{120}{0.45} \approx 266.7$ rad — identical to the kinematic result. The two methods must agree because the kinematic relation $\omega^2 = \omega_0^2 + 2\alpha\theta$, multiplied by $\tfrac12 I$, \emph{is} the work--energy theorem:
\[
\tfrac12 I\omega^2 - \tfrac12 I\omega_0^2 = I\alpha\,\theta = C\theta.
\]

\textbf{6. The lump of clay.}\par
During the very short time $\Delta t$ in which the clay is dropped and sticks, the impulse of the resisting couple $C_r'\,\Delta t$ is negligible. The forces between clay and wheel are internal to the system (wheel + clay), so the moment of momentum about the axle is conserved:
\[
I\omega_0 = (I + mR^2)\,\omega_1.
\]
Here $mR^2 = 2 \times 0.09 = 0.18\ \mathrm{kg\,m^2}$, so $I_{\text{total}} = 0.972 + 0.18 = 1.152\ \mathrm{kg\,m^2}$.
\[
\omega_1 = \frac{I}{I + mR^2}\,\omega_0 = \frac{0.972}{1.152} \times 15.7 = \frac{27}{32} \times 5\pi = \frac{135\pi}{32} \approx 13.25\ \mathrm{rad\,s^{-1}} \approx 126.5\ \mathrm{rev/min}.
\]
Energy before: $E_0 = 120$ J. Energy after:
\[
E_1 = \tfrac12 (I + mR^2)\omega_1^2 = \tfrac12 \times 1.152 \times (13.25)^2 \approx 101.25\ \mathrm{J}.
\]
Energy lost: $\Delta E = 120 - 101.25 = 18.75$ J. This energy is dissipated as heat and sound during the inelastic ``collision'' while the clay slips on the wheel before sticking. Note that angular momentum is conserved but kinetic energy is \emph{not}.

\textbf{7. Disc or ring?}\par
With the same mass $M = 21.6$ kg concentrated at the rim (a heavy ring with light spokes), the moment of inertia is
\[
I_{\text{ring}} = MR^2 = 21.6 \times 0.09 = 1.944\ \mathrm{kg\,m^2},
\]
exactly \emph{double} the disc value. For the same $120$ J kick:
\[
\omega_0' = \sqrt{\frac{2E_0}{I_{\text{ring}}}} = \sqrt{\frac{240}{1.944}} \approx 11.11\ \mathrm{rad\,s^{-1}} \approx 106\ \mathrm{rev/min},
\]
a gentler working speed. The number of revolutions is unchanged ($\approx 42.4$ rev), since it depends only on $E_0$ and $C_r'$; but the stopping time increases by a factor $\sqrt{2}$ to
\[
t' = \frac{\omega_0'}{|\alpha'|} = \frac{\omega_0'}{C_r'/I_{\text{ring}}} = \frac{I_{\text{ring}}\omega_0'}{C_r'} = \sqrt{2}\, \frac{I_{\text{disc}}\omega_0}{C_r'} \approx 1.414 \times 33.9 \approx 48.0\ \mathrm{s},
\]
giving over a minute of useful spinning. The lower speed is more comfortable for shaping clay. Disadvantages: a rim-loaded wheel is harder to build and balance, the spokes must be strong enough to transmit the torque, and the rim is more exposed to knocks.

\textbf{8. Final recommendation.}\par
Improve the bearings to bring the resisting couple down to about $0.45\ \mathrm{N\,m}$; build a rim-loaded wheel (hoop) of mass $\approx 22$ kg and radius $0.30$ m. One $120$ J kick then gives $\approx 106$ rev/min initially, more than $40$ revolutions, and about $48$ seconds of working time — with the bonus that dropping a $2$ kg lump of clay on the rim slows the wheel by only about $16\%$ (from $106$ to $\approx 89$ rev/min).

\begin{savaistu}[Flywheels everywhere]
Pa Nformi's wheel is a \cle{flywheel}: a store of rotational kinetic energy. The same physics governs the flywheels that smooth the motion of engines, the massive stone wheels of grinding mills, and modern flywheel energy-storage systems used to stabilise electrical grids — including experimental installations paired with solar mini-grids in rural areas. The trade-off you discovered — for a fixed energy, more inertia means lower speed but longer running time — is exactly the trade-off engineers face when sizing any flywheel.
\end{savaistu}

\newpage

\coursec{Methods and worked examples}

\coursec{Moment of Inertia and Radius of Gyration}

\courssub{Definition and Computation by Integration}

\begin{methode}[Moment of inertia by summation or integration]
\begin{enumerate}
\item Write the definition $I=\sum m_i r_i^2$ for discrete masses, or $I=\int r^2\,dm$ for a continuous body, where $r$ is the \emph{perpendicular} distance from the axis.
\item For a continuous body, choose a thin element (strip, ring, shell) whose every point is at the same distance $r$ from the axis; express its mass $dm$ using the density (linear $\lambda$, surface $\sigma$, or volume $\rho$).
\item Express $dm$ in terms of a single variable and integrate between the correct limits.
\item Check the dimensions ($\mathrm{kg\,m^2}$) and, where possible, compare with a standard result.
\end{enumerate}
\textbf{Reflexes:} for a thin rod of length $2a$ about its centre use strips; for a disc use concentric rings; for a sphere use slices. Always state the axis clearly.
\end{methode}

\begin{popfigure}
\begin{tikzpicture}[scale=1.5]
\draw[thick] (-1.5,0) -- (1.5,0);
\draw[fill] (0,0) circle (1pt) node[below] {$O$};
\draw[<->] (-1.5,-0.3) -- (1.5,-0.3) node[midway,below] {$2a$};
\draw[<->] (0,0.2) -- (0.8,0.2) node[midway,above] {$x$};
\draw[thick,popblue] (0.8,-0.1) -- (0.8,0.1);
\draw[thick,popblue] (1.0,-0.1) -- (1.0,0.1);
\node[popblue,above] at (0.9,0.15) {$dx$};
\node[left] at (-1.5,0) {$A$};
\node[right] at (1.5,0) {$B$};
\end{tikzpicture}
\legende{Element of a uniform rod for moment of inertia calculation.}
\end{popfigure}

\begin{exoresolu}[Moment of inertia of a uniform rod about its centre]
A uniform thin rod $AB$ has mass $M$ and length $2a$. By integration, find its moment of inertia about an axis through its centre, perpendicular to its length.
\tcblower
Let the axis pass through the centre $O$ of the rod, perpendicular to $AB$. The linear density is
\[
\lambda=\frac{M}{2a}.
\]
Take a small element of length $dx$ at distance $x$ from $O$. Its mass is
\[
dm=\lambda\,dx=\frac{M}{2a}\,dx,
\]
and every point of this element is at distance $x$ from the axis, so its contribution is $dI=x^2\,dm$.

Integrating over the whole rod, from $x=-a$ to $x=a$:
\begin{align*}
I&=\int_{-a}^{a}x^2\,\frac{M}{2a}\,dx
=\frac{M}{2a}\left[\frac{x^3}{3}\right]_{-a}^{a}\\[2pt]
&=\frac{M}{2a}\cdot\frac{2a^3}{3}
=\frac{1}{3}Ma^2.
\end{align*}
Hence $\boxed{I=\tfrac13 Ma^2}$ about a perpendicular axis through the centre.

\textbf{Check:} by the parallel axis theorem, about an end the moment of inertia is $I=\tfrac13 Ma^2+Ma^2=\tfrac43 Ma^2$, the standard result.
\end{exoresolu}

\courssub{Radius of Gyration and the Two Axis Theorems}

\begin{methode}[Radius of gyration, parallel and perpendicular axes]
\begin{enumerate}
\item Write $I=Mk^2$: the radius of gyration $k$ concentrates all the mass at distance $k$ from the axis.
\item \textbf{Parallel axis theorem:} $I=I_G+Mh^2$, where $I_G$ is about the parallel axis through the centre of mass $G$ and $h$ is the distance between the axes.
\item \textbf{Perpendicular axis theorem} (lamina only, axis $Oz$ perpendicular to the lamina): $I_z=I_x+I_y$.
\item Combine: to move between two axes neither through $G$, go via $G$: $I_2=I_1-Mh_1^2+Mh_2^2$.
\end{enumerate}
\textbf{Reflex:} $I_G$ is always the \emph{smallest} moment of inertia among parallel axes.
\end{methode}

\begin{exoresolu}[Radius of gyration of a rectangular lamina]
A uniform rectangular lamina has sides $2a$ and $2b$ and mass $M$. Find (a) its radius of gyration about the axis through its centre perpendicular to its plane; (b) its moment of inertia about a parallel axis through a corner.
\tcblower
For the lamina in the $xy$-plane with centre at the origin, the standard results about axes through $G$ in the plane are
\[
I_x=\tfrac13 Mb^2,\qquad I_y=\tfrac13 Ma^2.
\]
(a) By the perpendicular axis theorem,
\[
I_z=I_x+I_y=\tfrac13 M(a^2+b^2).
\]
Since $I_z=Mk^2$,
\[
\boxed{k=\sqrt{\frac{a^2+b^2}{3}}}.
\]
(b) The corner is at distance $h=\sqrt{a^2+b^2}$ from $G$. By the parallel axis theorem,
\begin{align*}
I_{\text{corner}}&=I_z+Mh^2
=\tfrac13 M(a^2+b^2)+M(a^2+b^2)\\
&=\boxed{\tfrac43 M(a^2+b^2)}.
\end{align*}
\end{exoresolu}

\coursec{Standard Moments of Inertia by Integration}

\courssub{Uniform Disc about its Symmetry Axis}

\begin{popfigure}
\begin{tikzpicture}[scale=1.2]
\draw[thick] (0,0) circle (2);
\draw[thick,popblue] (0,0) circle (1.5);
\draw[thick,popblue] (0,0) circle (1.6);
\draw[fill=popblue!20] (0,0) circle (1.6);
\draw[fill=white] (0,0) circle (1.5);
\node at (0,0) {$O$};
\draw[<->] (0,0) -- (1.55,0) node[midway,above] {$r$};
\draw[<->] (1.5,0) -- (1.6,0) node[midway,above] {$dr$};
\node[right] at (2,0) {$R$};
\end{tikzpicture}
\legende{Concentric ring element for a uniform disc.}
\end{popfigure}

\begin{exoresolu}[Moment of inertia of a uniform disc]
A uniform disc of mass $M$ and radius $R$ rotates about an axis through its centre perpendicular to its plane. Derive $I=\frac12 MR^2$ by integration.
\tcblower
Surface density $\sigma = \dfrac{M}{\pi R^2}$. Consider a concentric ring of radius $r$ and width $dr$. Its area is $2\pi r\,dr$, so its mass is
\[
dm = \sigma \cdot 2\pi r\,dr = \frac{2M}{R^2}\,r\,dr.
\]
Every point of the ring is at distance $r$ from the axis, hence
\[
dI = r^2\,dm = \frac{2M}{R^2}\,r^3\,dr.
\]
Integrate from $r=0$ to $r=R$:
\[
I = \frac{2M}{R^2}\int_0^R r^3\,dr = \frac{2M}{R^2}\cdot\frac{R^4}{4} = \frac12 MR^2.
\]
Thus $\boxed{I=\tfrac12 MR^2}$.
\end{exoresolu}

\courssub{Uniform Rectangular Lamina about Perpendicular Axis through Centre}

\begin{exoresolu}[Moment of inertia of a rectangular lamina by integration]
A uniform rectangular lamina of mass $M$ has sides $2a$ and $2b$. Derive the moments of inertia about axes through its centre parallel to the sides, and hence the moment of inertia about the perpendicular axis through the centre.
\tcblower
Place the lamina in the $xy$-plane with centre at the origin, sides parallel to the axes. Surface density $\sigma = \dfrac{M}{4ab}$.

\textbf{About $Ox$ (parallel to side $2a$):} Take a strip parallel to $Ox$ at distance $y$, width $dy$, length $2a$. Its mass $dm = \sigma \cdot 2a\,dy$. Every point is at distance $|y|$ from $Ox$, so
\[
I_x = \int_{-b}^{b} y^2\,dm = 2a\sigma\int_{-b}^{b} y^2\,dy = 2a\sigma\cdot\frac{2b^3}{3} = \frac{4a\sigma b^3}{3}.
\]
Substitute $\sigma = M/(4ab)$:
\[
I_x = \frac{4a}{3}\cdot\frac{M}{4ab}\cdot b^3 = \frac13 Mb^2.
\]

\textbf{About $Oy$:} By symmetry, $I_y = \frac13 Ma^2$.

\textbf{About $Oz$ (perpendicular through centre):} By the perpendicular axis theorem,
\[
I_z = I_x + I_y = \frac13 M(a^2+b^2).
\]
These are the standard results used in the previous section.
\end{exoresolu}

\courssub{Uniform Solid Sphere about a Diameter}

\begin{exoresolu}[Moment of inertia of a uniform solid sphere]
A uniform solid sphere of mass $M$ and radius $R$. Derive $I=\frac25 MR^2$ about a diameter.
\tcblower
Volume density $\rho = \dfrac{M}{\frac43\pi R^3} = \dfrac{3M}{4\pi R^3}$. Slice the sphere into discs perpendicular to the chosen diameter (the $x$-axis). A disc at distance $x$ from the centre has radius $y = \sqrt{R^2-x^2}$ and thickness $dx$. Its volume is $\pi y^2\,dx$, mass
\[
dm = \rho\pi(R^2-x^2)\,dx.
\]
The moment of inertia of this disc about its own axis (the $x$-axis) is $\frac12 dm\,y^2$. Hence
\[
dI = \frac12 \rho\pi (R^2-x^2)^2\,dx.
\]
Integrate from $x=-R$ to $x=R$:
\begin{align*}
I &= \frac12 \rho\pi \int_{-R}^{R} (R^2-x^2)^2\,dx
= \rho\pi \int_{0}^{R} (R^4 - 2R^2x^2 + x^4)\,dx \\
&= \rho\pi \left[ R^4x - \frac{2}{3}R^2x^3 + \frac{1}{5}x^5 \right]_0^R
= \rho\pi R^5\left(1 - \frac23 + \frac15\right) \\
&= \rho\pi R^5\cdot\frac{8}{15}.
\end{align*}
Substitute $\rho = \dfrac{3M}{4\pi R^3}$:
\[
I = \frac{3M}{4\pi R^3}\cdot\pi R^5\cdot\frac{8}{15} = \frac{24}{60}MR^2 = \frac25 MR^2.
\]
Thus $\boxed{I=\tfrac25 MR^2}$.
\end{exoresolu}

\courssub{Uniform Spherical Shell about a Diameter}

\begin{exoresolu}[Moment of inertia of a uniform spherical shell]
A uniform thin spherical shell of mass $M$ and radius $R$. Derive $I=\frac23 MR^2$ about a diameter.
\tcblower
Surface density $\sigma = \dfrac{M}{4\pi R^2}$. Use spherical coordinates: a ring at polar angle $\theta$ has radius $R\sin\theta$, width $R\,d\theta$. Its area is $2\pi R\sin\theta \cdot R\,d\theta = 2\pi R^2\sin\theta\,d\theta$, mass
\[
dm = \sigma \cdot 2\pi R^2\sin\theta\,d\theta = \frac{M}{2}\sin\theta\,d\theta.
\]
Distance from the axis (diameter) is $R\sin\theta$, so
\[
dI = (R\sin\theta)^2\,dm = \frac{M}{2}R^2\sin^3\theta\,d\theta.
\]
Integrate $\theta$ from $0$ to $\pi$:
\[
I = \frac{M}{2}R^2 \int_0^\pi \sin^3\theta\,d\theta
= \frac{M}{2}R^2 \int_0^\pi \sin\theta(1-\cos^2\theta)\,d\theta.
\]
Let $u=\cos\theta$, $du=-\sin\theta\,d\theta$:
\[
\int_0^\pi \sin^3\theta\,d\theta = \int_{1}^{-1} -(1-u^2)\,du = \int_{-1}^{1}(1-u^2)\,du = \left[u-\frac{u^3}{3}\right]_{-1}^{1} = \frac43.
\]
Hence
\[
I = \frac{M}{2}R^2 \cdot \frac43 = \frac23 MR^2.
\]
Thus $\boxed{I=\tfrac23 MR^2}$.
\end{exoresolu}

\begin{aretenir}[Summary of standard moments of inertia]
\begin{center}
\begin{tabularx}{\linewidth}{|l|X|X|}\hline
\rowcolor{popblueL}
\textbf{Body} & \textbf{Axis} & \textbf{Moment of Inertia} \\ \hline
Uniform thin rod, mass $M$, length $2a$ & Through centre, $\perp$ to rod & $\frac13 Ma^2$ \\ \hline
Uniform thin rod, mass $M$, length $2a$ & Through one end, $\perp$ to rod & $\frac43 Ma^2$ \\ \hline
Uniform disc, mass $M$, radius $R$ & Through centre, $\perp$ to plane & $\frac12 MR^2$ \\ \hline
Uniform disc, mass $M$, radius $R$ & Diameter (in plane) & $\frac14 MR^2$ \\ \hline
Uniform rectangular lamina, mass $M$, sides $2a$, $2b$ & Through centre, $\perp$ to plane & $\frac13 M(a^2+b^2)$ \\ \hline
Uniform solid sphere, mass $M$, radius $R$ & Diameter & $\frac25 MR^2$ \\ \hline
Uniform spherical shell, mass $M$, radius $R$ & Diameter & $\frac23 MR^2$ \\ \hline
Uniform solid cylinder, mass $M$, radius $R$ & Axis of symmetry & $\frac12 MR^2$ \\ \hline
\end{tabularx}
\end{center}
\end{aretenir}

\coursec{Equation of Motion under a Constant Torque}

\courssub{Newton's Second Law for Rotation}

\begin{methode}[Rigid body rotating about a fixed axis under constant torque]
\begin{enumerate}
\item Compute the total moment of inertia $I$ about the fixed axis (sum contributions of all parts).
\item Compute the resultant torque $\Gamma$ about the axis (watch the signs: choose a positive sense of rotation).
\item Apply the equation of motion $\Gamma=I\ddot{\theta}$, i.e. $\ddot{\theta}=\Gamma/I$ (constant angular acceleration).
\item Use the constant-angular-acceleration formulae, analogues of the linear ones:
\[
\omega=\omega_0+\alpha t,\quad \theta=\omega_0 t+\tfrac12\alpha t^2,\quad \omega^2=\omega_0^2+2\alpha\theta .
\]
\end{enumerate}
\textbf{Reflex:} a torque from a tangential force $F$ at distance $r$ is $\Gamma=Fr$; a couple of forces $F$ separated by $d$ gives $\Gamma=Fd$.
\end{methode}

\begin{exoresolu}[Flywheel accelerated by a constant torque]
A flywheel of moment of inertia $I=0.8\ \mathrm{kg\,m^2}$ is mounted on a fixed horizontal axle. A constant torque of $6\ \mathrm{N\,m}$ is applied while a frictional couple of $1\ \mathrm{N\,m}$ opposes the motion. Starting from rest, find (a) the angular acceleration; (b) the angular velocity after $10\ \mathrm{s}$; (c) the number of revolutions made in that time.
\tcblower
(a) Resultant torque: $\Gamma=6-1=5\ \mathrm{N\,m}$. Then
\[
\alpha=\ddot{\theta}=\frac{\Gamma}{I}=\frac{5}{0.8}=6.25\ \mathrm{rad\,s^{-2}}.
\]
(b) With $\omega_0=0$:
\[
\omega=\alpha t=6.25\times 10=\boxed{62.5\ \mathrm{rad\,s^{-1}}}.
\]
(c) Angle turned through:
\[
\theta=\tfrac12\alpha t^2=\tfrac12(6.25)(100)=312.5\ \mathrm{rad}.
\]
Number of revolutions:
\[
N=\frac{\theta}{2\pi}=\frac{312.5}{2\pi}\approx \boxed{49.7\ \text{revolutions}}.
\]
\end{exoresolu}

\coursec{Moment of Momentum and Its Conservation; Impulse}

\courssub{Angular Momentum and Angular Impulse}

\begin{methode}[Angular impulse and moment of momentum]
\begin{enumerate}
\item Moment of momentum (angular momentum) about a fixed axis: $L=I\omega$.
\item Angular impulse of a torque: $\displaystyle \int \Gamma\,dt=\Delta L=I\omega_2-I\omega_1$.
\item For a particle of mass $m$ moving with speed $v$ at perpendicular distance $r$ from the axis, $L=mvr$.
\item For connected systems (string over a pulley with inertia), write \emph{separate} equations: $ma=mg-T$ type for each particle, and $(T_1-T_2)r=I\alpha$ for the pulley, with $a=r\alpha$ (no slipping).
\end{enumerate}
\textbf{Reflex:} when the pulley has inertia, the tensions on the two sides are \emph{different}.
\end{methode}

\begin{exoresolu}[Masses hanging over a massive pulley]
A light inextensible string passes over a pulley of mass $2\ \mathrm{kg}$ and radius $0.1\ \mathrm{m}$, modelled as a uniform disc ($I=\tfrac12 Mr^2$). Masses of $3\ \mathrm{kg}$ and $1\ \mathrm{kg}$ hang from the ends. The string does not slip on the pulley. Taking $g=9.8\ \mathrm{m\,s^{-2}}$, find the acceleration of the masses and the two tensions.
\tcblower
Moment of inertia of the pulley:
\[
I=\tfrac12(2)(0.1)^2=0.01\ \mathrm{kg\,m^2}.
\]
Let $T_1$ act on the $3\ \mathrm{kg}$ side, $T_2$ on the $1\ \mathrm{kg}$ side, and let $a$ be the acceleration of the $3\ \mathrm{kg}$ mass downwards.

\textbf{Particle equations} (Newton's second law):
\begin{align*}
3g-T_1&=3a,\\
T_2-g&=a.
\end{align*}
\textbf{Pulley equation} (torque $=I\alpha$, with $\alpha=a/r$):
\[
(T_1-T_2)r=I\frac{a}{r}
\quad\Longrightarrow\quad
T_1-T_2=\frac{Ia}{r^2}=\frac{0.01}{0.01}a=a.
\]
Adding the three equations eliminates $T_1,T_2$:
\[
3g-g=3a+a+a=5a
\quad\Longrightarrow\quad
a=\frac{2g}{5}=\frac{19.6}{5}=\boxed{3.92\ \mathrm{m\,s^{-2}}}.
\]
Then
\[
T_1=3(g-a)=3(5.88)=\boxed{17.64\ \mathrm{N}},\qquad
T_2=g+a=13.72\ \mathrm{N}.
\]
Check: $T_1-T_2=3.92=a$ as required. Note $T_1\neq T_2$ because the pulley has inertia.
\end{exoresolu}

\courssub{Conservation of Angular Momentum}

\begin{methode}[Conservation of angular momentum]
\begin{enumerate}
\item Check that the \emph{external} torque about the axis is zero (internal forces, however large, do not change $L$).
\item Write $I_1\omega_1=I_2\omega_2$ where $I$ may change (masses moving in/out, body changing shape).
\item For a point mass $m$ added to or removed from a rotating system, adjust $I$ accordingly: $I_{\text{new}}=I_{\text{old}}+mr^2$.
\item Kinetic energy is \emph{not} conserved in such processes (internal forces do work): expect $KE$ to change.
\end{enumerate}
\end{methode}

\begin{exoresolu}[Mass dropped onto a rotating turntable]
A horizontal turntable of moment of inertia $0.5\ \mathrm{kg\,m^2}$ rotates freely at $4\ \mathrm{rad\,s^{-1}}$ about a vertical axis. A small bag of sand of mass $2\ \mathrm{kg}$ is dropped vertically onto it and sticks at a distance $0.3\ \mathrm{m}$ from the axis. Find the new angular velocity and the loss of kinetic energy.
\tcblower
No external torque acts about the vertical axis, so moment of momentum is conserved.

New moment of inertia:
\[
I_2=0.5+mr^2=0.5+2(0.3)^2=0.5+0.18=0.68\ \mathrm{kg\,m^2}.
\]
Conservation of $L$:
\[
I_1\omega_1=I_2\omega_2
\quad\Longrightarrow\quad
\omega_2=\frac{0.5\times 4}{0.68}=\frac{2}{0.68}\approx \boxed{2.94\ \mathrm{rad\,s^{-1}}}.
\]
Kinetic energies:
\[
KE_1=\tfrac12 I_1\omega_1^2=\tfrac12(0.5)(16)=4\ \mathrm{J},
\]
\[
KE_2=\tfrac12(0.68)(2.94)^2\approx 2.94\ \mathrm{J}.
\]
Loss of kinetic energy: $4-2.94\approx \boxed{1.06\ \mathrm{J}}$, dissipated as heat and sound when the bag grips the turntable (an inelastic process).
\end{exoresolu}

\coursec{Energy of a Rotating Body and the Force on the Axis}

\courssub{Rotational Kinetic Energy and Conservation of Energy}

\courssub{Force Exerted on the Axis of Rotation}

\begin{methode}[Conservation of energy for a rotating rigid body]
\begin{enumerate}
\item Write the kinetic energy as $KE=\tfrac12 I\omega^2$ (rotation about a fixed axis); add $\tfrac12 Mv_G^2$ if the body also translates.
\item Write the change in gravitational potential energy from the vertical displacement of the \emph{centre of mass}: $\Delta PE=Mg\,\Delta h_G$.
\item If only gravity does work, $KE+PE=$ constant; use this to find $\omega$ in any position.
\item To find the \textbf{force on the axis}: use the motion of the centre of mass,
\[
\vec{F}_{\text{axis}}+M\vec{g}=M\vec{a}_G,
\]
with $a_G$ having radial component $h\omega^2$ (towards the axis) and tangential component $h\alpha$, where $h=OG$.
\end{enumerate}
\end{methode}

\begin{exoresolu}[Rod falling from the vertical; force on the hinge]
A uniform rod of mass $M$ and length $2a$ is freely hinged at one end $O$ and released from rest in the vertical (unstable) position. Find (a) its angular velocity when it has turned through an angle $\theta$; (b) the horizontal and vertical components of the reaction at the hinge when $\theta=90^{\circ}$.
\tcblower
Moment of inertia about $O$: $I_O=\tfrac43 Ma^2$; the centre of mass $G$ is at distance $a$ from $O$.

(a) When the rod makes angle $\theta$ with the upward vertical, $G$ has fallen a distance $a(1-\cos\theta)$. Energy conservation:
\[
\tfrac12 I_O\omega^2=Mga(1-\cos\theta)
\quad\Longrightarrow\quad
\omega^2=\frac{3g}{2a}(1-\cos\theta).
\]
(b) At $\theta=90^{\circ}$: $\omega^2=\dfrac{3g}{2a}$.

Angular acceleration from $\Gamma=I\alpha$ at $\theta=90^{\circ}$ (weight has moment $Mga$ about $O$):
\[
Mga=\tfrac43 Ma^2\alpha
\quad\Longrightarrow\quad
\alpha=\frac{3g}{4a}.
\]
Acceleration of $G$ (at $\theta=90^\circ$, $G$ moves horizontally; radial direction is horizontal, tangential is vertical):
\[
a_{\text{rad}}=a\omega^2=\frac{3g}{2}\ (\text{horizontal, towards }O),\qquad
a_{\text{tan}}=a\alpha=\frac{3g}{4}\ (\text{vertically downwards}).
\]
Newton's second law for $G$, with hinge reaction $(X,Y)$:
\[
X=Ma_{\text{rad}}=\frac{3}{2}Mg,\qquad
Mg-Y=Ma_{\text{tan}}=\frac{3}{4}Mg
\ \Longrightarrow\ 
Y=\frac{1}{4}Mg.
\]
Hence the hinge exerts $\boxed{X=\tfrac32 Mg}$ horizontally and $\boxed{Y=\tfrac14 Mg}$ vertically upwards; magnitude $R=\tfrac{Mg}{4}\sqrt{37}$.
\end{exoresolu}

\coursec{The Compound Pendulum and Synthesis}

\courssub{Period of a Compound Pendulum}

\begin{methode}[Period of a compound pendulum]
\begin{enumerate}
\item Identify the pivot $O$; compute $I_O$ about the pivot (often via the parallel axis theorem: $I_O=I_G+Mh^2=M(k_G^2+h^2)$, with $h=OG$).
\item Write the restoring torque for angular displacement $\theta$: $\Gamma=-Mgh\sin\theta$.
\item For small oscillations, $\sin\theta\approx\theta$, giving $\ddot{\theta}=-\dfrac{Mgh}{I_O}\theta$: SHM with
\[
T=2\pi\sqrt{\frac{I_O}{Mgh}}=2\pi\sqrt{\frac{k_G^2+h^2}{gh}}.
\]
\item The \textbf{equivalent simple pendulum} has length $\ell=\dfrac{k_G^2+h^2}{h}$; the centre of oscillation $O'$ lies at distance $\ell$ from $O$, and $O$ and $O'$ are interchangeable (Kater's idea).
\end{enumerate}
\textbf{Reflex:} always check that the small-angle approximation is stated; the formula is only valid for small oscillations.
\end{methode}

\begin{exoresolu}[Compound pendulum: uniform rod swinging about one end]
A uniform rod of length $1.2\ \mathrm{m}$ swings as a compound pendulum about a horizontal axis through one end. Taking $g=9.8\ \mathrm{m\,s^{-2}}$, find (a) the period of small oscillations; (b) the length of the equivalent simple pendulum; (c) the distance from the pivot of the centre of oscillation.
\tcblower
Let $2a=1.2\ \mathrm{m}$, so $a=0.6\ \mathrm{m}$ and $h=OG=a=0.6\ \mathrm{m}$.

About the centre: $I_G=\tfrac13 Ma^2$, so $k_G^2=\tfrac13 a^2=0.12\ \mathrm{m^2}$.

(a) Period:
\[
T=2\pi\sqrt{\frac{k_G^2+h^2}{gh}}
=2\pi\sqrt{\frac{0.12+0.36}{9.8\times 0.6}}
=2\pi\sqrt{\frac{0.48}{5.88}}.
\]
\[
\frac{0.48}{5.88}=0.08163\ \mathrm{s^2},\qquad
T=2\pi(0.2857)\approx \boxed{1.79\ \mathrm{s}}.
\]
(b) Equivalent simple pendulum length:
\[
\ell=\frac{k_G^2+h^2}{h}=\frac{0.48}{0.6}=\boxed{0.8\ \mathrm{m}}.
\]
(c) The centre of oscillation is at distance $\ell=0.8\ \mathrm{m}$ below the pivot, i.e. $0.2\ \mathrm{m}$ below $G$.

\textbf{Check (Kater's property):} if the rod were swung about an axis through this point $O'$, with $h'=0.2\ \mathrm{m}$, then
\[
\frac{k_G^2+h'^2}{h'}=\frac{0.12+0.04}{0.2}=0.8\ \mathrm{m}=\ell,
\]
giving the same period, as expected: pivot and centre of oscillation are interchangeable.
\end{exoresolu}

\courssub{Combined Rotation and Translation: Rolling Motion}

\begin{methode}[Rigid body rolling or rotating with connected motion]
\begin{enumerate}
\item Draw a clear diagram; mark all forces and the axis/ instantaneous constraints.
\item Choose your tool: energy conservation (cleanest when only gravity does work and no slipping dissipates energy) or equations of motion ($F=Ma_G$ and $\Gamma_G=I_G\alpha$).
\item For rolling without slipping, impose the constraint $v_G=r\omega$, $a_G=r\alpha$.
\item Total kinetic energy: $KE=\tfrac12 Mv_G^2+\tfrac12 I_G\omega^2$.
\end{enumerate}
\end{methode}

\begin{exoresolu}[Cylinder rolling down an inclined plane]
A solid cylinder of mass $M$ and radius $r$ rolls without slipping down a plane inclined at $30^{\circ}$ to the horizontal, starting from rest. Find (a) its acceleration; (b) its speed after travelling $2\ \mathrm{m}$ down the slope ($g=9.8\ \mathrm{m\,s^{-2}}$); (c) the minimum coefficient of friction for rolling to be maintained.
\tcblower
For a solid cylinder, $I_G=\tfrac12 Mr^2$. Rolling constraint: $a=r\alpha$.

(a) Equations of motion along the plane and about $G$:
\[
Mg\sin 30^{\circ}-F=Ma,\qquad Fr=I_G\alpha=\tfrac12 Mr^2\cdot\frac{a}{r}=\tfrac12 Mra.
\]
From the second equation $F=\tfrac12 Ma$. Substituting:
\[
Mg\sin30^{\circ}=Ma+\tfrac12 Ma=\tfrac32 Ma
\quad\Longrightarrow\quad
a=\frac{2g\sin30^{\circ}}{3}=\frac{g}{3}=\boxed{3.27\ \mathrm{m\,s^{-2}}}.
\]
(b) With $u=0$, $s=2\ \mathrm{m}$: $v^2=2as=2\times3.27\times2=13.07$, so
\[
v\approx \boxed{3.62\ \mathrm{m\,s^{-1}}}.
\]
(Energy check: $Mgh=\tfrac12 Mv^2+\tfrac12 I_G\omega^2=\tfrac34 Mv^2$ with $h=2\sin30^{\circ}=1\ \mathrm{m}$ gives $v^2=\tfrac43 g=13.07$.)

(c) Friction needed: $F=\tfrac12 Ma=\tfrac16 Mg$. Normal reaction: $N=Mg\cos30^{\circ}$. Rolling requires $F\le\mu N$:
\[
\mu\ge \frac{F}{N}=\frac{Mg/6}{Mg\cos30^{\circ}}=\frac{1}{6\times0.866}\approx \boxed{0.19}.
\]
\end{exoresolu}

\begin{attention}[Examiner's warnings]
\begin{itemize}
\item Never confuse $I_G$ with $I_O$: apply the parallel axis theorem carefully before writing $\Gamma=I\ddot\theta$ about a pivot.
\item In pulley problems with a massive pulley, the two tensions are different; writing $T_1=T_2$ is the most common error.
\item Conservation of moment of momentum does \emph{not} imply conservation of kinetic energy when bodies stick or $I$ changes.
\item For the compound pendulum, the formula $T=2\pi\sqrt{I_O/(Mgh)}$ requires \emph{small} oscillations; state this assumption.
\item Keep all angles in radians when using $\omega^2=\omega_0^2+2\alpha\theta$ and torque equations.
\end{itemize}
\end{attention}

\newpage

\coursec{Exercises}

\courssub{I check my knowledge}

\begin{exercice}[Moment of inertia of point masses]
Three particles of masses \SI{1}{kg}, \SI{2}{kg} and \SI{3}{kg} are fixed at the vertices of a light equilateral triangular frame of side \SI{0.6}{m}. Find the moment of inertia and the radius of gyration about an axis through the \SI{1}{kg} mass perpendicular to the plane of the triangle.
\end{exercice}

\begin{exercice}[Radius of gyration of a uniform rod]
A uniform rod of mass \SI{4}{kg} and length \SI{2}{m} rotates about an axis through one end perpendicular to the rod. State the moment of inertia and calculate the radius of gyration about this axis.
\end{exercice}

\begin{exercice}[Standard moment of inertia by integration]
Prove by integration that the moment of inertia of a uniform disc of mass $M$ and radius $a$ about a diameter is $\frac{1}{4}Ma^2$. You may quote the moment of inertia of a disc about its central axis perpendicular to its plane.
\end{exercice}

\begin{exercice}[Constant torque and angular motion]
A flywheel of moment of inertia \SI{0.8}{kg.m^2} rotates at \SI{300}{rev/min}. It is brought to rest in \SI{15}{s} by a constant braking couple. Find the magnitude of the couple and the number of revolutions made while stopping.
\end{exercice}

\courssub{I practise}

\begin{exercice}[Masses on a pulley with moment of inertia]
Two masses of \SI{3}{kg} and \SI{5}{kg} hang from a light inextensible string passing over a pulley of mass \SI{4}{kg} and radius \SI{0.2}{m} modelled as a uniform disc. The system is released from rest. Find the acceleration of the masses and the tension in each part of the string.
\end{exercice}

\begin{exercice}[Rod released from horizontal]
A uniform rod AB of length \SI{1.2}{m} and mass \SI{2}{kg} is free to rotate in a vertical plane about a smooth horizontal axis through A. The rod is released from rest when horizontal. Find its angular speed when vertical and the magnitude of the reaction at A at that instant.
\end{exercice}

\begin{exercice}[Disc with added mass on rim]
A uniform disc of mass \SI{2}{kg} and radius \SI{0.5}{m} rotates freely at \SI{4}{rad/s} about its fixed vertical axis. A particle of mass \SI{0.5}{kg} is gently placed on the rim and adheres to it. Find the new angular speed and the loss of kinetic energy.
\end{exercice}

\begin{exercice}[Compound pendulum: rod and disc]
A compound pendulum consists of a uniform rod AB of mass $M$ and length $2a$, with a uniform disc of mass $M$ and radius $a$ fixed to the rod at B, the plane of the disc being perpendicular to the rod. The pendulum oscillates in a vertical plane about a horizontal axis through A. Find the moment of inertia about A, the period of small oscillations, and the reaction on the axis when the pendulum passes through its lowest position with angular speed $\omega$.
\end{exercice}

\courssub{I prepare for the GCE}

\begin{exercice}[Small oscillations of a rod]
A uniform rod of mass $M$ and length $2a$ makes small oscillations in a vertical plane about one end. Show that the period is $2\pi\sqrt{\frac{4a}{3g}}$ and state the length of the equivalent simple pendulum. [7 marks]
\end{exercice}

\begin{exercice}[Impulsive tension and energy]
A light string is wrapped round a pulley of moment of inertia $I$ and radius $r$. A mass $m$ is attached to the free end of the string. The mass is released from rest and falls a distance $h$ before the string suddenly becomes taut. Using energy and impulsive tension, find the angular speed of the pulley just after the jerk. [8 marks]
\end{exercice}

\begin{exercice}[GCE structured question: compound pendulum]
A compound pendulum is made by rigidly attaching a uniform disc of mass $2m$ and radius $a$ to the end B of a uniform rod AB of mass $m$ and length $2a$. The disc lies in the same plane as the rod and its centre is at B. The pendulum is free to rotate in a vertical plane about a smooth horizontal axis through A.
\begin{enumerate}
\item Find the moment of inertia of the pendulum about the axis through A. [4 marks]
\item The pendulum is released from rest with AB horizontal. Find the angular speed when AB is vertical. [5 marks]
\item Find the magnitude and direction of the reaction at the axis A when the pendulum passes through the vertical position. [6 marks]
\item Determine the period of small oscillations about the equilibrium position. [5 marks]
\end{enumerate}
\end{exercice}

\newpage

\coursec{Solutions to the exercises}

\begin{corrige}[Exercise 1 — Moment of inertia of point masses]
The moment of inertia of a system of point masses about an axis is \(I = \sum m_i r_i^2\), where \(r_i\) is the perpendicular distance from the axis.

The axis passes through the \(1\,\mathrm{kg}\) mass and is perpendicular to the plane of the triangle. Hence the \(1\,\mathrm{kg}\) mass lies on the axis (\(r=0\)). The other two masses are at the other vertices of the equilateral triangle of side \(0.6\,\mathrm{m}\); each is at distance \(0.6\,\mathrm{m}\) from the axis.

\[
I = (2\,\mathrm{kg})(0.6\,\mathrm{m})^2 + (3\,\mathrm{kg})(0.6\,\mathrm{m})^2 = (2+3)\times 0.36 = 1.8\,\mathrm{kg\,m^2}.
\]

Total mass \(M = 1+2+3 = 6\,\mathrm{kg}\). The radius of gyration \(k\) is defined by \(I = M k^2\):

\[
k = \sqrt{\frac{I}{M}} = \sqrt{\frac{1.8}{6}} = \sqrt{0.3} \approx 0.548\,\mathrm{m}.
\]

\emph{Answer:} \(I = 1.8\,\mathrm{kg\,m^2}\), \(k = \sqrt{0.3}\,\mathrm{m} \approx 0.548\,\mathrm{m}\).
\end{corrige}

\begin{corrige}[Exercise 2 — Radius of gyration of a uniform rod]
For a uniform rod of mass \(M\) and length \(L\) about an axis through one end perpendicular to the rod, the standard moment of inertia is

\[
I = \frac{1}{3} M L^2.
\]

Here \(M = 4\,\mathrm{kg}\), \(L = 2\,\mathrm{m}\):

\[
I = \frac{1}{3} \times 4 \times (2)^2 = \frac{16}{3}\,\mathrm{kg\,m^2}.
\]

The radius of gyration \(k\) satisfies \(I = M k^2\):

\[
k = \sqrt{\frac{I}{M}} = \sqrt{\frac{16/3}{4}} = \sqrt{\frac{4}{3}} = \frac{2}{\sqrt{3}} = \frac{2\sqrt{3}}{3}\,\mathrm{m} \approx 1.155\,\mathrm{m}.
\]

\emph{Answer:} \(I = \frac{16}{3}\,\mathrm{kg\,m^2}\), \(k = \frac{2\sqrt{3}}{3}\,\mathrm{m}\).
\end{corrige}

\begin{corrige}[Exercise 3 — Standard moment of inertia by integration]
We prove that the moment of inertia of a uniform disc of mass \(M\) and radius \(a\) about a diameter is \(\frac{1}{4}Ma^2\).

\emph{Method 1 (Perpendicular axis theorem):}
The moment of inertia of the disc about its central axis perpendicular to its plane is \(I_z = \frac{1}{2}Ma^2\). Let \(I_x\) and \(I_y\) be the moments of inertia about two perpendicular diameters in the plane of the disc. By symmetry \(I_x = I_y\). The perpendicular axis theorem states \(I_z = I_x + I_y = 2I_x\). Hence

\[
I_x = \frac{1}{2} I_z = \frac{1}{2} \times \frac{1}{2}Ma^2 = \frac{1}{4}Ma^2.
\]

\emph{Method 2 (Direct integration):}
Place the disc in the \(xy\)-plane with centre at the origin. The diameter is the \(x\)-axis. Consider a thin strip parallel to the \(x\)-axis at distance \(y\) from it, of width \(\mathrm{d}y\). The strip has length \(2\sqrt{a^2-y^2}\). Its area is \(\mathrm{d}A = 2\sqrt{a^2-y^2}\,\mathrm{d}y\). Mass per unit area \(\sigma = M/(\pi a^2)\). The moment of inertia of the strip about the \(x\)-axis is \(\mathrm{d}I = y^2\,\mathrm{d}m = y^2 \sigma\,\mathrm{d}A\). Integrating from \(y=-a\) to \(y=a\):

\[
I = \sigma \int_{-a}^{a} 2y^2\sqrt{a^2-y^2}\,\mathrm{d}y.
\]

Using the substitution \(y = a\sin\theta\), \(\mathrm{d}y = a\cos\theta\,\mathrm{d}\theta\), \(\sqrt{a^2-y^2} = a\cos\theta\), the integral becomes

\[
I = \sigma \int_{-\pi/2}^{\pi/2} 2a^2\sin^2\theta \cdot a\cos\theta \cdot a\cos\theta\,\mathrm{d}\theta = 2\sigma a^4 \int_{-\pi/2}^{\pi/2} \sin^2\theta\cos^2\theta\,\mathrm{d}\theta.
\]

Since \(\sin^2\theta\cos^2\theta = \frac{1}{4}\sin^2 2\theta = \frac{1}{8}(1-\cos 4\theta)\), the integral evaluates to \(\frac{\pi}{8}\). Thus

\[
I = 2\sigma a^4 \cdot \frac{\pi}{8} = \frac{\pi}{4}\sigma a^4 = \frac{\pi}{4} \cdot \frac{M}{\pi a^2} \cdot a^4 = \frac{1}{4}Ma^2.
\]

\emph{Answer:} \(I = \frac{1}{4}Ma^2\) (proved).
\end{corrige}

\begin{corrige}[Exercise 4 — Constant torque and angular motion]
Initial angular speed: \(\omega_0 = 300\,\mathrm{rev/min} = 300 \times \frac{2\pi}{60} = 10\pi\,\mathrm{rad/s}\).
Final angular speed \(\omega = 0\) after time \(t = 15\,\mathrm{s}\).
Constant angular acceleration \(\alpha\):

\[
\omega = \omega_0 + \alpha t \quad\Rightarrow\quad 0 = 10\pi + 15\alpha \quad\Rightarrow\quad \alpha = -\frac{2\pi}{3}\,\mathrm{rad/s^2}.
\]

Magnitude of braking couple (torque): \(\tau = I|\alpha| = 0.8 \times \frac{2\pi}{3} = \frac{1.6\pi}{3} \approx 1.676\,\mathrm{N\,m}\).

Angular displacement during stopping:

\[
\theta = \omega_0 t + \frac{1}{2}\alpha t^2 = 10\pi \times 15 + \frac{1}{2}\left(-\frac{2\pi}{3}\right) \times 225 = 150\pi - 75\pi = 75\pi\,\mathrm{rad}.
\]

Number of revolutions: \(N = \frac{\theta}{2\pi} = \frac{75\pi}{2\pi} = 37.5\).

\emph{Answer:} Couple \(= \frac{1.6\pi}{3}\,\mathrm{N\,m} \approx 1.68\,\mathrm{N\,m}\); revolutions \(= 37.5\).
\end{corrige}

\begin{corrige}[Exercise 5 — Masses on a pulley with moment of inertia]
Let \(T_1\) be the tension on the \(3\,\mathrm{kg}\) side, \(T_2\) on the \(5\,\mathrm{kg}\) side. The pulley is a uniform disc: \(I = \frac{1}{2} M R^2 = \frac{1}{2} \times 4 \times (0.2)^2 = 0.08\,\mathrm{kg\,m^2}\). Let \(a\) be the downward acceleration of the \(5\,\mathrm{kg}\) mass (and upward acceleration of the \(3\,\mathrm{kg}\) mass). The angular acceleration of the pulley is \(\alpha = a/R\).

Equations of motion:
\begin{align*}
5g - T_2 &= 5a \quad &(1)\\
T_1 - 3g &= 3a \quad &(2)\\
(T_2 - T_1)R &= I\alpha = I\frac{a}{R} \quad\Rightarrow\quad T_2 - T_1 = \frac{I}{R^2}a = \frac{0.08}{0.04}a = 2a \quad &(3)
\end{align*}

Add (1) and (2): \(5g - T_2 + T_1 - 3g = 8a \;\Rightarrow\; 2g - (T_2 - T_1) = 8a\).
Substitute (3): \(2g - 2a = 8a \;\Rightarrow\; 2g = 10a \;\Rightarrow\; a = \frac{g}{5} = 1.96\,\mathrm{m/s^2}\) (using \(g=9.8\,\mathrm{m/s^2}\)).

Then
\[
T_2 = 5g - 5a = 5 \times 9.8 - 5 \times 1.96 = 49 - 9.8 = 39.2\,\mathrm{N},
\]
\[
T_1 = 3g + 3a = 3 \times 9.8 + 3 \times 1.96 = 29.4 + 5.88 = 35.28\,\mathrm{N}.
\]

\emph{Answer:} \(a = 1.96\,\mathrm{m/s^2}\), \(T_1 = 35.28\,\mathrm{N}\), \(T_2 = 39.2\,\mathrm{N}\).
\end{corrige}

\begin{corrige}[Exercise 6 — Rod released from horizontal]
Rod: mass \(M = 2\,\mathrm{kg}\), length \(L = 1.2\,\mathrm{m}\). Moment of inertia about end \(A\):
\[
I = \frac{1}{3} M L^2 = \frac{1}{3} \times 2 \times (1.2)^2 = 0.96\,\mathrm{kg\,m^2}.
\]

Loss of gravitational potential energy when the rod falls from horizontal to vertical: the centre of mass drops by \(L/2 = 0.6\,\mathrm{m}\).
\[
\Delta U = M g \frac{L}{2} = 2 \times 9.8 \times 0.6 = 11.76\,\mathrm{J}.
\]

This equals the gain in rotational kinetic energy:
\[
\frac{1}{2} I \omega^2 = 11.76 \quad\Rightarrow\quad \omega^2 = \frac{2 \times 11.76}{0.96} = 24.5 \quad\Rightarrow\quad \omega = \sqrt{24.5} \approx 4.95\,\mathrm{rad/s}.
\]

\emph{Reaction at \(A\) when vertical:}
At the lowest point, the centre of mass moves in a circle of radius \(L/2\). The centripetal force required is
\[
F_c = M \frac{v_{\mathrm{cm}}^2}{L/2} = M \frac{(\omega L/2)^2}{L/2} = M \omega^2 \frac{L}{2}.
\]
Forces on the rod: weight \(Mg\) downwards, reaction \(R\) upwards. Net upward force provides centripetal acceleration:
\[
R - Mg = M \omega^2 \frac{L}{2}.
\]
Using \(\omega^2 = 3g/L\) (from energy: \(\frac{1}{2} \cdot \frac{1}{3}ML^2 \omega^2 = Mg\frac{L}{2} \Rightarrow \omega^2 = 3g/L\)):
\[
M \omega^2 \frac{L}{2} = M \cdot \frac{3g}{L} \cdot \frac{L}{2} = \frac{3}{2}Mg.
\]
Hence
\[
R = Mg + \frac{3}{2}Mg = \frac{5}{2}Mg = 2.5 \times 2 \times 9.8 = 49\,\mathrm{N}.
\]
Direction: vertically upward.

\emph{Answer:} \(\omega \approx 4.95\,\mathrm{rad/s}\), \(R = 49\,\mathrm{N}\) upward.
\end{corrige}

\begin{corrige}[Exercise 7 — Disc with added mass on rim]
Disc: \(M = 2\,\mathrm{kg}\), \(R = 0.5\,\mathrm{m}\). Initial moment of inertia:
\[
I_{\mathrm{disc}} = \frac{1}{2} M R^2 = \frac{1}{2} \times 2 \times (0.5)^2 = 0.25\,\mathrm{kg\,m^2}.
\]
Initial angular speed \(\omega_0 = 4\,\mathrm{rad/s}\). Initial angular momentum:
\[
L = I_{\mathrm{disc}} \omega_0 = 0.25 \times 4 = 1\,\mathrm{kg\,m^2/s}.
\]

Particle of mass \(m = 0.5\,\mathrm{kg}\) placed on the rim. Its moment of inertia about the axis: \(m R^2 = 0.5 \times 0.25 = 0.125\,\mathrm{kg\,m^2}\).
Total moment of inertia after adhesion:
\[
I_{\mathrm{total}} = 0.25 + 0.125 = 0.375\,\mathrm{kg\,m^2}.
\]

Angular momentum is conserved (no external torque about the axis):
\[
\omega_{\mathrm{new}} = \frac{L}{I_{\mathrm{total}}} = \frac{1}{0.375} = \frac{8}{3} \approx 2.667\,\mathrm{rad/s}.
\]

\emph{Answer:} New angular speed \(\omega = \frac{8}{3}\,\mathrm{rad/s} \approx 2.67\,\mathrm{rad/s}\).
\end{corrige}

\newpage

\coursec{Summary sheet}

\begin{popfigure}
\begin{tikzpicture}[
  mindmap,
  grow cyclic,
  every node/.style={concept, text=white, font=\small},
  root concept/.append style={concept color=popink, font=\bfseries\small, minimum size=2.6cm},
  level 1/.append style={level distance=3.4cm, sibling angle=60},
  level 1 concept/.append style={font=\footnotesize, minimum size=2.1cm},
]
\node[root concept]{Rotational Dynamics}
  child[concept color=popblue]{ node {Moment of inertia $I=\sum m r^{2}$} }
  child[concept color=poporange]{ node {Radius of gyration $I=Mk^{2}$} }
  child[concept color=popgreen]{ node {Standard results by integration} }
  child[concept color=poppurple]{ node {Constant torque $G=I\ddot{\theta}$} }
  child[concept color=poppink]{ node {Moment of momentum $L=I\omega$; impulse; conservation} }
  child[concept color=popgold]{ node {Energy $\tfrac12 I\omega^{2}$; force on axis} }
  child[concept color=popdark]{ node {Compound pendulum $T=2\pi\sqrt{\frac{k^{2}+h^{2}}{gh}}$} };
\end{tikzpicture}
\legende{Mind map of the chapter Rotational Dynamics.}
\end{popfigure}

\begin{bilan}

\textbf{Key definitions}
\begin{itemize}
\item \cle{Moment of inertia} of a particle: $I=mr^{2}$ ($r$ = perpendicular distance from the axis). For a system or rigid body: $I=\sum m_i r_i^{2}=\int r^{2}\,dm$. Unit: $\mathrm{kg\,m^{2}}$.
\item \cle{Radius of gyration} $k$: $I=Mk^{2}$, i.e. $k=\sqrt{I/M}$ — the distance at which the whole mass could be concentrated to give the same $I$.
\item \cle{Torque} (moment of a force) about an axis: $G=Fd$ where $d$ is the perpendicular distance from the axis to the line of action.
\item \cle{Moment of momentum} (angular momentum) about a fixed axis: $L=I\omega$.
\item \cle{Compound pendulum}: a rigid body swinging freely about a fixed horizontal axis.
\end{itemize}

\textbf{Standard moments of inertia (proved by integration)}
\begin{itemize}
\item Uniform rod, mass $M$, length $2a$, about centre: $I=\tfrac13 Ma^{2}$; about one end: $I=\tfrac43 Ma^{2}$.
\item Uniform hoop (ring) radius $a$ about central axis: $I=Ma^{2}$.
\item Uniform disc radius $a$ about central axis: $I=\tfrac12 Ma^{2}$.
\item Uniform sphere radius $a$ about a diameter: $I=\tfrac25 Ma^{2}$.
\item \cle{Parallel axes theorem}: $I=I_G+Mh^{2}$ ($I_G$ about the parallel axis through the centre of mass, $h$ = distance between axes).
\item \cle{Perpendicular axes theorem} (lamina): $I_z=I_x+I_y$.
\end{itemize}

\textbf{Laws and formulas to know}
\begin{itemize}
\item Equation of rotational motion about a fixed axis: $\boxed{G=I\ddot{\theta}}$ (rotational analogue of $F=ma$).
\item Constant torque $G$: constant angular acceleration; the constant-acceleration formulas apply with $s\to\theta$, $u\to\omega_0$, $v\to\omega$, $a\to\ddot\theta$:
\[ \omega=\omega_0+\ddot\theta\,t,\qquad \theta=\omega_0 t+\tfrac12\ddot\theta\,t^{2},\qquad \omega^{2}=\omega_0^{2}+2\ddot\theta\,\theta. \]
\item Angular impulse: change of moment of momentum, $J=I\omega_2-I\omega_1$.
\item \cle{Conservation of moment of momentum}: if the total external moment about the axis is zero, $I_1\omega_1=I_2\omega_2$ (e.g. two bodies engaging on the same axis).
\item Kinetic energy of rotation: $E_k=\tfrac12 I\omega^{2}$; work done by a constant torque: $W=G\theta$; power: $P=G\omega$.
\item Work–energy principle: $G\theta=\tfrac12 I\omega^{2}-\tfrac12 I\omega_0^{2}$; conservation of energy: $\tfrac12 I\omega^{2}+Mgh=\text{constant}$.
\item Compound pendulum: period of small oscillations
\[ T=2\pi\sqrt{\frac{I}{Mgh}}=2\pi\sqrt{\frac{k^{2}+h^{2}}{gh}}, \]
where $h$ = distance from axis to centre of mass; equivalent simple pendulum length $\ell=\dfrac{k^{2}+h^{2}}{h}$. Minimum period when $h=k$.
\item Force on the axis: resolve along and perpendicular to $GO$ using $\sum F=M\ddot x,\ M\ddot y$ with $\ddot x=-h\omega^{2}$, $\ddot y=h\ddot\theta$ (components of acceleration of $G$), together with $G=I_O\ddot\theta$.
\end{itemize}

\textbf{Methods}
\begin{enumerate}
\item \emph{Finding $I$ by integration}: choose a strip/element of mass $dm=\rho\,dV$ at distance $r$ from the axis; write $I=\int r^{2}\,dm$; express everything in one variable; integrate between correct limits.
\item \emph{Constant torque problems}: compute $I$ (use parallel axes if needed), then $\ddot\theta=G/I$, then use the angular kinematics formulas.
\item \emph{Connected particles}: write $F=ma$ for each translating particle and $G=I\ddot\theta$ for the rotating body; link with $a=r\ddot\theta$; solve simultaneously.
\item \emph{Impulse problems}: apply angular impulse = change of $I\omega$ to each body; use conservation of moment of momentum at engagement.
\item \emph{Compound pendulum}: find $I$ about the pivot (parallel axes), then $T=2\pi\sqrt{I/(Mgh)}$.
\end{enumerate}

\textbf{Common mistakes}
\begin{itemize}
\item Confusing $I$ about the centre with $I$ about an end — always check the axis and use parallel axes correctly ($I=I_G+Mh^{2}$, never $I_G-Mh^{2}$).
\item Mixing units: angles must be in \textbf{radians} in all formulas.
\item Using $G=I\ddot\theta$ with $I$ about the wrong axis (the axis must be the fixed axis of rotation).
\item Forgetting that energy is \emph{not} conserved in inelastic engagements (impulse problems), even though moment of momentum is.
\item Taking moments of forces that pass through the axis (their moment is zero).
\item Sign errors in the force on the axis: draw a clear diagram and fix positive directions first.
\end{itemize}

\textbf{GCE tips}
\begin{itemize}
\item Learn the four standard results and be ready to \emph{prove} any of them by integration — this is a frequent examinable proof.
\item Show the equation $G=I\ddot\theta$ explicitly before substituting numbers; method marks are awarded for it.
\item In compound pendulum questions, quote $T=2\pi\sqrt{I/(Mgh)}$ and define $h$ clearly on a diagram.
\item Keep answers in exact form (surds, $\pi$) unless a decimal is requested; give units ($\mathrm{kg\,m^{2}}$, $\mathrm{rad\,s^{-1}}$, $\mathrm{N\,m}$).
\item Check dimensional consistency: $I\omega^{2}$ has units of energy — a quick way to validate a result.
\end{itemize}
\end{bilan}

\findecours

\end{document}
